% CUE_GCT isolation track. Run CUE-GCT-20261121-I.
% Model: GPT-6 Astra Pro; configured effort is not exposed to this agent.
% Bibliographic keys required: Hoza2026, AGKS2014, ForbesShpilka2013, MulmuleyGCTV, AllenderOgihara1996.

\section{Isolation: the exact transfer, the quotient, and the surviving gate}
\label{sec:isolation-track}

We use a field $F$, $n$ variables, individual degree at most $\delta$, and
an ROABP expression
\[
 f(\mathbf x)=u^{\mathsf T}A_1(x_{\pi(1)})\cdots A_n(x_{\pi(n)})v.
\]
The order $\pi$ may be unknown.  A polynomial-size hitting set here means
one set depending only on $(n,w,\delta,F)$, not on the matrices or their
coefficients, which hits every nonzero member of the class.  Arithmetic
complexity and bit complexity are distinguished throughout.

\paragraph{Background I-B01.}
The stipulated new graph theorem is Hoza (2026), Theorem 1.3, pp.~2--3;
its quantitative version is Theorem 2.4, p.~5.  For every simple layered
digraph of width $W$ and length $N$, some seed in a family of size
$(WN)^{O(1)}$ gives a minimum-weight path unique for every reachable pair.
The weights are nonnegative integers bounded by $B=(WN)^{O(1)}$ and are
computable in logarithmic space.  A path weight sums its internal vertex
weights.  With a binary code $C$ of length $\ell$ and distance $d$, the
formula is
\[
 \rho_x(a,i)=6^{\nu_2(i)}\sum_{j=1}^{\ell}6^{\ell-j}C(a)_j x_j
 \quad(i>0),\qquad \rho_x(a,0)=0,
\tag{I.1}\label{eq:I-hoza}
\]
and the probability of a bad uniform seed is at most
$2^{-d}W^4N^2$.  The graph is fixed before the seed is chosen.
We use this result as supplied; the present work does not claim an
independent verification of its complete theorem.

\paragraph{Background I-B02.}
Agrawal, Gurjar, Korwar, and Saxena (2014), \S3, Definition~5, call
$\omega\in\mathbb Z_{\geq0}^{n}$ basis isolating for
$D=\sum_a v_a\mathbf x^a$ if there is a coefficient basis indexed by $B$
whose monomial weights are distinct, and each $v_a$ with $a\notin B$
belongs to the span of basis coefficients of strictly smaller weight.
Their Lemma~6 preserves every nonzero scalar projection under
$x_i=t^{\omega_i}$; Lemma~8 and Theorem~1 give their quasipolynomial
construction, including unknown variable order.  Those known results
are cited, not claimed as findings here.

\subsection{I-01. A genuine polynomial transfer for vertex-variable path polynomials}
\label{I-01}

\paragraph{Statement.}
Let $G$ range over simple layered digraphs with fixed vertex slots,
width $W$, length $N\ge1$, fixed source $s$ and target $t$, and arbitrary
scalar edge labels $c_e\in F$.  Give every internal vertex slot its own
variable $z_v$, and form
\[
 P_G(\mathbf z)=\sum_{p:s\leadsto t}\left(\prod_{e\in p}c_e\right)
                        \prod_{v\in\operatorname{int}(p)}z_v.
\]
There is a uniform hitting set of cardinality at most
\[
 2^{\ell}\bigl((N-1)B+1\bigr)=(WN)^{O(1)}
\tag{I.2}
\]
over any field containing $(N-1)B+1$ distinct elements.  Its points have
coordinates $z_v=a^{\rho_x(v)}$, where $x$ runs over Hoza's seeds and
$a$ runs over a fixed set of $(N-1)B+1$ field elements.  The same conclusion
holds for independent edge variables after subdividing edges; the new
width is at most $W^2$ and the new length is $2N$.

\paragraph{Proof.}
Delete zero-labelled edges.  In a simple layered digraph two different
$s$--$t$ paths have different sequences of internal vertices and hence
different vertex monomials, except in the length-one case, where there
is at most one edge.  Thus $P_G\ne0$ exactly when some path survives.
Choose a min-isolating seed for the surviving graph using I-B01.  In
$P_G(t^{\rho_x(v)})$, the coefficient of the least path weight is the
product of labels on its unique minimum path, and is nonzero in the
field.  The degree is at most $(N-1)B$.  A nonzero polynomial of that
degree has at most $(N-1)B$ roots.  This proves (I.2), uniformly in the
edge labels.  For edge variables, insert one distinct intermediate
vertex on each edge and give it that edge's variable; give the original
internal vertices auxiliary independent variables as well.  The same
argument provides a hitting set for the expanded vertex-variable
polynomial.  To obtain points only on edge variables, assign to an edge
the product of its intermediate-vertex value and the value of its head
when the head is an internal original vertex.  Each original path then
receives exactly the product of all its expanded internal-vertex values.
This combines exponents additively, leaves a polynomial degree bound,
and proves the edge-variable statement. \hfill$\square$

\paragraph{Scope.}
This is an exact consequence of the new graph theorem.  It does not
identify several vertex or edge variables with the same ROABP variable.
That identification is the substantive next step, treated below.
Over a smaller finite field, the argument gives an extension-field
hitting set after constructing a sufficiently large extension; it is
not a statement that the same number of base-field points exists.

\subsection{I-02. The quotient obstruction is already present in a two-path diamond}
\label{I-02}

\paragraph{Statement.}
Let the monomial map $q:F[z_1,\ldots,z_s]\longrightarrow
F[x_1,\ldots,x_n]$ send $z_j$ to $\mathbf x^{a_j}$, and let $A$ be
the matrix with columns $a_j$.  A diagonal one-parameter action
$z_j\mapsto t^{r_j}z_j$ is induced by a diagonal action
$x_i\mapsto t^{\omega_i}x_i$ precisely when
$r=A^{\mathsf T}\omega$.  If the $a_j$ generate a proper sublattice,
requiring $r$ to annihilate $\ker(A:\mathbb Z^s\to\mathbb Z^n)$
only guarantees a character on the image lattice; an integral
extension to all $\mathbb Z^n$ is an additional lattice condition.
Consequently path isolation does not, in general, descend through $q$.

\paragraph{Proof and counterexample.}
Equivariance on each generator reads
$t^{r_j}\mathbf x^{a_j}=t^{\langle\omega,a_j\rangle}\mathbf x^{a_j}$,
which is equivalent to the displayed matrix equation.  The kernel
description follows from the factorization property of homomorphisms
of free abelian groups, with the stated extension qualification.

For the counterexample take the simple diamond
$s\to u\to t$ and $s\to v\to t$.  Label the two paths so that their
vertex-variable polynomial is $z_u-z_v$.  The map $q(z_u)=q(z_v)=x$
gives the zero polynomial.  Every min-isolating weight on this diamond
has $r_u\ne r_v$, so that its weighted polynomial
$t^{r_u}-t^{r_v}$ is nonzero.  Such a weight does not descend, because
$z_u-z_v\in\ker q$ is not homogeneous for it.  This remains a valid
counterexample in characteristic two, where the two distinct powers
still have nonzero coefficients.  Alternatively give both paths label
$1$ in characteristic zero: the original polynomial is $2x\ne0$,
but every weight induced from $x\mapsto t^{\omega}$ ties its two paths.
Thus compatible minimum-path isolation is not even necessary for
nonzero scalar output, while arbitrary path isolation is not sufficient
to certify nonzero output after the quotient. \hfill$\square$

\paragraph{What is refuted.}
The refuted assertion is the direct identification of a Hoza seed with
an ROABP variable-weight assignment.  Neither polynomial ROABP hitting
sets nor their possible derivation by a different construction is
refuted by this example.

\subsection{I-03. Three exact characterizations of basis isolation}
\label{I-03}

\paragraph{Statement.}
Let $S\subset\mathbb Z_{\geq0}^n$ index distinct monomials and let
$V=\operatorname{span}_F\{v_a:a\in S\}$ have dimension $r$ over an arbitrary field.
For $\omega\in\mathbb Z^n$ set $w_a=\langle\omega,a\rangle$.
For integer weights we use the same order-theoretic definition as
I-B02, allowing negative variable weights; the hitting-set
applications below use nonnegative weights.
The following are equivalent.
\begin{enumerate}
\item $\omega$ is a basis-isolating weight assignment in the sense of
this extended definition.
\item Every nonzero linear functional $\lambda\in V^*$ makes
$\sum_a\lambda(v_a)\mathbf x^a$ have a unique monomial of minimum weight.
\item At each weight level $c$, at most one coefficient $v_a$ with
$w_a=c$ survives modulo
$V_{<c}=\operatorname{span}\{v_b:w_b<c\}$.
\end{enumerate}
For a fixed coefficient basis $B$, write uniquely
$v_a=\sum_{b\in B}\lambda_{ab}v_b$.  It is the basis isolated by
$\omega$ exactly when
\[
 \langle\omega,a-b\rangle>0
 \quad(a\notin B,\ \lambda_{ab}\ne0),
 \qquad
 \langle\omega,b-b'\rangle\ne0\quad(b\ne b'\in B).
\tag{I.3}\label{eq:I-cone}
\]
The first inequalities alone characterize $B$ as the unique
minimum-total-weight coefficient basis.

\paragraph{Proof.}
Assume (1), with isolated basis $B$.  For a nonzero $\lambda$, take
the least-weight basis member $b$ on which $\lambda$ is nonzero.
It is unique.  Every nonbasis coefficient of no greater weight lies
in the span of smaller basis coefficients and is annihilated by
$\lambda$.  This proves (2), including its stronger unique-monomial
assertion.

If (3) fails, choose two nonzero coefficient classes at a level $c$.
There is a functional on $V/V_{<c}$ nonzero on both: for independent
classes assign value $1$ to both and extend a basis; for dependent
nonzero classes use any functional nonzero on one of them.  Pull it
back to $V$.  It vanishes below $c$ and is nonzero on two monomials
at $c$, contradicting (2).  This argument works also over finite
fields.  Thus (2) implies (3).

Under (3), choose the unique surviving monomial at every level where
the coefficient span grows.  These coefficients form a basis, one
per growth level.  All remaining coefficients at a level already
belong to the span of smaller levels.  Iterating through the finitely
many levels proves (1).  Unique coordinates in a fixed basis prove
the equivalence with (I.3).

If $B$ is the unique minimum-total-weight basis and
$\lambda_{ab}\ne0$, then $B-b+a$ is also a basis.  Its total weight
is strictly larger, so $w_a>w_b$.  Conversely, suppose these strict
inequalities hold.  Expand any other basis $C$ in the basis $B$.
Its change-of-basis determinant is nonzero, so at least one
permutation in its determinant expansion is nonzero.  Every common
member of $B\cap C$ is forced to match itself; every member
$a\in C\setminus B$ matches some $b\in B\setminus C$ with
$\lambda_{ab}\ne0$.  Summing the strict inequalities over this
matching gives $\sum_{a\in C}w_a>\sum_{b\in B}w_b$.
Thus $B$ is the unique minimum-total-weight basis. \hfill$\square$

\paragraph{Geometric consequence.}
For a fixed ordering of the weights within $B$, (I.3) is an open
rational polyhedral cone in the cocharacter space $\mathbb R^n$.
The union of these cones is precisely the set of BIWAs.  Equivalently,
the one-parameter subgroup simultaneously gives a monomial initial
form for every nonzero scalar projection of $D$.  Isolating one path
or one preselected scalar output has weaker quantifiers.

\subsection{I-04. Projected matroid vertices and polynomial random weights}
\label{I-04}

\paragraph{Statement.}
Let $\mathcal B$ be the coefficient bases of the data in I-03, and put
\[
 \beta(B)=\sum_{a\in B}a,\qquad
 Q=\operatorname{conv}\{\beta(B):B\in\mathcal B\}\subset\mathbb R^n.
\]
Every vertex of $Q$ is the image of exactly one indexed basis.
Assume $S\subseteq\{0,\ldots,\delta\}^{n}$, $r\ge1$, $\delta\ge1$.
Choose independently
$a_i\in\{1,\ldots,M_1\}$ and $b_i\in\{1,\ldots,M_2\}$.
With $L=n\delta M_2+1$, the weights
\[
 \omega_i=L a_i+b_i
\tag{I.4}
\]
are basis isolating with probability at least
\[
 1-\frac{nr\delta}{M_1}-\frac{\binom r2}{M_2}.
\tag{I.5}
\]
For example $M_1=4nr\delta$ and $M_2=4\max(1,r^2)$ give a positive
constant success probability and polynomially bounded weights.  This
is a randomized existence statement, not an explicit small family.

\paragraph{Proof of the vertex assertion.}
Suppose two distinct bases $B,C$ have the same image $\beta$.
Take $e\in B\setminus C$.  Let $\ell$ be the coordinate functional
of $v_e$ in the basis $B$, and expand $v_e=\sum_{f\in C}c_fv_f$.
Because $1=\ell(v_e)=\sum c_f\ell(v_f)$, some
$f\in C\setminus B$ has both $c_f\ne0$ and $\ell(v_f)\ne0$.
Hence $B-e+f$ and $C-f+e$ are both bases.  Their images are
$\beta+(f-e)$ and $\beta-(f-e)$, which are distinct because the
ground-set exponents are distinct.  Their midpoint is $\beta$,
so $\beta$ is not a vertex.  Every vertex has at least one preimage
from the finite defining set, proving the claim.

\paragraph{Proof of the probability bound.}
First consider any finite set $T\subseteq\{0,\ldots,D\}^{n}$ and
uniform independent weights $a_i\in[M_1]$.  Fix all weights except
$a_i$.  The least weight as a function of a real parameter $t=a_i$
is the lower envelope of at most $D+1$ affine functions, one for
each possible $i$-th coordinate.  Its slopes are among
$0,\ldots,D$, and there are at most $D$ breakpoints at which
different slopes minimize.  If two different points of $T$ minimize,
they differ in some coordinate $i$, and that coordinate weight is
one of these breakpoints.  Conditioning on the other coordinates
and applying a union bound gives probability at most $nD/M_1$.
Apply this with $T=\{\beta(B)\}$ and $D=r\delta$.
With probability at least $1-nr\delta/M_1$, there is a unique
minimizing point of $Q$.  It is a vertex and, by the preceding
argument, has a unique basis preimage $B$.

Conditional on the first weights and this $B$, two different
monomials in $B$ have equal $b$-weight with probability at most
$1/M_2$: condition on every $b_i$ but one whose exponent difference
is nonzero.  There is at most one possible value of that remaining
integer weight.  A union bound over pairs costs
$\binom r2/M_2$.  By I-03, the unique first-weight minimum basis
satisfies $a\cdot m>a\cdot b$ whenever the $b$-coordinate of
$v_m$ is nonzero and $m\notin B$.  Each difference is an integer
at least one.  The perturbation changes it by at most
$n\delta M_2$, so multiplication by $L=n\delta M_2+1$ preserves
strict positivity.  It also preserves inequalities between basis
members with unequal first weights, and separates any remaining
ties by the second weights.  Equation (I.3) proves (I.5).
All randomness and weights in this proof are integers; no assumption
on the characteristic of the coefficient field enters the argument.
\hfill$\square$

\paragraph{The remaining derandomization problem.}
The desired generator must hit these cones for coefficient matroids
arising from products of univariate matrices, with a family and
maximum weights both polynomial in $(n,w,\delta)$.  Hoza's theorem
hits path-polytope cones in independent vertex coordinates.  I-02
shows why those cocharacter spaces cannot be silently identified.
I-04 narrows the required object to normal cones of a specified
projected representable-matroid polytope, with the finite
within-basis tie hyperplanes removed.

\subsection{I-05. Positive Gram determinants certify bases, not compressed rank}
\label{I-05}

\paragraph{Statement.}
For complex coefficient vectors, choose coordinates $V=\mathbb C^r$
and form
\[
 G_D(\mathbf z)=\det\left(\sum_{a\in S}v_av_a^*\mathbf z^a\right).
\]
Then
\[
 G_D(\mathbf z)=\sum_{B\in\mathcal B}
                  |\det C_B|^2\mathbf z^{\beta(B)},
\qquad \operatorname{Newt}(G_D)=Q,
\tag{I.6}
\]
where $C$ has columns $v_a$.  At a vertex of $Q$, the coefficient
is the squared volume of its unique basis from I-04.  Nevertheless,
positivity of $G_D(t^\omega)$ does not imply that the coefficient
span of $D(t^\omega)$ retains its dimension.

\paragraph{Proof.}
Apply the Cauchy--Binet identity to
$C\operatorname{diag}(\mathbf z^a)C^*$.
Every nonzero summand is strictly positive in its coefficient, so
there are no cancellations between bases with the same exponent
sum; their convex hull is exactly the Newton polytope.  I-04
identifies its vertex coefficients.  For the negative assertion,
take $D(x,y)=e_1x+e_2y$ and $\omega=(1,1)$.  Its Gram determinant
is $xy$ and specializes to $t^2$, whereas
$D(t,t)=t(e_1+e_2)$ has coefficient span of dimension one instead
of two.  The unique coefficient basis has two equal member weights;
this is exactly the extra condition in I-03. \hfill$\square$

\paragraph{Field and geometric qualification.}
The conjugate transpose and positivity in (I.6) are statements on
a complex vector space with a Hermitian structure.  They are not
positivity statements in an arbitrary characteristic.  The
determinantal/Cauchy--Binet identity itself has an integral
two-matrix form with independent dual variables.  Specializing the
second matrix to a conjugate transpose is a real, not a complex
algebraic, operation.  This qualification is essential in translating
the formula into a scheme over $\mathbb Z$.

\subsection{I-06. Exact black-box implementation of nonlinear exponent weights can be exponential}
\label{I-06}

\paragraph{Statement.}
Let $K$ be any field, let $t$ be transcendental over it, and define
a linear functional on polynomials of individual degree at most $2d$
by
\[
 \mathcal L(\mathbf x^\alpha)=\prod_{i=1}^n h_i(\alpha_i),
 \qquad h_i(e)\in K(t).
\]
Put $H_i=(h_i(a+b))_{0\le a,b\le d}$.
If one exact universal implementation is
\[
 \mathcal L(f)=\sum_{j=1}^M c_j f(\xi_j),
\tag{I.7}
\]
where the points and coefficients are fixed independently of $f$
over any extension of $K(t)$, then
\[
 M\ \ge\ \prod_{i=1}^n\operatorname{rank}(H_i).
\tag{I.8}
\]
In particular $\mathcal L(\mathbf x^\alpha)=
t^{\sum_i\alpha_i^2}$ requires $M\ge(d+1)^n$.
For degree two, the nonlinear exponent filters supplied by every
Hoza seed that separates three exponent symbols also have the
lower bound $M\ge2^n$.

\paragraph{Proof.}
Index a moment matrix by
$\alpha,\beta\in\{0,\ldots,d\}^n$ and set
$H_{\alpha\beta}=\mathcal L(\mathbf x^{\alpha+\beta})$.
It is the tensor product $\bigotimes_iH_i$, with rank the product
of the ranks.  Each evaluation in (I.7) contributes the rank-one
matrix $c_j(\xi_j^\alpha)_\alpha(\xi_j^\beta)_\beta^{\mathsf T}$,
so its rank is at most $M$.  Extension of scalars does not change
the rank, proving (I.8).

For $h_i(e)=t^{e^2}$, factor the $(a,b)$ entry as
$t^{a^2}t^{b^2}(t^{2a})^b$.  The remaining matrix is Vandermonde
on the distinct elements $1,t^2,\ldots,t^{2d}$ and has full rank,
in every characteristic.

For the Hoza assertion take three distinct binary codewords
$C(0),C(1),C(2)$ and put
$h_x(e)=\sum_j6^{\ell-j}C(e)_j x_j$ as integers.
If $h_x(0)+h_x(2)=2h_x(1)$, then
\[
 \sum_j6^{\ell-j}x_j\bigl(C(0)_j+C(2)_j-2C(1)_j\bigr)=0.
\]
The digits lie between $-2$ and $2$.  The first nonzero digit
dominates the absolute sum of all later digits, since
$2\sum_{k=0}^{q-1}6^k<6^q$.  Thus every digit vanishes, which for
binary digits says the three selected code bits agree in every
selected coordinate.  Hence all three values $h_x(e)$ coincide.
Conversely constant values give that equality.  Any seed isolating
the three-arm graph is nonconstant on its three internal vertices,
so the equality fails.  At layer $i$, the scalar multiplier
$6^{\nu_2(i)}$ in (I.1) preserves this failure.  Consequently the
$2\times2$ Hankel determinant of the exponent filter
\[
 e\longmapsto t^{6^{\nu_2(i)}h_x(e)}
\]
is the difference of two distinct monomials, and has nonzero
determinant over $K(t)$.  Equation (I.8) now gives $2^n$.
\hfill$\square$

\paragraph{Scope of the lower bound.}
This rules out the proposed exact linear interpolation implementation
of independently reweighted monomial letters by a polynomial number
of evaluations of the original black box.  It is not a lower bound
on general ROABP PIT, adaptive reconstruction, or approximate
integration.  Even requiring (I.7) only for width-one ROABPs implies
the same identities on every monomial and hence the same lower
bound by linearity; width-one PIT itself remains elementary.

\subsection{I-07. Haar averaging repairs invariance by discarding state isolation}
\label{I-07}

\paragraph{Statement.}
For a fixed complex $m\times m$ matrix $D$, the function
$A\mapsto\operatorname{tr}(DA)$ is invariant under simultaneous
conjugation if and only if $D$ is scalar.  Moreover
\[
 \int_{U(m)}UDU^*\,dU=\frac{\operatorname{tr}D}{m}I.
\tag{I.9}
\]
Thus independently averaging fixed diagonal state-weight insertions
in a trace word replaces each insertion by a scalar and recovers
a scalar multiple of the original trace word.  If the same unitary
is used for two insertions, the exact formula for $m\ge2$ is
\[
 \int\operatorname{tr}(A UDU^* B UEU^*)\,dU
 =\alpha\operatorname{tr}(AB)+\beta\operatorname{tr}A\operatorname{tr}B,
\tag{I.10}
\]
where
\[
 \alpha=\frac{m\operatorname{tr}D\operatorname{tr}E-
                    \operatorname{tr}(DE)}{m(m^2-1)},\qquad
 \beta=\frac{m\operatorname{tr}(DE)-
                    \operatorname{tr}D\operatorname{tr}E}{m(m^2-1)}.
\]
Even positive semidefinite $D,E$ need not make these coefficients
nonnegative: for orthogonal rank-one projections,
$\beta=-1/[m(m^2-1)]$.

\paragraph{Proof.}
Conjugation invariance for every $A$ gives
$\operatorname{tr}((g^{-1}Dg-D)A)=0$ for all $A$.  The trace pairing
is nondegenerate, so $D$ commutes with every invertible matrix and
is scalar.  The reverse direction is immediate.  Haar invariance
makes the left side of (I.9) commute with every unitary; the
commutant is scalar, and its trace determines the scalar.  Repeated
integration gives the independent-insertion claim.

For completeness, in (I.10) define
$T=\int(UDU^*)\otimes(UEU^*)\,dU$.  The commutant of
$U\otimes U$ is spanned by the identity $I$ and tensor flip $S$.
Here is an elementary verification: commutation with diagonal unitaries
forces matrix entries to preserve the unordered pair of tensor indices;
permutation unitaries make the two-dimensional off-diagonal blocks
$\left(\begin{smallmatrix}b&c\\c&b\end{smallmatrix}\right)$ and
the diagonal blocks a common scalar $a$.  A rotation mixing two basis
vectors forces $a=b+c$.  The resulting operator is $bI+cS$.
Thus $T=\alpha I+\beta S$.  Its trace and its trace after multiplication
by $S$ are respectively
\[
 m^2\alpha+m\beta=\operatorname{tr}D\operatorname{tr}E,
 \qquad m\alpha+m^2\beta=\operatorname{tr}(DE).
\]
Solving gives the stated values.  Finally
$\operatorname{tr}(A X B Y)=
\operatorname{tr}((A\otimes B)(X\otimes Y)S)$;
contracting $T$ gives (I.10).  For orthogonal rank-one projections,
$\operatorname{tr}D=\operatorname{tr}E=1$ and
$\operatorname{tr}(DE)=0$, proving the sign assertion.
\hfill$\square$

\paragraph{Meaning for the proposed transfer.}
Hoza's internal-state weights enter matrix products as fixed
diagonal insertions.  Their lack of conjugation invariance cannot
be repaired by independent Haar averaging while preserving the
state information.  Correlated averaging changes the trace word
into a combination of different trace invariants and can reintroduce
negative coefficients.  Neither averaging is the missing
basis-isolation construction.

\subsection{I-08. The exact matrix-invariant consequence, conditional on the remaining gate}
\label{I-08}

\paragraph{Statement.}
Suppose one has an explicit hitting set $\mathcal H\subset F^{m^2}$
of polynomial size, constructible in polynomial arithmetic time,
for ROABPs of depth $m^2$, width at most $2m^2$, and individual
degree less than $r$.  Let $F$ have characteristic zero and put
\[
 M(x)=\sum_{j=0}^{r-1}M_jx^j,\qquad
 T_{\ell,\alpha}(\mathbf M)=
 \operatorname{tr}\bigl(M(\alpha_1)\cdots M(\alpha_\ell)\bigr),
 \quad1\le\ell\le m^2,quad\alpha\in\mathcal H.
\tag{I.11}
\]
These at most $m^2|\mathcal H|$ invariants are an explicit
polynomial-size separating family for simultaneous conjugation.
Over algebraically closed $F$, the invariant ring is finite over
the subring they generate.  This is precisely the specialization
of Forbes and Shpilka (2013), Construction~3.7, Theorem~3.8 and
Corollary~3.10, to a polynomial hitting set.  It supplies the
matrix-invariant instance of explicit Noether normalization in
the sense of that paper.  It does not by itself extract an
algebraically independent system of parameters of cardinality
equal to the quotient dimension.

\paragraph{Derivation and exact source dependence.}
The coefficient of $x_1^{a_1}\cdots x_\ell^{a_\ell}$ in
$\operatorname{tr}(M(x_1)\cdots M(x_\ell))$ is the trace word
$\operatorname{tr}(M_{a_1}\cdots M_{a_\ell})$.  The latter words,
for $\ell\le m^2$, generate the characteristic-zero invariant ring
as recorded in Forbes--Shpilka Theorem~1.3 and Corollary~3.6.
For two tuples $A,B$, the difference of their trace-generating
polynomials is an ROABP of width at most $2m^2$, by their Lemma~3.2.
The hitting-set assumption therefore detects exactly whether
some trace generator separates $A$ and $B$.  Their Theorem~3.9
(the stated specialization of Derksen--Kemper Theorem~2.3.12)
then gives integrality, and finite generation upgrades this to
a finite morphism.  The only new hypothesis in this specialization
is the polynomial hitting-set gate; it has not been proved here.
\hfill$\square$

\paragraph{Consequences that do not follow.}
This does not prove general arithmetic PIT, arbitrary explicit
Noether normalization, permanent circuit lower bounds, or
$P\ne NP$.  Forbes--Shpilka already prove deterministic NC
orbit-closure intersection in Theorem~1.15 using white-box PIT.
Their discussion after Theorems~1.13--1.15 explicitly separates
the matrix-invariant case from the stronger general explicit-variety
problem.  In the supplied revised GCT~V source, the general-PIT
equivalence is Theorem~5.14, and the matrix-specific ROABP replacement
is \S7.5, Lemma~7.9.  The earlier GCT~V theorem numbering cited by
Forbes--Shpilka must not be transferred to the revised source.  No new polynomial hitting set
may be inserted into (I.11) until its own gate is closed.

\subsection{I-09. An exact positive transfer-matrix certificate remains white box}
\label{I-09}

\paragraph{Statement.}
For a complex ROABP in known matrix-product form
\[
 f(\mathbf x)=u^{\mathsf T}\prod_{i=1}^n
                  \left(\sum_{a=0}^{\delta}A_{i,a}x_i^a\right)v
             =\sum_{\alpha}c_{\alpha}\mathbf x^{\alpha},
\]
its coefficient squared-norm generating polynomial is
\[
 \sum_{\alpha}|c_\alpha|^2\mathbf z^\alpha
 = (u\otimes\bar u)^{\mathsf T}
     \prod_{i=1}^n\left(\sum_{a=0}^{\delta}
                    A_{i,a}\otimes\overline{A_{i,a}}\,z_i^a\right)
                     (v\otimes\bar v).
\tag{I.12}
\]
In particular $f=0$ if and only if the right side at
$\mathbf z=\mathbf1$ is zero.  Its matrices have width $w^2$.
The equality also gives the usual torus identity
\[
 \int_{[0,2\pi]^n}|f(e^{i\theta_1},\ldots,e^{i\theta_n})|^2
                \prod_i\frac{d\theta_i}{2\pi}
       =\sum_\alpha|c_\alpha|^2.
\tag{I.13}
\]

\paragraph{Proof.}
The coefficient $c_\alpha$ is
$u^{\mathsf T}A_{1,\alpha_1}\cdots A_{n,\alpha_n}v$.
Multiplicativity of tensor products identifies its squared
absolute value with the product of the corresponding tensor
matrices and boundary vectors.  Summing independently over each
$\alpha_i$ gives (I.12).  The sum is a finite sum of nonnegative
real numbers, and it is zero exactly when every coefficient is
zero.  In (I.13), expansion of the square and integration of
$e^{i\langle\alpha-\beta,\theta\rangle}$ kills exactly the terms
$\alpha\ne\beta$. \hfill$\square$

\paragraph{Why this does not close the black-box gate.}
The transfer matrices use the individual coefficient matrices
$A_{i,a}$, which a black box does not supply.  The torus integral
is another exact expression, but the product quadrature that is
exact for all such degree bounds uses $(\delta+1)^n$ nodes.
No lower bound on all specialized ROABP quadratures is asserted.
I-06 proves a lower bound for the different, explicitly specified
nonlinear exponent filter.  Positivity in (I.12) closes the
white-box zero-testing question only; it does not supply the
missing uniform polynomial set of original evaluations.

\subsection{I-10. Schemes, cocharacters, flags, and flat limits}
\label{I-10}

\paragraph{Statement and dictionary.}
Let $R$ be a domain carrying the relevant coefficients, and let
$T=\mathbb G_{m,R}^{n}$.  The integer weight vector $\omega$
is the cocharacter
$\lambda_\omega:\mathbb G_{m,R}\to T$,
$t\mapsto(t^{\omega_1},\ldots,t^{\omega_n})$.
For nonzero $f=\sum_a c_a\mathbf x^a$ and
$\mu=\min_{c_a\ne0}\langle\omega,a\rangle$, define
\[
 \widetilde f(t,\mathbf x)=
 \sum_a c_a t^{\langle\omega,a\rangle-\mu}\mathbf x^a.
\tag{I.14}
\]
This is a polynomial over $R[t]$, its special fiber is the
weight-initial form, and its generic fiber is the original
polynomial under a torus change of coordinates.  A unique
minimizing exponent gives a monomial special fiber.  Over a field,
the hypersurface family defined by (I.14) is flat over the
$t$-line.  For a homogeneous ideal, the corresponding weight
homogenization is the usual flat Rees/initial-ideal degeneration;
this is stronger than merely selecting a least monomial of one
polynomial and must not be conflated with it.

For a vector polynomial, I-03 identifies a complete coefficient
flag whose successive quotients have one supported weight
monomial.  With an explicit coefficient matrix, nonvanishing
maximal minors define a matroid stratum of the appropriate
Grassmannian.  The exponents of its Pl\"ucker coordinates under
$\lambda_\omega$ are the sums $\beta(B)$.  Its torus orbit has
weight polytope $Q$; I-04 says a vertex with nonzero Pl\"ucker
coordinate corresponds to a unique indexed coefficient basis.
The BIWA cones refine its normal fan by the individual
within-basis weight comparisons in (I.3).

\paragraph{Proof of the flatness assertion and the dictionary.}
Formula (I.14) has no negative power of $t$ and is not divisible
by $t$.  Over a field $F$, for every irreducible polynomial
$p(t)$ the image of $\widetilde f$ modulo $p(t)$ is nonzero:
for $p=t$ use its nonzero initial form; otherwise $t$ is invertible
in $F[t]/(p)$ and its monomial coefficients remain nonzero.
Thus no nonzero $p(t)$ divides $\widetilde f$ in the UFD
$F[t,\mathbf x]$.  If $p(t)g\in(\widetilde f)$, primality of
the irreducible factors of $\widetilde f$ implies
$g\in(\widetilde f)$.  Hence the quotient by
$(\widetilde f)$ is torsion-free over the PID $F[t]$, and
therefore flat.  The generic-fiber statement follows by
inverting $t$ in (I.14).  Exterior products of coefficient
columns are precisely Pl\"ucker coordinates, and their weights
add, giving $\beta(B)$.  Equation (I.3) then identifies the
normal-cone refinement.  This proves each part of the stated
dictionary directly. \hfill$\square$

\paragraph{Integral qualification.}
The torus and monomial maps are defined over $\mathbb Z$.
A specified coefficient matroid and its rank flag persist on an
open localization where the finitely many designated nonzero
minors remain nonzero; they need not persist in every residue
characteristic.  Over a general base $R$, flatness of a
hypersurface quotient also needs the relevant coefficient and
fiber hypotheses.  The field proof above is not a claim of
automatic flatness over arbitrary $R$.  The Grothendieck locator
table in the main paper should be used for the exact SGA/FGA/EGA
placement of these notions, rather than inventing a theorem
number for these elementary derived statements.

\subsection{I-11. Exterior and subspace-state repairs do not yet produce a graph generator}
\label{I-11}

\paragraph{Checked finite assertions.}
For coefficient ambient dimension $k$, the standard coordinate
space $\bigwedge^rF^k$ has dimension $\binom kr$.  Over
$\mathbb F_q$, there are at least $q^{r(k-r)}$ distinct
$r$-dimensional subspaces of $F_q^k$: the graphs of the
$q^{r(k-r)}$ linear maps $F_q^r\to F_q^{k-r}$ are all distinct.
Thus a literal graph with one state for each such subspace is
not polynomial in $k$ when $r$ grows with $k$.

Exterior coordinates also retain cancellation.  For two equal
vectors $v_1=v_2=e_1+e_2$, the determinant is
$1\cdot1-1\cdot1=0$, although its expansion has two nonzero
permutation terms.  Making one term uniquely minimum by
path-dependent weights changes the determinant rather than
certifying its nonvanishing.  Tensoring $r$ independently chosen
word trajectories gives a standard realization with $w^{O(r)}$
states; for fixed $r$ this removes that particular size cost,
but the cancellation and weight-descent problems remain.

\paragraph{Proof and scope.}
The exterior dimension follows from the basis wedges indexed
by $r$-element subsets of a fixed $k$-element basis.  The graph
of a linear map is an $r$-dimensional subspace, and projection
onto the first summand recovers the map, proving distinctness.
The determinant example is immediate.  The tensor product
construction records one of $w$ state positions for each of
the $r$ trajectories, so has $w^r$ positions before additional
boundary coordinates; determinant signs are imposed at the
boundary and do not disappear.  These are size and correctness
facts about the specified constructions, not lower bounds on
all possible reductions.  A scalar output can itself have
coefficient-space dimension one, so constant coefficient rank
alone cannot make the general ROABP hitting-set task disappear.
\hfill$\square$

\subsection{I-12. Resolving coefficient support is an arithmetic counting gate}
\label{I-12}

\paragraph{Background I-B05.}
Allender and Ogihara (1996), \S2, p.~4, define the accepting-minus-rejecting
gap model; \S5, Definition~2, p.~11, defines $C_{=}L$ by the vanishing
of a $\operatorname{GapL}$ function.  Their Theorem~1, p.~2, identifies
$\operatorname{GapL}$ with logarithmic-space reductions to integer
determinant, and Theorem~14, p.~12, gives the corresponding singularity
completeness result.  The following ROABP formulation is proved directly.

\paragraph{Statement.}
Restrict to explicitly given layered ROABPs over $\mathbb Z$
(equivalently over $\mathbb Q$) with coefficients in
$\{0,1,-1\}$ and polynomially bounded width and length.  The input
lists the coefficient entries and the queried exponent vector;
univariate edge labels may be given densely or as canonical sparse
lists with binary exponents.  Deciding
whether a specified monomial coefficient is nonzero is complete,
under deterministic logarithmic-space reductions, for the class
\[
 \operatorname{coC}_{=}L
   =\{\{x:g(x)\ne0\}:g\in\operatorname{GapL}\}.
\]
Here $\operatorname{GapL}$ is defined by the difference of the
numbers of accepting and rejecting paths of polynomial-time
nondeterministic logarithmic-space machines.  Hardness already
holds for ROABPs whose polynomial is one monomial times a scalar.
Therefore a logarithmic-space uniform polynomial-size reduction
which represents exactly nonzero coefficient support by ordinary
directed reachability would imply
$\operatorname{coC}_{=}L\subseteq NL$.

\paragraph{Proof.}
For membership, guess a state path in the layered ROABP.  At layer
$i$ choose the term of the prescribed exponent from its univariate
edge label.  A zero term contributes nothing; otherwise track the
parity of the negative signs.  The difference between positive and
negative paths is exactly the specified monomial coefficient.
State indices, a layer counter, and a sign bit require logarithmic
space.  Neutral or invalid branches can be arranged to contribute
one accepting and one rejecting path, hence zero gap.

For hardness, start with a polynomial-time nondeterministic
logarithmic-space machine computing a gap $g(x)$.  Its time-unrolled
configuration graph has polynomial width and length and is
constructible from $x$ in logarithmic space.  Pad halted computations
by deterministic transitions, retaining their accepting or rejecting
tag, so every path has the same length.  Mark the final edge
of accepting paths by $1$ and of rejecting paths by $-1$;
all preceding edges have label $1$.  Separate parallel transitions
by an intermediate layer if necessary.  Multiply every edge label
in layer $i$ by a fresh variable $x_i$.  Every source-to-target path
then has the same variable monomial $\prod_i x_i$, with its
accepting or rejecting sign, so the output is
$g(x)\prod_i x_i$.  This is a known-order multilinear ROABP with
coefficients $0,1,-1$.  Its unique possible monomial coefficient
is nonzero exactly when $g(x)\ne0$.  Finally, composition with the
postulated graph reduction puts every such language in $NL$.
\hfill$\square$

\paragraph{Relation to the other gates.}
This is a precise complexity formulation of the operation which
would have to cancel equal-monomial paths before graph isolation.
It is not a separation of complexity classes.  Nor does it rule
out a different polynomial-time algebraic construction: the
explicit ROABP supplies its matrices, and white-box polynomial
identity testing is already possible.  In scheme language the
coefficient is a regular function of the edge coefficients;
testing its vanishing uses the kernel of a coefficient map, not
merely whether its parameterizing path set is empty.  The diamond
in I-02 is the smallest instance of this distinction.

\subsection{I-13. Exact universal energy cubature needs the full coefficient dimension}
\label{I-13}

\paragraph{Statement.}
Let $n\ge1$, $\delta\ge0$, and $D=(\delta+1)^n$.  Fix points
$\xi_1,\ldots,\xi_Q\in\mathbb C^n$ and a complex $Q\times Q$ matrix $C$,
independent of the polynomial.  Suppose that every width-one product
\[
 f(\mathbf x)=\prod_{i=1}^n\left(\sum_{a=0}^{\delta}z_{i,a}x_i^a\right),
 \qquad z_{i,a}\in\mathbb C,
\]
satisfies the exact identity
\[
 \sum_{\alpha\in\{0,\ldots,\delta\}^n}|[\mathbf x^\alpha]f|^2
  =\sum_{j,k=1}^Q C_{jk}\,\overline{f(\xi_j)}f(\xi_k).
\tag{I.13}\label{eq:I-cubature}
\]
Then $Q\ge D$.  In particular, signed cubature
$\sum_j c_j|f(\xi_j)|^2$ requires at least $D$ nodes.  More generally,
a formula with $M$ fixed paired terms
$\sum_{j=1}^M c_j\overline{f(\xi_j)}f(\eta_j)$ requires $M\ge D$.
The bound is attained for the diagonal formula by the Cartesian product
of the $(\delta+1)$st roots of unity in each coordinate, all with weight
$1/D$.

\paragraph{Proof.}
Index coordinates by $\alpha\in\{0,\ldots,\delta\}^n$ and put
$a_\alpha=\prod_i z_{i,\alpha_i}$.  Let the evaluation matrix be
$E_{j\alpha}=\xi_j^\alpha$.  The asserted identity is
\[
 a^*(I_D-E^*CE)a=0
\quad\text{for every decomposable coefficient vector }
a=\bigotimes_i(z_{i,0},\ldots,z_{i,\delta}).
\]
Expanding gives a polynomial identity in the variables
$z_{i,a}$ and their conjugates.  Its monomials
\[
 \prod_i\overline{z_{i,\alpha_i}}z_{i,\beta_i},
 \qquad (\alpha,\beta)\in\{0,\ldots,\delta\}^{2n},
\]
are linearly independent as functions of the complex variables.
Indeed, writing each $z=u+iv$ identifies the polynomial ring in
formal $z,\bar z$ with the polynomial ring in $u,v$ by an invertible
complex linear change of variables; a polynomial vanishing on every
real tuple $(u,v)$ has every coefficient zero.  Distinct pairs
$(\alpha,\beta)$ give distinct displayed monomials.  Consequently
$I_D=E^*CE$, so
$D=\operatorname{rank}(I_D)\le\operatorname{rank}(E)\le Q$.
For paired terms the same coefficient comparison represents $I_D$
as a sum of $M$ rank-at-most-one matrices, proving $M\ge D$.

Finally, averaging $\xi^\alpha\overline{\xi^\beta}$ over the stated
root-of-unity grid gives $1$ when $\alpha=\beta$ and $0$ otherwise:
in each coordinate the difference lies between $-\delta$ and
$\delta$, whose only multiple of $\delta+1$ is zero.  Expanding
$|f|^2$ therefore proves the matching formula, for every polynomial
of the stated individual degree, including the width-one class.
\hfill$\square$

\paragraph{Scope and consequence for I-09.}
Torus orthogonality identifies the left side of (I.13) with
$\int_{(S^1)^n}|f(z)|^2\,dz$.  Thus the positive white-box energy
computation in I-09 cannot be replaced by a polynomial number of
fixed black-box samples through any exact quadratic formula of the
displayed kind, even for width one and fixed $\delta\ge1$ as $n$ grows.
This is a lower bound for universal exact cubature, not for hitting
sets: for width one the substitution $x_i=t$ gives the nonzero
univariate product $\prod_i f_i(t)$ of degree at most $n\delta$.
Evaluation at $n\delta+1$ distinct scalars already gives a polynomial
hitting set for that class.  The rank argument makes no claim about
adaptive procedures or other nonlinear uses of evaluations.

\subsection{I-14. Joint compression is a Grassmannian transversality condition}
\label{I-14}

\paragraph{Statement.}
Write $D=\sum_{a\in S}v_a\mathbf x^a$, choose coordinates on
$V=\operatorname{span}\{v_a\}$, and let $C=(v_a)_{a\in S}$ be the
resulting full-row-rank $r\times s$ matrix, where $s=|S|$.
For a finite family $\mathcal W$ of integer weights, form a matrix
$P_{\mathcal W}$ with rows indexed by $S$ and columns indexed by
the pairs $(\omega,c)$ that occur as a monomial weight.  Substitutions
are interpreted as Laurent polynomials if necessary.  Set
\[
 (P_{\mathcal W})_{a,(\omega,c)}
   =\begin{cases}1&\langle\omega,a\rangle=c,\\0&\text{otherwise.}\end{cases}
\]
The following conditions are equivalent:
\begin{enumerate}
\item the coefficients of all the substitutions
$D(t^{\omega_1},\ldots,t^{\omega_n})$, taken jointly, span $V$;
\item $\operatorname{rank}(CP_{\mathcal W})=r$;
\item $\operatorname{im}(C^{\mathsf T})\cap
\ker(P_{\mathcal W}^{\mathsf T})=\{0\}$;
\item for every nonzero $\lambda\in V^*$, at least one substituted
scalar polynomial $\lambda(D(t^\omega))$ is nonzero.
\end{enumerate}
Equivalently, for some $r$-element set $T$ of columns of $P_{\mathcal W}$,
\[
 \sum_{\substack{B\subseteq S\\|B|=r}}
   \det(C_B)\det((P_{\mathcal W})_{B,T})\ne0.
\tag{I.14}\label{eq:I-joint-Plucker}
\]
Thus the failure locus on $\operatorname{Gr}(r,F^s)$ is the
intersection with the linear equations obtained by setting every
expression in (I.14) to zero.  These depend on Pl\"ucker coordinates,
including their values, rather than only on the coefficient matroid.

\paragraph{Proof.}
The $(\omega,c)$ column of $CP_{\mathcal W}$ is exactly
$\sum_{\langle\omega,a\rangle=c}v_a$, the coefficient at that
weight after substitution.  This proves (1)$\Leftrightarrow$(2).
Since $C^{\mathsf T}:F^r\to F^s$ is injective, the transpose of
$CP_{\mathcal W}$ is injective exactly when the intersection in
(3) is zero.  A proper subspace of $V$ has a nonzero annihilating
linear functional, proving (1)$\Leftrightarrow$(4).  Finally,
Cauchy--Binet applied to each $r$-column submatrix of
$CP_{\mathcal W}$ gives (I.14).  These are linear forms in the
Pl\"ucker coordinates $\det(C_B)$, proving the geometric description.

Joint preservation can hold even though no individual member is
basis isolating.  Take
\[
 D=e_1(x_1-x_2)+e_2(x_2-x_3),\qquad
 \mathcal W=\{(0,0,1),(1,0,0)\}.
\]
The two substituted vector polynomials are respectively
$(1-t)e_2$ and $(t-1)e_1$.  Each has coefficient rank one and so
fails I-03, but their coefficients jointly span the original
two-dimensional space.  To see why the matroid alone is insufficient,
over a field of characteristic different from two compare coefficient
rows $(1,1)$ and $(1,-1)$ on the support $\{x_1,x_2\}$ under weight
$(1,1)$.  They have the same rank-one coefficient matroid, while
their compressed rows are respectively $2$ and $0$.
\hfill$\square$

\paragraph{A sufficient black-box gate.}
If $S\subseteq\{0,\ldots,\delta\}^n$, all weights are nonnegative
and at most $B$, and joint preservation holds, evaluating the
$|\mathcal W|$ curves at $n\delta B+1$ distinct field elements
hits every nonzero scalar projection of $D$.  Indeed, condition
(4) supplies a nonzero univariate polynomial of degree at most
$n\delta B$.  This exact sufficient condition is weaker than
requiring one BIWA.  It does not remove cancellation from the
linear Pl\"ucker expressions in (I.14).

\subsection{I-15. Bounded affine decodings of the seed masks have width-two obstructions}
\label{I-15}

\paragraph{Statement.}
Suppose a proposed family of variable weights has the form
\[
 \omega(s)=b+Ts,\qquad
 b\in\mathbb Z^n,\quad T\in\mathbb Z^{n\times k},
\]
where every entry of $b$ and $T$ has absolute value at most $M$.
The seed set can be any subset of $\mathbb Z^k$.  If
\[
 2^n>(2nM+1)^{k+1},
\tag{I.15}\label{eq:I-affine-mask-count}
\]
there are distinct $\alpha,\beta\in\{0,1\}^n$ such that the
nonzero multilinear binomial $\mathbf x^\alpha-\mathbf x^\beta$
vanishes identically on every curve $\mathbf x=t^{\omega(s)}$.
This binomial has a width-two ROABP in every variable order.
In particular the whole family fails even the joint condition I-14.

\paragraph{Proof.}
Map the $2^n$ Boolean exponent vectors $\alpha$ to
\[
 \left(b^{\mathsf T}\alpha,T^{\mathsf T}\alpha\right)\in
 \{-nM,\ldots,nM\}^{k+1}.
\]
By (I.15), two distinct vectors have the same image.  Then
$\langle\omega(s),\alpha-\beta\rangle=0$ for every seed $s$,
so the two substituted powers of $t$ coincide.  This is a Laurent
identity if negative weights are allowed, and a polynomial identity
when the weights are nonnegative.  Distinct monomials have nonzero
coefficients in their difference over any field, including
characteristic two.  The diagonal matrix product with factors
$\operatorname{diag}(x_i^{\alpha_i},x_i^{\beta_i})$ and boundaries
$(1,-1)$ and $(1,1)^{\mathsf T}$ computes the binomial.  Reordering
the commuting diagonal factors gives every variable order.
\hfill$\square$

\paragraph{A necessary dimension for any bounded weight family.}
More generally, let $\mathcal W\subseteq\mathbb Z^n$ have all
coordinates of absolute value at most $B\ge1$, and let
$k=\dim_{\mathbb Q}\operatorname{span}_{\mathbb Q}\mathcal W$.
If $\mathcal W$ hits every multilinear binomial by monomial curves,
then necessarily
\[
 k\ge\frac{n}{\log_2(2nB+1)}.
\]
Indeed, select $k$ members of $\mathcal W$ forming a rational basis
and apply the same pigeonhole argument to their integer dot products
with Boolean exponent vectors.  Equality on this basis implies
equality on every member of $\mathcal W$, independently of how
the seed was decoded.  Thus polynomially bounded output weights
computed as polynomials of any fixed degree $h$ in
$\ell=O(\log n)$ Boolean seed bits also fail for sufficiently large
$n$: their span has dimension at most
$\sum_{j=0}^h\binom{\ell}{j}=O((\log n)^h)$, below the necessary
bound.  The dimension estimate follows by grouping the coefficient
vectors of the multilinear seed monomials; on Boolean seeds powers
of a bit reduce to that bit.

\paragraph{Application to the actual mask formula.}
Every weight in Hoza's formula (I.1) is an integer linear form in
the common $\ell$ seed bits.  Hence assigning those forms to
variables, or taking any fixed linear averaging or decoding with
polynomially bounded integer coefficients, gives $\omega(s)=b+Ts$
with $k=\ell$ and polynomially bounded entries.  A rational average
with a common polynomially bounded denominator has the same property
after clearing that denominator.  Since $\ell=O(\log(WN))$, if
$W,N$ and the resulting entry bound are polynomial in the number
$n$ of variables, condition (I.15) holds for all sufficiently large
$n$.  This rules out this specified seed-to-variable decoding,
even if all its members are used jointly.

For the most direct alphabet decoding the counterexample is explicit
without counting.  Let
$H_s(a)=\sum_j6^{\ell-j}C(a)_js_j$.  Matching letters $0,1$ by
an affine exponent map at layer $i$ gives
\[
 \omega_i(s)=6^{\nu_2(i)}\bigl(H_s(1)-H_s(0)\bigr).
\]
Layers $1$ and $3$ receive the same weight, so $x_1-x_3$ vanishes
for every seed.  This uses the actual layer dependence in (I.1)
and holds for every choice of its code.  There is no contradiction
with minimum-path isolation: the graph weights themselves need not
descend to these variable weights.  The argument does not rule out
nonlinear decoders whose outputs span sufficiently many independent
directions, or unbounded integer weights.

\subsection{I-16. The finite-field row-state lift and its uniform evaluation limit}
\label{I-16}

\paragraph{Statement.}
Over $F=\mathbb F_q$, an explicitly given width-$w$ ROABP
\[
 f=u^{\mathsf T}\prod_{i=1}^n
       \left(\sum_{a=0}^{\delta}A_{i,a}x_i^a\right)v
\]
has an exact deterministic coefficient-support automaton with
at most $q^w$ states per layer.  A state is a row vector
$r\in F^{1\times w}$; its letter-$a$ transition at layer $i$ is
$r\mapsto rA_{i,a}$.  The initial state is $u^{\mathsf T}$, and
a final state is accepting exactly when $rv\ne0$.
Nevertheless, if $0\le\delta<q$, every base-field hitting set
$H\subseteq F^n$ for all width-one individual-degree-$\delta$
polynomials satisfies
\[
 |H|\ \ge\ \left(\frac{q}{q-\delta}\right)^n.
\tag{I.16}\label{eq:I-small-field-lower}
\]
For $\delta=q-1$ all $q^n$ points are necessary and sufficient.
Thus polynomial state-space size alone cannot produce a polynomial
uniform base-field hitting set, even when $w=1$ and $q$ is fixed.

\paragraph{Proof.}
On the word $(a_1,\ldots,a_n)$, the automaton's final scalar is
$u^{\mathsf T}A_{1,a_1}\cdots A_{n,a_n}v$, precisely the monomial
coefficient.  Matrix multiplication has already performed all path
cancellation before the state is tested.  To obtain a simple layered
graph, subdivide each transition with an intermediate vertex labelled
by its current row and its letter; append a common sink after all
accepting states.  The width is at most $(\delta+1)q^w$ and the
length is $2n+1$.  It is constructible from the coefficient matrices
in time polynomial in this displayed size and the field-operation
cost.  Hoza applies to this coefficient-dependent graph, with the
corresponding polynomial-in-$q^w$ seed family; its path weights
still depend on the row states.

For the lower bound, choose independently a $\delta$-element subset
$R_i\subseteq F$ for each coordinate and consider
\[
 f_{R_1,\ldots,R_n}(\mathbf x)
     =\prod_{i=1}^n\prod_{a\in R_i}(x_i-a).
\]
These are $\binom{q}{\delta}^{n}$ distinct nonzero width-one polynomials.
A fixed point $\xi\in F^n$ hits exactly
$\binom{q-1}{\delta}^n$ of them, since each $R_i$ must avoid
$\xi_i$.  Covering all these polynomials therefore requires
\[
 |H|\binom{q-1}{\delta}^n\ge\binom{q}{\delta}^{n},
\]
which is (I.16).  The full grid $F^n$ hits every polynomial of
individual degree less than $q$: induction on the number of variables,
using the univariate root bound, shows that a polynomial vanishing
on that grid must be zero.  This also gives the matching upper bound
when $\delta=q-1$.  For $\delta\ge q$, even that grid cannot be a
hitting set for all formally nonzero polynomials, because
$x_1^q-x_1$ vanishes on it. \hfill$\square$

\paragraph{Precise scope.}
The lower bound concerns evaluations in the coefficient field.
It does not apply to extension-field evaluations, which can behave
very differently for a fixed base-field coefficient class.
The row-state lift is useful as an exact finite arithmetic simulation
and avoids the erroneous unweighted support graph of I-02.
It does not construct the graph from an evaluation oracle or make
its state-dependent path weights into variable cocharacters.

\subsection{I-17. A concrete uniform hitting set for invertible factors}
\label{I-17}

\paragraph{Background and statement.}
Invertible-factor ROABP hitting sets are already a result of
Agrawal et al. (2014), Theorem~3, pp.~5 and~19, and their
Lemma~27, p.~22, gives matrix coefficient concentration from
invertible constant terms.  The following is a self-contained
row-space specialization with constant boundary vectors; no
historical-priority claim is made.

Let $n,w\ge1$ and $\delta\ge1$.  Suppose
\[
 f=u^{\mathsf T}A_1(x_{\pi(1)})\cdots A_n(x_{\pi(n)})v
\]
has square $w\times w$ factors of individual degree at most
$\delta$, each with nonzero determinant polynomial.
Choose sets $S,B\subseteq F$ with
$|S|=nw\delta+1$, $|B|=\delta+1$, and $0\in B$.
Then the following single set, independent of the factors and order,
hits every nonzero such $f$:
\[
 \mathcal H=\{a\mathbf1+b: a\in S,\ b\in B^n,
                         |\operatorname{supp}(b)|\le w-1\}.
\tag{I.17}\label{eq:I-invertible-hit}
\]
Its cardinality is at most
\[
 (nw\delta+1)\sum_{j=0}^{\min(w-1,n)}\binom nj\delta^j.
\]
It is polynomial in $n,\delta$ for every fixed $w$, works in arbitrary
characteristic with the stated field-size condition, and requires
no knowledge of $\pi$.

\paragraph{Proof.}
The nonzero polynomial $\prod_i\det A_i(t)$ has degree at most
$nw\delta$, so some $a\in S$ makes every $A_i(a)$ invertible.
Reindex the variables in their true order for this proof only.
Set $P_0=I$, $P_i=A_1(a)\cdots A_i(a)$, and
\[
 M_i(y)=P_{i-1}A_i(a+y)P_i^{-1}.
\]
Then $M_i(0)=I$ and telescoping gives
\[
 f(a\mathbf1+\mathbf y)
    =u^{\mathsf T}M_1(y_1)\cdots M_n(y_n)P_nv.
\]
For nonzero $f$ we have $u\ne0$.  Let $R_i\subseteq F^{1\times w}$
be the span of the coefficient rows of
$u^{\mathsf T}M_1(y_1)\cdots M_i(y_i)$, and put $d_i=\dim R_i$.
Because the next constant term is $I$, these spaces are nested and
$d_0=1$, $d_i\le w$.

We claim that $R_i$ has a basis consisting of actual coefficient
rows whose monomials have support at most $d_i-1$.  This is true for
$i=0$.  Inductively retain a basis for $R_{i-1}$.  If the dimension
does not grow, that basis suffices.  Otherwise extend it using rows
$rM_{i,a}$, where $r$ is a retained coefficient row and $a>0$.
These rows span all new coefficients, and each is itself a coefficient
row with support at most $d_{i-1}$.  Since $d_i>d_{i-1}$, every
retained or added basis row has support at most $d_i-1$.
This proves the claim.

At $i=n$, at least one coefficient row has nonzero product with
$P_nv$.  Therefore some basis coefficient of support at most $w-1$
also has nonzero product.  Restricting the shifted scalar polynomial
to those variables, with the others zero, gives a nonzero polynomial
of individual degree at most $\delta$.  The grid $B$ on those
variables hits it by repeated univariate root counting.  This point
lies in (I.17).  The set (I.17) is invariant under coordinate
permutations, so its construction never needed the true order.
Counting supports proves its size bound. \hfill$\square$

The support bound is sharp for the normalized row-space argument:
for $w\ge2$ and $n=w-1$, factors $I+x_iE_{i,i+1}$ and boundaries $e_1,e_w$
compute $x_1\cdots x_{w-1}$.  For $w=1$, the theorem is the familiar
diagonal univariate test.  Generically singular factors require a
different argument, as the following explicit family shows.

\subsection{I-18. Regularization must preserve a unit Pl\"ucker coordinate}
\label{I-18}

\paragraph{Statement and counterexample.}
Let $m\ge2$, $n=2m$, and
\[
 f(\mathbf x)=\prod_{j=1}^m(x_{2j}-x_{2j-1}).
\]
This is a nonzero multilinear width-two ROABP in the displayed
variable order.  It vanishes at every point $a\mathbf1+b$ with
$|\operatorname{supp}(b)|<m$, for every scalar $a$.
Consequently the sampling restriction in I-17 cannot extend to
all width-two ROABPs by regularizing their matrices and specializing
the regularization parameter to zero.

More precisely, let $R=F[[\varepsilon]]$, $K=F((\varepsilon))$,
$C(\varepsilon)\in R^{r\times s}$ have
$\operatorname{rank}C(0)=r$, and let $P\in R^{s\times h}$ be any
integral compression.  Then
\[
 \operatorname{rank}_K(CP)=r
 \quad\Longleftrightarrow\quad
 \text{some $r$-minor of $CP$ has finite valuation},
\]
whereas
\[
 \operatorname{rank}_F(C(0)P(0))=r
 \quad\Longleftrightarrow\quad
 \text{some $r$-minor of $CP$ has valuation zero}.
\tag{I.18}\label{eq:I-unit-Plucker}
\]
The latter is equivalent to $\operatorname{im}_R(CP)=R^r$.
Generic span equality alone is insufficient.

\paragraph{Proof.}
Use the two singular matrix factors
\[
 L(x)=\begin{pmatrix}-x&1\\0&0\end{pmatrix},\qquad
 R(y)=\begin{pmatrix}1&0\\y&0\end{pmatrix}.
\]
Their product has $(1,1)$ entry $y-x$ and all other entries zero.
Taking the ordered product of these pairs with boundaries $e_1,e_1$
computes $f$.  At $a\mathbf1+b$ each pair difference is
$b_{2j}-b_{2j-1}$.  If fewer than $m$ coordinates of $b$ are
nonzero, one of the $m$ disjoint pairs has both entries zero,
so the product vanishes.

Replace every matrix factor by that factor plus $\varepsilon I$.
All regularized factors are generically invertible over $K$,
because their determinant is monic in $\varepsilon$.
Call the resulting scalar polynomial $f_\varepsilon$.
At the zero point a regularized pair is
\[
 (L(0)+\varepsilon I)(R(0)+\varepsilon I)
  =\begin{pmatrix}\varepsilon(1+\varepsilon)&\varepsilon\\
                   0&\varepsilon^2\end{pmatrix},
\]
so $f_\varepsilon(0)=\varepsilon^m(1+\varepsilon)^m\ne0$ in $K$.
Yet every common-shift sample of perturbation support less than $m$
has zero specialization at $\varepsilon=0$.

For the matrix assertion, a minor is nonzero over $K$ exactly when
it has finite valuation, and its reduction is nonzero over $F$
exactly when its valuation is zero.  This proves the two rank
criteria.  A unit $r$-minor gives $r$ columns invertible over $R$,
so their span is $R^r$.  Conversely, if the columns span $R^r$,
their reductions span $F^r$, which supplies a unit minor.
\hfill$\square$

\paragraph{The resulting special-fiber gate.}
In the scalar example the module generated by all coefficients of
$f_\varepsilon$ is $R$, since some coefficient of $f_0=f$ is nonzero.
For any finite collection of the displayed fixed $F$-valued samples
that contains the zero point, their values generate a nonzero ideal $N\subseteq
\varepsilon R$.  Hence $K N=K$ and its saturation
$(KN)\cap R$ is $R$, although all original sampled values are zero.
Saturating by dividing by a power of $\varepsilon$ thus does not
supply a nonzero evaluation of the original black box.
For vector coefficients the correct next target is a compression
with a unit Pl\"ucker coordinate, as in (I.18), or an equivalent
proof of coefficient-span preservation on the special fiber.
This refines the generic Grassmannian condition in I-14; it is not
established merely by flatness of a surrounding Rees family.

\paragraph{Background I-B06: the successful width-two repair.}
The width-two examples in I-15 and I-18 are failures of those
specified samplers.  Polynomial hitting sets for arbitrary
unknown-order univariate width-two ROABPs are already provided by
Agrawal et al. (2014), \S5.3, Theorem~34, p.~26, with the usual
field-size or extension convention required by the evaluation sets.
Their Lemma~32, p.~25, factors a nonzero singular $2\times2$ matrix
as a rank-one outer product after multiplying by one nonzero entry.
Applied at every singular factor, this writes a nonzero polynomial
multiple of the ROABP output as a product of invertible-factor
ROABP outputs.  The scalar factors may share the variables at the
cuts; product closure does not require disjointness.

Lemma~33, p.~26, supplies that product closure: interpolate a curve
through a common $h$-point hitting set for the factors.  Each nonzero
factor remains nonzero on the curve, so their product remains
nonzero.  If there are $m$ factors, each of total degree at most $D$,
the curve can have coordinate degree $h-1$, giving degree at most
$mD(h-1)$ and requiring that degree plus one parameter values.
The paper states the harmlessly looser $mDh+1$ bound.  For the
univariate-factor specialization of Theorem~34, the factor sparsity
is at most $\delta+1$ and maximum monomial support is one, so the
stated bound is polynomial.  These are cited positive results,
not new claims established by the present counterexamples.

\subsection{I-19. Higher-rank cuts and what block compression actually preserves}
\label{I-19}

\paragraph{Statement.}
Let $A(x)\in F[x]^{w\times w}$ have rank $1\le r<w$ over $F(x)$
and entry degrees at most $\delta$.  Choose $r$ rows $I$ and
$r$ columns $J$ with $M=A_{I,J}$ and $\Delta=\det M\ne0$.
Then
\[
 \Delta A=UV,\qquad
 U=A_{*,J},\quad V=\operatorname{adj}(M)A_{I,*},
\tag{I.19}\label{eq:I-rank-cut}
\]
where $U$ is $w\times r$, $V$ is $r\times w$,
$\deg U\le\delta$, $\deg V\le r\delta$, and
$\deg\Delta\le r\delta$.

Apply (I.19) to the singular factors at positions
$i_1<\cdots<i_m$ of an ROABP.  Write $X$ for their variables,
$Y_j$ for the disjoint intervening variable blocks, and $B_j$
for the product of the original matrices within block $Y_j$.
With $K=F(X)$, multiplication gives the exact identity
\[
 \left(\prod_{j=1}^m\Delta_j\right)f
 =u^{\mathsf T}B_0U_1
    \left(\prod_{j=1}^{m-1}V_jB_jU_{j+1}\right)V_mB_mv.
\tag{I.20}\label{eq:I-outer-cut}
\]
Over $K$, this is a matrix product of width at most
$\max_j r_j$ over disjoint blocks $Y_j$, with the corresponding
rectangular boundary factors.  Over $F$, neighboring displayed
factors generally share a cut variable; that dependence has not
been removed.

If polynomial substitutions into independent parameter sets preserve
the coefficient span of every block factor over $K$, then they
preserve the coefficient span of the full product, including scalar
nonvanishing.  This assertion requires independent parameter sets
for the different blocks.

\paragraph{Proof.}
The columns indexed by $J$ are a basis of the column space over
$F(x)$.  Expressing every column in that basis and restricting to
rows $I$ identifies its coordinate vector as $M^{-1}$ times its
row restriction.  Thus $A=A_{*,J}M^{-1}A_{I,*}$, and clearing
$\Delta$ proves (I.19) with the displayed degree bounds.
Substituting these factorizations into the ordered matrix product
and grouping the intervening factors proves (I.20).  Treating $X$
as elements of the coefficient field leaves exactly the disjoint
blocks $Y_j$; the dimensions in the product are the cut ranks.

For the span assertion, write each block factor as a sum of its
coefficient matrices times monomials in that block.  Since the
blocks are disjoint, the coefficients of the full product are
exactly the ordered products of one coefficient matrix from each
block, with the stated boundary contractions.  Replacing any
coefficient family by another with the same linear span leaves
the span of all these ordered products unchanged, by repeated
multilinearity.  Independent parameter sets preserve this coefficient
description after substitution, proving the claim.
\hfill$\square$

\paragraph{A valid local compression and the remaining coupling.}
For a fixed block, a hitting set for every nonzero scalar linear
combination of its matrix entries preserves its entire coefficient
span: otherwise a nonzero functional would annihilate all sampled
values but not the coefficient span.  A Lagrange curve through
these sample points also preserves that span.  Its coefficient
span contains all the sampled values and is contained in the
original coefficient span, so both inclusions are equalities.
Thus local compression is a concrete valid construction.

Two steps still require proof for a uniform black-box application:
the substitutions must accommodate the hidden block partition, and
the resulting nonzero rational dependence on $X$ must be hit after
returning from $K$ to the original field.  A simple example shows
why the latter cannot be silently dropped.  The rank-one matrices
\[
 A(x)=\begin{pmatrix}1&x\\x&x^2\end{pmatrix},\qquad
 A(y)=\begin{pmatrix}1&y\\y&y^2\end{pmatrix}
\]
with boundaries $e_1,e_1$ give $1+xy$.  After both variables are
put in $K=F(x,y)$, there are no remaining block variables to sample;
all nonvanishing information lies in the coefficient field.

For rank one, the positive product-closure result I-B06 tolerates
the shared cut variables.  For higher rank, (I.20) is a sum over
outer-state paths and retains cancellation.  Hoza isolation of
those paths supplies no proof of the joint-span condition I-14
for the original variables.  The supplied AGKS construction remains
a valid quasipolynomial completion of the general coefficient-span
task; the rank-cut construction above has not replaced its
coefficient and partition dependence by a polynomial procedure.

\paragraph{I-G02: a smaller structural sufficient gate.}
\label{I-G02}
For fixed width $w$, it would suffice to construct a uniform polynomial
hitting set for the expressions (I.20), with outer width at most
$w-1$, unknown block partition, and the shared cut variables kept
as genuine variables.  Their individual degrees are at most
$w\delta$: a cut variable occurs in $U_j$ and $V_j$, with combined
degree at most $(r_j+1)\delta$, and every other variable retains
degree at most $\delta$.  A hitting set for this class hits the
original nonzero $f$, because the multiplier $\prod_j\Delta_j$
is nonzero.  The case with no cuts is I-17, while outer width one
is covered by I-B06.  The first outer rank not handled by this
particular scalar-product repair is rank two.  This was the precise
sufficient gate reached by the rank-cut method.  The fixed-width gate for these original-coordinate cut expressions is
subsequently resolved by the rank-preservation and cone construction
I-31 below; I-30 gives the sharper six-pair result at width three.
No polynomial-in-width generator for the general coupled class is asserted.  For variable
$w$, the exponent depending on $w$ in I-17 is an additional barrier
to the desired polynomial bound.

\subsection{I-20. Bounded total corank admits a uniform extension of I-17}
\label{I-20}

\paragraph{Statement.}
For the constant-boundary ROABP in I-17, allow singular factors and
assume only that their total generic corank is at most an integer
$\kappa\ge0$:
\[
 \sum_{i=1}^n\bigl(w-\operatorname{rank}_{F(x_i)}A_i(x_i)\bigr)
 \le\kappa.
\]
Choose $S,G,B\subseteq F$ of sizes respectively
$nw\delta+1$, $w\delta+1$, and $\delta+1$, with $0\in B$.
For every subset $T\subseteq[n]$ of size at most $\kappa$, assign
the coordinates in $T$ arbitrary values in $G$.  Assign all other
coordinates the values $a+b_i$, where $a\in S$, $b_i\in B$, and
at most $w-1+\kappa$ of these $b_i$ are nonzero.  The union
$\mathcal H_\kappa$ of all such points hits every nonzero polynomial
with the stated representation.  Its size is at most
\[
 (nw\delta+1)
 \left(\sum_{c=0}^{\min(n,\kappa)}\binom nc(w\delta+1)^c\right)
 \left(\sum_{j=0}^{\min(n,w-1+\kappa)}\binom nj\delta^j\right),
\tag{I.21}\label{eq:I-corank-hit}
\]
which is $(nw(\delta+1))^{O(w+\kappa)}$.
The construction is independent of the variable order and of the
locations of singular factors.  It is polynomial for fixed
$w+\kappa$.

\paragraph{Proof.}
For a nonzero output every factor has positive generic rank.
Let $T$ be the true set of singular-factor variables; then
$|T|\le\kappa$.  For each such factor choose a nonzero maximal
minor $\Delta_i(x_i)$ of degree at most
$r_i\delta\le(w-1)\delta$.  Regard $f$ as a polynomial in the
remaining variables and choose any nonzero coefficient $g(X_T)$.
The nonzero polynomial
$g(X_T)\prod_{i\in T}\Delta_i(x_i)$ has individual degree at most
$w\delta$.  The grid $G^T$ therefore contains an assignment
$\gamma$ where it is nonzero.  At that assignment the remaining
polynomial is nonzero and each specialized singular factor has
exactly its generic rank $r_i$.

For the other factors, one common $a\in S$ makes all constant
terms $A_i(a)$ invertible, as in I-17.  Gauge these constant terms
to $I$ by products of invertible matrices, omitting the singular
constant factors from those products.  Concretely, set $P_0=I$,
let $P_i=P_{i-1}A_i(a)$ at a nonsingular layer and
$P_i=P_{i-1}$ at a specialized singular layer, and replace every
factor by $P_{i-1}A_iP_i^{-1}$.  Telescoping adjusts only the final
boundary.  Nonsingular layers now have constant term $I$;
singular layers are constant matrices of rank $r_i$.

Follow the coefficient row spans in this normalized product.
At a nonsingular layer the dimension can only grow.  At a constant
rank-$r_i$ layer it can only decrease, by at most $w-r_i$, because
the kernel of right multiplication has that dimension in the full
row space.  These spans are all nonzero, since the final polynomial
is nonzero.  Starting at dimension one and ending at dimension at
most $w$, the sum $U$ of all dimension increases is consequently
at most
\[
 U\le w-1+\sum_{i\in T}(w-r_i)\le w-1+\kappa.
\]
The greedy argument of I-17 now selects actual coefficient rows
whose support is at most the cumulative number of dimension
increases.  At a constant cut retain a basis of the images of the
old basis; this introduces no new support.  At a layer where the
span grows, every new candidate is an old basis row times one
nonconstant coefficient, so its support grows by at most one,
while the cumulative dimension increase grows by at least one.
This proves the claimed support bound inductively.

Some selected final row has nonzero scalar boundary contraction.
The grid $B$ on its support supplies a nonzero value of the
remaining shifted polynomial.  Together with the chosen cut
values $\gamma$, this is a point of $\mathcal H_\kappa$.
Enumerating all subsets and grid values gives (I.21).  The family
is invariant under coordinate permutations, proving the claimed
unknown-order uniformity. \hfill$\square$

The restriction is explicit: arbitrary ROABPs can have total
generic corank growing with $n$.  I-20 does not bound that parameter
or supply a polynomial procedure when it grows.  It shows that
rank cuts can be absorbed uniformly when the total dimension loss
is bounded; the remaining problem is to compress unboundedly many
such losses without discarding the original coefficient relations.

\subsection{I-B07. Literature scope and the source of the new rank-preservation attempt}
\label{I-B07}

The distinction between a representation in one hidden order and a
width bound in every order is essential.  Gurjar, Korwar, and Saxena
(2017), Theorem~3.6 and Corollary~3.7, give the polynomial
constant-width result with the order known; their Theorem~4.9 treats
the separate every-order class.  Guo and Gurjar (2020), Theorem~1,
give an improved quasipolynomial bound for one unknown order,
whereas their Theorem~3 concerns every-order representations.
Hu, van Melkebeek, and Morgan (2024), Theorem~1.9, give a width
bound with an additional low-degree-monomial promise; that promise
cannot be omitted from the statement.  Shpilka (2026), ECCC
TR26-314, Theorem~1.1 and Remark~1.3, gives a polynomial bound for
multilinear every-order ROABPs and their border.  It does not
assert the same bound for one hidden order.

Two 2026 preprints encountered in the current search make claims
about one unknown order.  Singh and Singh, arXiv:2602.13449v1, is
not used as an accepted theorem here: the subsequent source audit
I-22 gives explicit counterexamples to two foundational reductions.  Liu, Hexagon 2611.00634v1, supplies the rank-profile
mechanism specialized and proved in full as I-21 below.  Its
unrefereed status and disclosed AI provenance are retained in the
source ledger.  The accepted mathematical assertion in this report
is the explicitly proved specialization I-21, not an endorsement
of either preprint as a whole.

\subsection{I-21. A checked width-three specialization of the rank-preservation construction}
\label{I-21}

\paragraph{Source and status.}
The mechanism below is a specialization of Liu (2026),
\emph{Black-box identity testing for affine orbits of read-once
oblivious algebraic branching programs}, Hexagon
\texttt{2611.00634v1}, Section~3.1, Theorem~3.2 and Lemma~3.3,
and Section~3.2, Lemmas~3.4 and~3.6, printed pp.~5--7.
This is an unrefereed preprint dated November 19, 2026; its host
discloses primarily AI-generated text and no human-understood parts.
Its identifier is a Hexagon identifier, not an arXiv identifier.
We give a complete proof of the narrower original-coordinate,
width-three statement.  The affine-orbit assertions, the general-width
bounds, and the other sections of that manuscript are not being
certified here.  No historical priority is claimed.

\paragraph{Statement.}
Let $F$ be any field, let $n,D\ge1$, and let $f\in F[x_1,\ldots,x_n]$
have total degree at most $D$.  Suppose that for some permutation
$\sigma$ it has a constant-boundary representation
\[
 f=uA_1(x_{\sigma(1)})\cdots A_n(x_{\sigma(n)})v,
 \qquad A_i\in F[x_{\sigma(i)}]^{3\times3}.
\]
Smaller widths and rectangular layers are included by zero padding.
Let $E/F$ be a working field with at least
\[
 M=\max\{n,D+1,(n-1)D+1\}
\]
distinct available elements.  Choose distinct $a_1,\ldots,a_n\in E$,
and let $\lambda_i(Y)$ be their degree-$(n-1)$ Lagrange polynomials.
Define the twenty-one-parameter map
\[
 \mathcal G_i(t,\mathbf y,\mathbf z)
   =t+\sum_{j=1}^{10}z_j\lambda_i(y_j),
 \qquad 1\le i\le n.
 \tag{I.22}\label{eq:I-width3-generator}
\]
Then $f\ne0$ implies $f(\mathcal G)\ne0$ as a polynomial over $E$.
No variable order is an input to the map.

Put $R=(n-1)D+1$.  Take any sets $T,Y\subset E$ with
$|T|=D+1$ and $|Y|=R$.  For distinct
$\theta_0,\ldots,\theta_D\in E$, put
\[
 \Sigma_{10,D}=
 \{(\theta_{b_1},\ldots,\theta_{b_{10}}):
       b_j\ge0,\ \sum_j b_j\le D\}.
\]
The evaluations of $\mathcal G$ on
$T\times Y^{10}\times\Sigma_{10,D}$ form a hitting set of size at most
\[
 (D+1)\bigl((n-1)D+1\bigr)^{10}\binom{D+10}{10}.
 \tag{I.23}\label{eq:I-width3-size}
\]
This is polynomial in $n,D$ and is valid in every characteristic.
The field convention is explicit: if $E\ne F$, the black box must
allow evaluation over $E$.  No base-field query simulation is being
assumed.  In particular this does not contradict the finite-field
obstruction in I-16.

\paragraph{1. The coefficient-space lemma needed at width three.}
For a field $L$, work in
\[
 \mathcal R=L[Y_1,\ldots,Y_m]/(Y_1^{e_1+1},\ldots,Y_m^{e_m+1}),
 \qquad e_i\ge0,
\]
and suppose $B_i$ depends only on $Y_i$.  A matrix has \emph{unit rank}
$r$ if its constant matrix has rank $r$ and all its $(r+1)$-minors
vanish in $\mathcal R$.  Unit rank zero means that the matrix is zero.
For any constant row $b$, if every interval of the word
$B_1\cdots B_m$ has unit rank in $\{0,2,3\}$, then
\[
 \operatorname{Coeff}(bB_1\cdots B_m)
   =\operatorname{span}_L
       \{[Y^\alpha](bB_1\cdots B_m):
                         |\operatorname{supp}\alpha|\le4\}.
 \tag{I.24}\label{eq:I-profile23}
\]
If every interval has unit rank three, the corresponding bound is two.
Here $\operatorname{Coeff}$ means the span of coefficient rows.

We prove this assertion, including the stable-projection step.
An interval of unit rank zero makes the whole product zero.  Otherwise
the full constant product has rank two or three.  In rank three all
constant factors are invertible.  Cut-dependent constant changes of
basis make them identities.  The coefficient row spaces of the
successive prefixes are then nested.  Starting with one row, they can
grow at most twice.  At a growth layer one new nonconstant local
coefficient increases support by at most one; at a stable layer
constant propagation preserves a spanning set.  This proves the
bound two, with actual coefficient rows, as in I-17.

In rank two write the constant matrices as $C_i$.  Choose a
two-dimensional row space $U_0$ on which right multiplication by
$C_1\cdots C_m$ is injective, and set $U_i=U_0C_1\cdots C_i$.
Choose a one-dimensional complement $V_m$ of $U_m$.  At cut $i$ take
\[
 V_i=\{z:zC_{i+1}\cdots C_m\in V_m\}.
\]
The suffix has rank at least two and maps $U_i$ isomorphically onto
$U_m$.  Thus $V_i$ has dimension one, complements $U_i$, and
$V_{i-1}C_i\subseteq V_i$.  Transport bases of the $U_i$ along the
constants and choose bases of the $V_i$.  The resulting constant
matrices are
\[
 C_i=\operatorname{diag}(I_2,q_i),\qquad q_i\in L.
 \tag{I.25}
\]
These basis changes preserve supports and all interval unit ranks.

Let $\pi$ denote projection of a row onto its first two coordinates.
For a starting constant row space $S_0$, propagate
$S_i=\operatorname{Coeff}(S_{i-1}B_i)$.  Since
$\pi(sC_i)=\pi(s)$, the spaces $\pi S_i$ are nested.
Consider first a segment on which they are all the same space $T$.
Choose basis lifts of $T$ in $S_0$, with initial cost zero, and a basis
of the kernel of $\pi|_{S_0}$, with initial cost one.  Define the cost
of an obtained coefficient row as its initial cost plus the number
of nonconstant local coefficients used.  We claim that cost at most
one spans the final space on this stable segment.

If every $q_i$ on the segment is nonzero, the constants are invertible.
All the spaces lie in $\pi^{-1}(T)$, of dimension $\dim T+1$.
If the initial kernel is present, the initial space is already this
entire ambient space and constant propagation suffices.  Otherwise
there is room for at most one dimension increase.  Before that
increase the propagated constant lifts span; at the increase one
nonconstant coefficient suffices; afterwards constant propagation
again spans.  The claimed cost bound follows.

If some $q_i=0$, let $c$ be the last such position, put
$P=B_c\cdots B_m$, and let $B$ be the matrix of the initial lifts
propagated through the constants before $c$.  The suffix $P$ has
unit rank two and constant matrix $\operatorname{diag}(I_2,0)$.
Write its upper blocks as $[U\ V]$.  The constant part of $U$ is
$I_2$, so $U$ is invertible in $\mathcal R$.  The vanishing
three-by-three determinant, or its Schur complement, shows that
every row $(x,y)$ in the row module of $P$ satisfies
$y=xU^{-1}V$.

Let $H=\pi B$, whose rows form a basis of $T$.  If $T=0$, this row
graph forces every propagated row through $P$ to be zero.  Otherwise
write
\[
 \pi(BP)=QH,\qquad \pi(sP)=\eta H\quad(s\in S_{c-1}).
\]
Such expressions exist because all coefficient projections stay in
$T$, and $Q(0)=I$.  Therefore $Q$ is invertible.  The row graph gives
\[
 sP=\eta Q^{-1}BP.
 \tag{I.26}
\]
Coefficients of a polynomial matrix times $BP$ are linear
combinations of coefficient rows of $BP$, also in the truncated
ring.  Conversely the rows of $B$ belong to $S_{c-1}$.  It follows
that the final coefficient space is exactly
$\operatorname{Coeff}(BP)$.
Restart just after $c$ from $\operatorname{Coeff}(BB_c)$.
Constant lifts have cost zero and a kernel basis has cost at most
one, since only one variable position was used.  The remaining
constants are invertible; the preceding case therefore proves the
same cost bound.  This completes the stable-segment claim.

For the original single starting row $b$, first suppose $\pi b\ne0$.
The active projected space starts in dimension one and has at most
one strict growth.  The first stable segment costs at most one,
the growth factor at most one more, and the final stable segment
at most one more: support at most three suffices.  If $\pi b=0$,
the first stable segment has zero projected space.  The
stable-segment claim, with the sole initial row assigned cost one,
says that this segment is spanned by coefficients with no
nonconstant positions, or is zero.  There can subsequently be at
most two projection-growth factors.  Each costs at most one, and
each following stable segment at most one.  Support at most four
therefore suffices.  At every restart, lifts and kernel vectors
are linear combinations of the coefficient rows already retained;
the span filtration by support is preserved.  This proves (I.24).

The bound four in (I.24) is sharp.  With $P=\operatorname{diag}(1,1,0)$,
take
\[
 B_1=P+Y_1E_{31},\quad B_2=I+Y_2E_{13},\quad
 B_3=P+Y_3E_{32},\quad B_4=I+Y_4E_{23},\quad b=e_3^T.
\]
Every interval containing $B_1$ or $B_3$ has generic and constant
rank two, and the other intervals are invertible.  The output row is
\[
 Y_1e_1^T+Y_1Y_2Y_3e_2^T+Y_1Y_2Y_3Y_4e_3^T.
\]
Its third coefficient direction requires support four.

\paragraph{2. Lowest homogeneous components.}
Suppose every interval of a polynomial matrix word has the same
generic rank as at zero, with these ranks in $\{2,3\}$.  For constant
endpoints $b,c$, every monomial in the lowest nonzero homogeneous
component of $bB_1\cdots B_mc$ has support at most four.  If all
interval ranks are three, the bound is two.

Indeed, choose such a monomial $Y^\alpha$ and reduce modulo
$Y_i^{\alpha_i+1}$ for every $i$.  Only monomials dividing
$Y^\alpha$ survive.  Lowest-degree minimality makes the surviving
scalar polynomial exactly its nonzero $Y^\alpha$ term.  All the
intervals retain their unit ranks: the pivot constant minor remains
invertible and the higher minors were already zero polynomially.
Apply (I.24), or its invertible version.  If the support of $\alpha$
were too large, the final boundary $c$ would annihilate every
spanning coefficient row and hence also the surviving one, a
contradiction.

\paragraph{3. A finite collection of local rank tests.}
Augment the original matrix word by the constant factors $e_1u$
and $ve_1^T$.  Its full product has only one possibly nonzero entry,
namely $f$ in position $(1,1)$.  All intervals have a well-defined
generic rank over the rational function field.  Translation of
each distinct variable by any center in an overfield preserves
these generic ranks.

Fix $r=1$ or $r=2$, and suppose all intervals of generic rank
greater than $r$ are already preserved at the current center.
For any target interval $[l,h]$ of rank $r$, let $[l,i_0]$ be its
first prefix of rank $r$ and require this prefix to be preserved.
For each $i_0\le i<h$, choose the shortest rank-$r$ suffix $[j,i]$
of $[l,i]$, and require $[j,i+1]$ to be preserved.  Each such test
has generic rank $r$.  Its interior $[j+1,i]$, and the prefix
$[l,i_0-1]$ of the first test, are empty or have generic rank
greater than $r$ by minimality.  Thus a test has one such chunk
and at most one factor at each end.

These tests suffice.  Once $[l,i]$ is preserved, its specialized
row space is contained in that of $[j,i]$.  Both have dimension
$r$: the first has rank $r$, and the second cannot exceed its
generic rank $r$.  Their row spaces are equal.  Multiplying by
the next specialized factor and using the tested rank of
$[j,i+1]$ proves preservation of $[l,i+1]$.  This proves the
assertion by induction, for all target intervals.  Select one
nonzero $r$-minor of each test; no efficient selection or knowledge
of the hidden word is needed in the black-box construction.

For $r=2$, each selected minor is an entry of the word of second
compound matrices, by Cauchy--Binet.  These compounds have size
three.  The interior chunk and every one of its subintervals have
rank three, and so do their constants.  In a lowest-degree
monomial of the selected minor, extract the prescribed powers at
the at most two end factors.  The remaining scalar chunk
contraction is still in its lowest degree.  Part~2 gives support
at most two inside the chunk.  The whole monomial consequently
has support at most four.

For $r=1$ the compounds are the original matrices.  The chunk and
its subintervals have ranks in $\{2,3\}$, already preserved at
the center.  The same argument gives support at most four
inside the chunk and at most six in the whole selected entry.
Empty chunks and constant end factors only reduce these bounds.

\paragraph{4. Simultaneous preservation with four then six pairs.}
For $k\ge1$ write
\[
 \operatorname{SV}_{k,i}(\mathbf y,\mathbf z)
      =\sum_{j=1}^kz_j\lambda_i(y_j).
\]
Over any coefficient overfield of $E$, this map detects any
nonzero polynomial having a nonzero monomial of support at most
$k$: set the selectors to the Lagrange nodes of that support,
and the unused amplitudes to zero.  The resulting restriction
to that coordinate subspace retains the chosen coefficient.
In particular it detects a nonzero homogeneous polynomial all
of whose monomials have support at most $k$.

Start at the common center $(t,\ldots,t)$ over $E(t)$.
An interval of generic rank three consists entirely of invertible
univariate factors.  Its determinant at the center is a product
of nonzero univariate determinant polynomials evaluated at the
transcendental $t$, and is therefore nonzero.  Thus all rank-three
intervals are preserved simultaneously.

Add an $\operatorname{SV}_4$ perturbation with fresh independent
selectors and amplitudes.  For every chosen rank-two test minor,
its expansion about the current center is a nonzero polynomial
whose lowest homogeneous component $H$ has support at most four
in each monomial, by part~3.  Hence $H(\operatorname{SV}_4)\ne0$.
The map is homogeneous of degree one in its amplitudes.  The
lowest amplitude degree in the full substituted minor is
therefore exactly the degree of $H$, and its component is
$H(\operatorname{SV}_4)$; higher homogeneous terms cannot cancel
it.  Each selected test minor is consequently nonzero in the
same rational function field of the new parameters.  This is
simultaneous preservation, not a separate numeric choice for
each test: finitely many nonzero polynomials have nonzero product
in the parameter polynomial ring.  Earlier preserved minors
remain nonzero because setting the new amplitudes to zero
recovers their previous values.  The local-test induction now
preserves every interval of ranks two and three.

Over that enlarged rational function field, add an independent
$\operatorname{SV}_6$ perturbation.  The rank-one test components
have support at most six by part~3, so the same argument preserves
all these tests and all earlier ranks.  The augmented full word
has generic rank one whenever $f\ne0$.  Its specialized rank is
therefore one, which says precisely that the resulting scalar is
nonzero.  The two perturbations use disjoint parameter pairs;
their sum is $\operatorname{SV}_{10}$.  This proves (I.22).
All uses of the variable order occurred only in this existential
proof about the hidden representation.  The actual map is the
same for every order.

For comparison with the failed repair I-18, the polynomial
$\prod_{j=1}^m(x_{2j}-x_{2j-1})$ already survives one interpolation
pair: its pullback is
$z^m\prod_j(\lambda_{2j}(y)-\lambda_{2j-1}(y))$, which is nonzero
because the distinct Lagrange basis polynomials have nonzero
differences.  Interpolation directions are generally dense;
the construction therefore does not require a sparse perturbation
point to witness every nonzero polynomial.

\paragraph{5. Finite sampling and complexity.}
Since $f$ has total degree at most $D$, its pullback under (I.22)
has degree at most $D$ in $t$, at most $(n-1)D$ in each selector
$y_j$, and total degree at most $D$ in all ten amplitudes together.
Repeated univariate interpolation therefore leaves a nonzero
amplitude polynomial at some point of $T\times Y^{10}$.

The set $\Sigma_{10,D}$ detects every nonzero amplitude polynomial
of total degree at most $D$, over any characteristic.  To see this,
use the Newton basis
\[
 N_{\mathbf b}(\mathbf z)
  =\prod_{j=1}^{10}\prod_{h=0}^{b_j-1}(z_j-\theta_h),
 \qquad \sum_jb_j\le D.
\]
Its leading monomials are exactly those of total degree at most
$D$.  Evaluation at the node indexed by $\mathbf c$ vanishes
unless $c_j\ge b_j$ for every $j$.  Ordered by total degree,
the evaluation matrix is triangular with nonzero diagonal.
It is invertible, proving the claimed detection and (I.23).
Given the field nodes, computing the Lagrange polynomials and
enumerating the displayed family uses polynomially many field
operations in $n,D$, since its ten-pair parameter count is fixed.
Over a finite field, one may use an explicitly represented
extension large enough to supply the stated nodes.  Over
$\mathbb Q$ the nodes can be chosen as integers of polynomial
magnitude, giving rational coordinates of polynomial bit length.
\hfill$\square$

\paragraph{What this closes, and what it leaves open.}
This supplies a concrete positive answer for the width-three,
one-hidden-order case of I-G02, including arbitrarily many
rank-two cuts.  The rank-cut expression in I-19 equals
$f\prod_j\Delta_j$; the extra univariate factors can be absorbed
into the width-three layers, so the same theorem also hits those
particular coupled expressions with their increased degree bound.
The construction uses rank preservation directly and does not
derive a BIWA or a Hoza cocharacter decoder.
The general polynomial-in-$w$ gate I-G01 remains open in this
report.  In particular, I-08 requires widths growing with the
matrix dimension; this fixed-width result does not provide that
polynomial Noether-normalization conclusion.


\subsection{I-22. Two explicit failures of the active-simple-block reduction}
\label{I-22}

\paragraph{Audited claim and scope.}
Singh and Singh (2026), arXiv:2602.13449v1, Theorem~7.4
(printed p.~38), claims the unknown-order bound
$n(wd)^{O(w^2)}$.  Its foundational reductions are false as stated:
the equivalence after Definition~2.1 (p.~4), repeated in
Section~3.1 (p.~6), and the active-block conclusions in
Lemma~3.1 (pp.~6--7), Remark~3.1.4 (p.~9), and Lemma~3.2
(p.~10).  The following counterexamples work over $\mathbb C$,
inside the source's characteristic-zero hypothesis.  They
invalidate those proof steps; they do not disprove the headline
existence of a hitting set with the claimed bound.

\paragraph{Full word algebra loses the prescribed ordered output.}
Let $N=E_{12}$ and $M=E_{21}$, and take the width-two factors
\[
 A_1(x_1)=M+x_1N,\qquad A_2(x_2)=M+x_2N.
\]
Since $N^2=M^2=0$, $NM=E_{11}$, and $MN=E_{22}$,
\[
 A_1(x_1)A_2(x_2)=x_1E_{11}+x_2E_{22}.
\]
The specified ROABP output $e_1^TA_1A_2e_2$ is identically zero.
Nevertheless the algebra generated by the coefficient matrices
is the full $\operatorname{Mat}_2(\mathbb C)$, and
$e_1^TNe_2=1$.  Thus vanishing on the actual ordered products
of one local coefficient per layer does not imply vanishing on
all words in the generated algebra.

This example also separates a nonzero simple-block image of the
matrix product from a nonzero prescribed scalar output.  The
algebra is already simple, so its identity representation is a
simple-block projection and the displayed matrix product survives
it, although the scalar output is zero.  After distinct nonzero
specializations of $x_1,x_2$, the two specialized factors still
generate the full matrix algebra: their difference gives a
nonzero multiple of $N$, and then either factor gives $M$.
Hence moving to a generic multiplicative specialization does not
repair this example.

\paragraph{A nonzero output can live entirely in the radical.}
Take four width-two factors
\[
 A_1=x_1E_{11},\quad A_2=I+x_2N,\quad
 A_3=I+x_3N,\quad A_4=x_4E_{22}.
\]
Direct multiplication gives
\[
 W=A_1A_2A_3A_4=x_1x_4(x_2+x_3)N,\qquad
 f=e_1^TWe_2=x_1x_4(x_2+x_3)\ne0.
 \tag{I.27}
\]
The coefficient algebra is the algebra $\mathcal T_2$ of upper
triangular two-by-two matrices.  Its Jacobson radical is
$J=\mathbb C N$, with $J^2=0$, and its semisimple quotient is
$\mathbb C\oplus\mathbb C$, given by the two diagonal entries.
For completeness, $\mathbb C N$ is a square-zero ideal, hence
lies in the radical; the quotient is a product of fields and
has zero radical, proving equality.  Every homomorphism onto
a simple quotient therefore kills $W$.  No active simple block
of the kind required by the source exists for this nonzero
computation.

All four variables occur, the product length is $4\ge w=2$,
and the example is not a width-one scalar product in disguise:
the coefficient matrix of $x_2+x_3$, with rows and columns
indexed by exponents $0,1$, is
$\left(\begin{smallmatrix}0&1\\1&0\end{smallmatrix}\right)$,
of rank two.  The same obstruction persists after any nonzero
character values for which $x_2+x_3\ne0$; the specialized
algebra remains $\mathcal T_2$ and the product remains a
nonzero radical element.  Constant identity factors can be
inserted to make the product arbitrarily long.

Nilpotency $J^w=0$ controls products containing $w$ radical
factors.  It does not imply that a long product of arbitrary
algebra elements vanishes whenever the final product belongs
to $J$.  Equation~(I.27) gives precisely this distinction.
These two counterexamples are enough to withhold the cited
Theorem~7.4 from the accepted background of this report.
They leave the independently proved I-21 unaffected.
\hfill$\square$


\subsection{I-23. A binary-code cost for preserving decomposable exterior states}
\label{I-23}

\paragraph{The specified next construction.}
One proposed repair is to keep the state of an exterior-power lift
on its Grassmann cone by requiring each Hoza seed to act diagonally
on the Pl\"ucker coordinates in a decomposability-preserving way.
We analyze exactly this requirement.  The statements concern
formal algebraic one-parameter actions over a field, not just
their values on the finitely many elements of a finite field.

Fix $2\le r\le w-2$, index Pl\"ucker coordinates by the
$r$-subsets $I\subseteq[w]$, and let $\rho_I\in\mathbb Z$.
The diagonal action $p_I\mapsto t^{\rho_I}p_I$ preserves all
decomposable $r$-vectors if and only if
\[
 \rho_I=c+\sum_{i\in I}a_i
 \quad\hbox{for some }c,a_1,\ldots,a_w\in\mathbb Z.
 \tag{I.28}\label{eq:I-plucker-additive}
\]
Consequently, if a binary code
$C:\binom{[w]}r\to\{0,1\}^{\ell}$ has minimum distance $d$
and every seed of the Hoza weighting
\[
 \rho_s(I)=\sum_{j=1}^{\ell}6^{\ell-j}C(I)_js_j,
 \qquad s\in\{0,1\}^{\ell},
\]
preserves decomposability, then
\[
 \ell\ge\frac{wd}{2}.
 \tag{I.29}\label{eq:I-plucker-code-cost}
\]
The additional layer multiplier $6^{\nu_2(i)}$ in I-B01
does not change this conclusion.

\paragraph{Proof of the cocharacter classification.}
Let $K$ have size $r-2$, and let $a,b,c,d$ be four different
indices outside $K$.  The Pl\"ucker relation on those indices
forces
\[
 \rho_{Kab}+\rho_{Kcd}
 =\rho_{Kac}+\rho_{Kbd}
 =\rho_{Kad}+\rho_{Kbc}.
\]
Here subscripts denote sets, with the usual orientation signs
in the coordinate relation.  To verify the force of this
condition directly, restrict to planes containing the fixed
coordinate $(r-2)$-plane on $K$.  In the remaining two rows,
use four columns $(e_1,e_2,e_1,e_2)$, then
$(e_1,e_1,e_2,e_2)$.  The relation has exactly two nonzero
terms in each case.  After the diagonal action their two powers
of the indeterminate $t$ must agree.  This gives the displayed
equalities, in every characteristic.

For distinct $i,j$, consider
$\rho_{S\cup\{i\}}-\rho_{S\cup\{j\}}$ with
$|S|=r-1$ and $S\cap\{i,j\}=\varnothing$.
The four-index equality shows that this difference is unchanged
when one element of $S$ is exchanged.  The graph of such subsets
under single exchanges is connected, so the difference is a
constant $d_{ij}$ independent of $S$.  A common $S$ disjoint
from three specified indices exists because $r\le w-2$;
it gives $d_{ij}+d_{jk}=d_{ik}$.  Take $a_1=0$ and
$a_i=d_{i1}$.  Then
$\rho_I-\sum_{i\in I}a_i$ is unchanged by every single exchange
of an $r$-subset.  All $r$-subsets are connected by such
exchanges, so this is one integer $c$.  This proves necessity
of (I.28).  Conversely (I.28) is the exterior action induced by
$\operatorname{diag}(t^{a_1},\ldots,t^{a_w})$, followed by
one common scalar, and plainly preserves decomposable vectors.

\paragraph{Binary-valued additive functions are individual membership bits.}
Suppose $b(I)\in\{0,1\}$ has the form (I.28).
Every exchange difference $a_i-a_j$ lies in $\{-1,0,1\}$.
If $b$ is nonconstant, the $a_i$ therefore have exactly two
levels differing by one.  Let $S$ index the higher level.
Apart from an additive constant, $b(I)=|I\cap S|$.
If both $|S|$ and $w-|S|$ are at least two, then among
$r$-subsets this intersection count ranges from
$\max(0,r-(w-|S|))$ to $\min(r,|S|)$, a range of length
at least two when $2\le r\le w-2$.  That contradicts binary
values.  Hence $|S|=1$ or $w-|S|=1$.
The binary values then force $b(I)$ to be either membership
of one element or its complement.  Constant bits are the only
remaining possibilities.

Apply this classification to the seed $s=e_j$.  The positive
integer factor $6^{\ell-j}$ does not affect the required
exchange equalities, so each code bit is a constant or a single
membership/complement test.  Let $m_i$ count the nonconstant
bits testing element $i$.  Bases differing only by exchange of
$i$ and $j$ have Hamming distance $m_i+m_j$.  Such bases exist
for every pair.  Thus $m_i+m_j\ge d$ for every $i\ne j$.
Summing over all pairs gives
\[
 (w-1)\sum_i m_i\ge\binom w2d,
 \qquad \ell\ge\sum_i m_i\ge wd/2.
\]
This proves (I.29). \hfill$\square$

\paragraph{Consequence for this repair.}
On the $W=\binom wr$ exterior states, using the stated Hoza
error bound $2^{-d}W^4N^2\le1/2$ requires
$d\ge\lceil\log_2(2W^4N^2)\rceil$.
The all-seeds decomposability requirement then forces
$\ell=\Omega(w\log(WN))$ and at least
$(2W^4N^2)^{w/2}$ seeds.  Already $r=2$ gives
$W=\Theta(w^2)$, so this loses the polynomial family size
without invoking an exponentially large exterior dimension.

The restriction is essential: this is not a lower bound for
all generators or all nonlinear decodings of Hoza seeds.
It does not rule out a separately proved successful subset
of coherent seeds from another code, or a construction that
leaves the Grassmann cone at intermediate stages and restores
the required coefficient information afterwards.  Nor would
mere Pl\"ucker compatibility by itself resolve the distinct
variable-quotient condition in I-02.


\subsection{I-24. Exponential support is possible inside the preserved-rank test class}
\label{I-24}

\paragraph{Statement.}
For every $w\ge3$, there is a word of $w\times w$ univariate
linear matrices on
\[
 h_w=3\cdot2^{w-2}-2
 \tag{I.30}\label{eq:I-reset-support}
\]
distinct variables with the following properties.  Every
nonempty interval has the same generic rank as its constant
matrix, and that rank is at least two.  Nevertheless one scalar
contraction of the full word is exactly the product of all
$h_w$ variables.  Its lowest homogeneous component therefore
has support $h_w$.  With constant rank-one factors appended
at both ends it is an actual final rank-one test of the
rank-descent construction used in I-21.

Thus a polynomial-in-$w$ bound on the support of every such
lowest component is false, even for $r=1$, where every nonzero
exterior state is automatically decomposable.  This is a
counterexample to that particular compression proposal;
it is not a hitting-set size lower bound, since a monomial is
easily hit by a point with all coordinates nonzero.

\paragraph{Construction.}
Let $P_r=\operatorname{diag}(I_r,0_{w-r})$ for $2\le r<w$.
For $a\le r$ define a two-position gadget, on fresh variables,
\[
 Q_{r,a}(s,t)=(P_r+sE_{w,a})(P_r+tE_{a,w}).
\]
Each of its factors has constant and generic rank exactly $r$:
its leading $r\times r$ block is $I_r$, and the extra row or
column is a multiple of a row or column already present.
All the constant matrices $P_r$ are nested projections.
Consequently any interval of the eventual word has constant
product $P_{\min r_i}$ and generic rank at most the minimum
rank of its factors.  The nonzero constant minor gives the
reverse inequality.  Every interval therefore has generic
rank exactly its constant rank $\min r_i\ge2$.
The same assertion is a unit-rank statement after any separate
truncations of the variables.

Define words $H_k$, using fresh variable copies on every
occurrence, by
\[
 H_w=\text{empty word},\qquad
 H_k=H_{k+1}\,Q_{k,k}\,H_{k+1}
       \quad(3\le k<w).
\]
The final word is
\[
 \mathcal W_w=H_3\,Q_{2,1}\,H_3\,Q_{2,2}\,H_3,
\]
with all five displayed blocks using disjoint fresh variables.
If $L_k$ is the number of positions in $H_k$, then
$L_w=0$ and $L_k=2L_{k+1}+2$, whence
$L_k=2(2^{w-k}-1)$.  Thus $\mathcal W_w$ has
$3L_3+4=h_w$ positions.

\paragraph{Proof of the scalar identity.}
If an input row $q$ has $q_a=0$, then the two-position gadget
acts by
\[
 qQ_{r,a}(s,t)
   =\sum_{j\le r}q_je_j^T+q_ws\,e_a^T+q_wst\,e_w^T.
 \tag{I.31}
\]
In particular it erases coordinates strictly between $r$ and
$w$, leaves the earlier coordinates unchanged, and transfers
the last coordinate through both fresh variables.

Inductively, suppose an input row has coordinates
$k,\ldots,w-1$ zero and last coordinate $c$; its earlier
coordinates may be arbitrary.  The word $H_k$ leaves those
earlier coordinates unchanged and makes its last coordinate
$c$ times the product of every variable in $H_k$.
This is immediate for $H_w$.  For $H_k$, the first
$H_{k+1}$ leaves coordinate $k$ zero.  Formula~(I.31) for
$Q_{k,k}$ multiplies the last coordinate by its two variables
and resets coordinates $k+1,\ldots,w-1$ to zero.  The second
$H_{k+1}$ therefore applies under the same induction
hypothesis.  This proves the claim.

Start the full word at $e_w^T$.  After the first $H_3$,
coordinates one and two remain zero.  The gadget $Q_{2,1}$
uses the zero first coordinate and resets all coordinates
$3,\ldots,w-1$.  The second $H_3$ preserves coordinate two
as zero, so $Q_{2,2}$ again satisfies (I.31) and resets the
higher coordinates.  The final $H_3$ completes the induction.
Therefore
\[
 e_w^T\mathcal W_we_w=\prod_{i=1}^{h_w}Y_i.
\]
This argument is valid over every field; there are no signed
cancellations or characteristic restrictions.

\paragraph{Membership in the rank-test class.}
Append $L=e_1e_w^T$ on the left and $R=e_we_1^T$ on the
right.  The full product is the displayed nonzero monomial
times $e_1e_1^T$, so it has generic rank one.  Every prefix
beginning with $L$ is nonzero, since the last row of each
partial gadget has a surviving coordinate; consequently these
prefixes have rank one.  Every interval lying wholly inside
$\mathcal W_w$ has rank at least two, as already proved.
For the prefix ending just before $R$, the shortest rank-one
suffix therefore starts at $L$.  Appending $R$ gives exactly
the full word as the last test in part~3 of I-21.
Its interior chunk satisfies every higher-rank preservation
hypothesis at the zero center.  The exponential support cannot
be attributed to failure to preserve one of those interval
ranks. \hfill$\square$

At $w=3$ this yields support four; at $w=4$ it yields support
ten; at $w=5$ it yields support twenty-two.  The construction
shows why a small intrinsic state-variety dimension alone does
not bound how many variable positions are needed after many
singular resets.  A useful Hoza repair would have to exploit
more than low support or decomposability: it must handle the
original-variable coefficients and their cancellations while
allowing such long supports.


\subsection{I-25. Intrinsic Grassmannian compression succeeds for invertible chunks}
\label{I-25}

The nested-image method is suggested by Liu (2026), Section~3.3,
printed pp.~7--8.  The following restricted refinement uses
the projective Grassmannian of one decomposable exterior row
as its ambient variety and is proved directly.

\paragraph{Statement.}
Let $A_i(Y_i)$ be $w\times w$ matrices over a field $L$, in
distinct variables, with every $A_i(0)$ invertible.  Fix a
nonzero decomposable row $p\in\bigwedge^r L^w$ and a linear
functional $\ell$ on $\bigwedge^rL^w$, where $1\le r<w$.
If
\[
 q(\mathbf Y)=\ell\bigl(p(\textstyle\bigwedge^r A_1(Y_1))
                  \cdots(\textstyle\bigwedge^r A_m(Y_m))\bigr)
 \ne0,
\]
then there is a set $S\subseteq[m]$ with
$|S|\le r(w-r)$ for which setting $Y_i=0$ outside $S$
keeps $q$ nonzero.  In particular $q$ has a nonzero monomial
of support at most $r(w-r)$, and $\operatorname{SV}_{r(w-r)}$
detects it.

For an $r$-minor test with an invertible-constant interior
chunk and at most one factor on either end, the whole minor
therefore has some nonzero monomial of support at most
$r(w-r)+2$.  This replaces the arbitrary compound-space bound
$\binom wr-1+2$ for this restricted class.  It bounds the
existence of a low-support monomial, not every monomial in the
lowest homogeneous component.

\paragraph{Proof.}
Use cut-dependent invertible constant changes of basis to make
every $A_i(0)=I$.  These changes arise from $\operatorname{GL}_w$
and hence preserve decomposability after taking exterior powers;
they adjust only the fixed endpoints.  Extend scalars to an
algebraic closure.  This does not change whether a polynomial
or its coordinate restriction is zero.

Restrict each parameter to the nonempty open set where its
factor is invertible.  Let $Z_i$ be the Zariski closure in
$\operatorname{Gr}(r,w)$ of the projective prefix states
\[
 [p(\textstyle\bigwedge^r A_1(Y_1))\cdots
        (\textstyle\bigwedge^r A_i(Y_i))].
\]
Each $Z_i$ is irreducible, as the closure of the image of an
irreducible product of open affine lines.  The value $Y_i=0$
is allowed and acts as the identity, so
$Z_{i-1}\subseteq Z_i$.  The starting variety $Z_0$ is one
point.  The Grassmannian has dimension $r(w-r)$: its dense
matrix chart consists of row spaces of $[I_r\ X]$ with
$X\in L^{r\times(w-r)}$.
Thus at most $r(w-r)$ of these inclusions are strict.

Whenever $Z_i=Z_{i-1}$, replace $A_i(Y_i)$ by $I$.
This preserves the closure of the prefix image and hence the
closure of every later image.  To justify the latter assertion,
on invertible factor values the exterior action is a regular
map on the projective Grassmannian; the closure of the image
under the next factor of a dense prefix set is the same as
the closure obtained from the whole prefix variety.  This
persists on composing successive factors.  Inductively one
may therefore remove every non-growth position while retaining
the final projective image closure.  At most $r(w-r)$ positions
remain.

The scalar $q$ vanishes identically exactly when that final
projective image lies in the hyperplane cut out by $\ell$.
It suffices to test this on the invertible-value open set,
since a polynomial vanishing there vanishes identically.
Preserving the image closure therefore preserves nonvanishing
of $q$.  A nonzero restriction to the retained coordinates
contains a monomial supported there, proving the first claim.

For the endpoint version, adjoin the at most two endpoint
variables to the coefficient field.  A selected row of an
$r$th compound endpoint is the wedge of $r$ rows of the
original endpoint matrix, hence is decomposable; it is nonzero
if the selected full minor is nonzero.  The right endpoint
supplies a linear functional.  The preceding argument applies
to the interior chunk over this rational function field.
The resulting restricted minor is still a polynomial in the
original variables and is nonzero there.  Restoring the at
most two endpoint positions gives the asserted support bound.
\hfill$\square$

\paragraph{Exact limitation.}
This construction exploits the exterior geometry instead of
treating $\bigwedge^r L^w$ as an arbitrary vector space of
dimension $\binom wr$.  It is an existence-and-enumeration
compression: the retained positions depend on the matrices,
and enumerating them or using the stated SV map still has an
exponent depending on $w$.  The proof also requires invertible
constant interior factors.  I-24 shows that replacing this
condition merely by preservation of all interval ranks at
least two cannot yield a polynomial-in-$w$ support bound.
The next useful rank-descent construction must accommodate
those singular resets and the original-variable cancellation
relations simultaneously; intrinsic dimension alone does
not supply that extension.


\subsection{Open gate I-G01 and the order of further work}
\label{I-G01}

The surviving main gate is an explicit polynomial family of
nonnegative integer cocharacters of polynomial magnitude which, for every
width-$w$ matrix product with disjoint univariate factors,
enters at least one cone (I.3), or satisfies the exact weaker joint
coefficient-span condition I-14 for the family.  If this is
proved, I-B02 or I-14 gives polynomial hitting sets, and I-08 gives the
precisely delimited matrix-invariant conclusion.  No finding
above claims that this gate is closed.

The constructions already attempted and resolved as stated are:
direct state weights (I-02); monomial-alphabet lifting followed
by exact evaluation interpolation (I-06); positive Gram basis
enumeration (I-05); independent and correlated Haar repair
(I-07); exterior/permutation expansion and literal subspace
configuration graphs (I-11); and positive coefficient-norm
transfer matrices (I-09), exact universal energy cubature (I-13),
and the arithmetic support test (I-12).  The further repairs now
resolved are joint span preservation (I-14), bounded affine mask
decoding (I-15), exact finite-field row states (I-16), the restricted
invertible-factor sampler (I-17), and epsilon regularization with
its necessary special-fiber condition (I-18).  Rank cuts give the exact
block reduction and coupled gate I-19/I-G02, and bounded total corank
has the explicit positive sampler I-20.  The subsequent I-21
rank-preservation construction resolves the original-coordinate
width-three case with one hidden order, in all characteristics
under its explicit working-field convention.  The positive result from
direct graph isolation is I-01.  The specified next constructions have now been tested further:
I-23 proves the code-length cost of making every Hoza seed preserve
decomposable exterior states; I-24 gives exponential support even
inside preserved-rank test chunks; I-25 supplies the valid intrinsic
Grassmannian compression for invertible chunks.  None is a general
PIT lower bound.  A next construction may use dense interpolation
directions to bypass low-support bounds, but must prove nonvanishing
for the actual rank-test coefficients and retain polynomial dependence
on variable width as well as on the number of variables.

The later refinements sharpen this status.  I-26 gives a centered dense-SV
annihilator family and a logarithmic necessary pair count; I-B08 records
the known matching order of pair count for the every-order subclass.
I-27 rules out filtering a short actual Hoza code into coherent exterior
weights.  I-28 and I-30 reduce the general width-three sampler to six
pairs.  I-29 gives a sharp two-rank support bound and a polynomial pair
budget with unbounded cuts at one singular rank.  I-31 proves an
original-coordinate sampler for every fixed width, with pair count
controlled by the number of rank levels and at most
$2^{O(w\log w)}$ for all ranks.  Its interpolation size remains far
from polynomial in variable width.  Thus the fixed-width obstruction
is resolved in this track, while I-G01's genuinely polynomial
family-size and weight-size requirements remain unproved.


\subsection{I-26. Genuine annihilators of the centered dense interpolation map}
\label{I-26}

\paragraph{Relation to prior work.}
Studying the vanishing ideal of the SV generator has an established
determinantal approach.  Hu, van Melkebeek, and Morgan (2024),
Theorem~1.3 (printed p.~6) and Corollary~1.5 (printed p.~7), give explicit
multilinear degree-$(k+1)$ annihilators of $\operatorname{SV}_k$
on $2k+1$ variables, settling the degree question discussed by
Fournier and Korwar (2018, WACT slides~13--19).
Their Theorems~1.9--1.10 retain a low-degree-monomial hypothesis
in the ROABP hitting/lower-bound assertion.  The construction
below follows this determinantal perspective, explicitly adds
invariance under the common translation, and proves an
invertible-factor ROABP representation.  No priority is claimed
for a vanishing-ideal construction in this setting.

\paragraph{Statement.}
Let $k\ge1$, let $n\ge3k+2$, and let a field $F$ contain
distinct Lagrange nodes $a_1,\ldots,a_n$.  For those nodes there
is an explicitly constructible nonzero homogeneous multilinear
polynomial $P\in F[X_1,\ldots,X_n]$ of degree $k+1$ such that
\[
 P\bigl(t\mathbf1+\operatorname{SV}_k(\mathbf y,\mathbf z)\bigr)=0.
 \tag{I.32}\label{eq:I-centered-sv-zero}
\]
It has a constant-boundary ROABP representation in every order
of width at most
\[
 W_k=\sum_{j=0}^{k+1}\binom{k+1}{j}^{\!2}
       =\binom{2k+2}{k+1},
 \tag{I.33}\label{eq:I-centered-sv-width}
\]
with commuting factors $I+X_iN_i$, where $N_i^2=0$.
Every interval of this representation is invertible, and its
constant matrix is the identity.  Thus, after constant
rank-one endpoints are added, $P$ is a genuine final rank-one
test whose interior satisfies all higher-rank preservation
requirements at zero.

In particular a centered $\operatorname{SV}_{k(w)}$ generator
valid for all widths at most $w$ and all $n$, over fields or
extensions supplying the required distinct nodes, must have
$W_{k(w)}>w$.  Equivalently, for each proposed $k$ the
counterexample requires an admissible $n\ge3k+2$.
Since $W_k\le4^{k+1}$, this implies
$k(w)=\Omega(\log w)$.  It does not rule out a logarithmic
number of pairs, and is not a lower bound for other generators.
All the general assertions hold in every characteristic under
the stated distinct-node hypothesis.

\paragraph{A determinant that removes the common center.}
Set $m=k+1$ and let
\[
 U_{i\ell}=a_\ell^i,\qquad 0\le i\le k,\quad1\le\ell\le n.
\]
In the $2k+2$ dimensional vector space of polynomials of degree
at most $2k+1$, impose the $k+1$ linear conditions
\[
 \sum_{\ell=1}^n a_\ell^i q(a_\ell)=0,
       \qquad0\le i\le k.
\]
Their common kernel has dimension at least $k+1$.
Choose linearly independent polynomials
$q_1,\ldots,q_m$ in that kernel by Gaussian elimination, and
put $V_{j\ell}=q_j(a_\ell)$.  Then $UV^T=0$.
Evaluation on the $n$ distinct nodes is injective on polynomials
of degree at most $2k+1$, since $n>2k+1$; hence $V$ has row
rank $m$.

Define
\[
 B(\mathbf X)=U\operatorname{diag}(X_1,\ldots,X_n)V^T,
 \qquad P(\mathbf X)=\det B(\mathbf X).
 \tag{I.34}
\]
Every maximal minor $U_S$ is a nonzero Vandermonde determinant.
At least one maximal minor $V_S$ is nonzero because $V$ has
full row rank.  Cauchy--Binet gives
\[
 P(\mathbf X)=\sum_{|S|=m}\det(U_S)\det(V_S)
                                  \prod_{\ell\in S}X_\ell.
\]
The nonzero minor of $V$ therefore supplies a nonzero coefficient.
This proves nonvanishing, homogeneity, and multilinearity in
every characteristic.

Because $UV^T=0$, one has $B(\mathbf X+t\mathbf1)=B(\mathbf X)$.
For every $i,j$, the polynomial $Y^i q_j(Y)$ has degree at most
$3k+1\le n-1$.  Exact Lagrange reproduction consequently gives
\[
 \sum_{\ell=1}^n a_\ell^i q_j(a_\ell)\lambda_\ell(Y)
    =Y^i q_j(Y).
\]
It follows that
\[
 B(\operatorname{SV}_k)
   =\sum_{s=1}^k z_s
       (1,y_s,\ldots,y_s^k)^T
       (q_1(y_s),\ldots,q_m(y_s)).
\]
This matrix is a sum of $k$ rank-one matrices but has size
$m=k+1$, so its determinant is identically zero.  The common
translation has no effect on $B$, proving (I.32).

\paragraph{An invertible, any-order ROABP of the stated width.}
For a generic $m\times m$ matrix $Z$, form the column vector
$\mathcal M(Z)$ of all square minors, including the empty
minor $1$.  There are $W_k$ such minors by the elementary
Vandermonde sum in (I.33).  They are linearly independent as
polynomials: distinct row/column index sets have distinct
row and column multidegrees, and each individual determinant
is nonzero.

For fixed vectors $u,v$ and an indeterminate $s$, expansion
of a minor under a rank-one update gives
\[
 \det(Z_{I,J}+s u_Iv_J^T)
 =\det Z_{I,J}
  +s\sum_{a\in I,\ b\in J}
        (-1)^{\operatorname{pos}_I(a)+\operatorname{pos}_J(b)}
        u_av_b\det Z_{I\setminus\{a\},J\setminus\{b\}}.
\]
Terms using two update columns vanish, since those columns are
proportional.  Thus there is a constant matrix $N(u,v)$ such that
\[
 \mathcal M(Z+suv^T)=(I+sN(u,v))\mathcal M(Z).
\]
Translation twice in the same direction and linear independence
of the minor polynomials imply $N(u,v)^2=0$ by comparing the
coefficient of two independent translation parameters.  Two
different rank-one translations commute, so their representing
matrices commute as well.  These arguments compare polynomial
identities and do not divide by factorials.

Let $u_\ell$ and $v_\ell$ be the $\ell$th columns of $U$ and
$V$, and set $N_\ell=N(u_\ell,v_\ell)$.  Starting at $Z=0$
and translating in these $n$ rank-one directions gives
\[
 \mathcal M(B(\mathbf X))
     =\prod_{\ell=1}^n(I+X_\ell N_\ell)\,\mathcal M(0).
\]
The coordinate of the full determinant is $P$.  The factors
commute and satisfy
$(I+X_\ell N_\ell)^{-1}=I-X_\ell N_\ell$.
This gives the advertised width-$W_k$ ROABP in any variable
order, with invertible factors and identity constants.
The starting row and final column may be taken as the
full-determinant selector and the empty-minor selector.
Adding their rank-one endpoint matrices produces an actual
rank-one test exactly as in I-21.  Every interior interval
has generic and constant rank $W_k$, so no missing rank
preservation hypothesis accounts for this annihilation.
\hfill$\square$

\paragraph{A width-four, one-pair example.}
Over $\mathbb Q$, use $n=5$ and nodes $0,1,2,3,4$.
The polynomials
\[
 q_1(Y)=Y^2-4Y+2,\qquad
 q_2(Y)=Y^3-\tfrac{77}{5}Y+\tfrac{54}{5}
\]
satisfy both kernel conditions.  Expanding (I.34) gives
$\det B=\frac65 Q$, where
\[
\begin{split}
 Q={}&3X_1X_2-4X_1X_3-15X_1X_4+16X_1X_5
       +4X_2X_3+8X_2X_4-15X_2X_5\\
    &+4X_3X_4-4X_3X_5+3X_4X_5.
\end{split}
\tag{I.35}
\]
Thus $Q(t\mathbf1+\operatorname{SV}_1)=0$ and $Q\ne0$.
For a smaller representation than the general width-six bound,
take row vectors and column vectors
\[
\begin{array}{c|ccccc}
i&1&2&3&4&5\\ \hline
a_i&(1,0)&(0,1)&(-28/97,-4/97)&(3/16,0)&(0,0)\\
b_i&(0,0)^T&(3,0)^T&(-4,4)^T&(-15,8)^T&(16,-15)^T.
\end{array}
\]
In blocks of sizes $1,2,1$, put
\[
 N_i=\begin{pmatrix}0&a_i&0\\0&0_{2\times2}&b_i\\0&0&0\end{pmatrix}.
\]
All triple products of these $N_i$ vanish, and
$e_1^TN_iN_je_4=a_i b_j$.
Consequently
\[
 e_1^T\prod_{i=1}^5(I_4+X_iN_i)e_4
      =\sum_{i<j}(a_ib_j)X_iX_j=Q.
\]
The factors are upper triangular with diagonal entries one,
so every interior interval remains invertible at every point.
This refutes the one-pair proposal already at width four;
it does not assert that width four is the smallest possible
width of a centered one-pair annihilator of arbitrary degree.

\paragraph{A concrete repair for this family, and its scope.}
The same counterexamples are detected by a different simple
map.  Since every $U_S$ is nonzero, the support of $P$ is
exactly the set of bases of the column matroid of $V$.
Pairwise distinct element weights have a unique minimum-weight
matroid basis, by the greedy exchange property.  Hence
\[
 P(T,T^2,\ldots,T^n)\ne0,
\]
with degree at most $n(k+1)$.
This is an instance of the basis argument in I-03 and gives a
polynomial scalar sampling grid for this particular family:
choose $n(k+1)+1$ distinct scalar values, in an extension if
the ground field is too small.  This additional grid-size
condition is separate from the $n$ distinct Lagrange nodes.
It demonstrates that the dense-SV counterexample is a precise
test of that map, not a general obstruction to cocharacter
sampling.  The open dense-map question is whether a
logarithmic or polynomial number of pairs suffices for the
full rank-test class, or whether a different structured family
can handle its cancellations with polynomial dependence on
width.  Neither assertion is proved by this example.


\subsection{I-B08. The logarithmic dense-generator bound already holds in the every-order subclass}
\label{I-B08}

Forbes, Saptharishi, and Shpilka (2014), full version
arXiv:1309.5668, Definition~1.1 (printed p.~3), use the term
\emph{commutative ROABP} for a polynomial admitting a width
bound in \emph{every} variable order.  This is not the class
with one hidden order.  Pairwise commuting layer coefficient
matrices, as in I-26, give a representation in this subclass
by reordering the same factors.

Their Corollary~4.3 (printed p.~22), using Theorem~4.1
(printed pp.~21--22), states that if $f$ has individual degree
less than $d$, admits width-$w$ representations in every order,
and $|F|>nd$, then
\[
 f\ne0\quad\Longleftrightarrow\quad
 f\bigl(\operatorname{SV}_{1+2\lfloor2\log_2w\rfloor}\bigr)
       \ne0.
\]
The stated hitting-set size is
$\operatorname{poly}(n,d)^{O(\log w)}$.
One may use a sufficiently large extension when required.
This known logarithmic pair bound applies to the general
commuting, unipotent representation constructed in I-26.
Consequently its logarithmic necessary pair count matches
the order of a known sufficient bound for that restricted
class.  The new question retained after I-26 concerns the
full one-hidden-order rank-test class, and even a logarithmic
pair bound would still require a separate efficient sampling
argument to give polynomial rather than quasipolynomial size.

The primary definition, theorem, and corollary, including
the local factorization use of the every-order assumption
and the additive composition of two SV maps, were inspected.
The complete article and all of its subsidiary rank-condenser
proofs are not certified here.


\subsection{I-27. Filtering the actual Hoza masks does not restore a short coherent code}
\label{I-27}

\paragraph{The precise repair being tested.}
I-23 required every seed of a code to preserve the Grassmann cone.
One could instead keep only those seeds that happen to preserve it.
For the actual base-six weighting, this filtering has an exact
description and cannot combine noncoherent bits into a coherent
action.  This is a statement about that formula and formal
cocharacters; it is not a lower bound for arbitrary nonlinear
decoders or dense interpolation maps.

Let $2\le r\le w-2$, let
$C:\binom{[w]}r\longrightarrow\{0,1\}^{\ell}$ be any binary
code, and put
\[
 \rho_s(I)=\sum_{j=1}^{\ell}6^{\ell-j}s_j C(I)_j.
 \tag{I.36}
\]
Call a bit $j$ coherent if $I\mapsto C(I)_j$ satisfies the
Pl\"ucker exchange equalities in I-23.  Call a seed coherent
if the diagonal action $p_I\mapsto t^{\rho_s(I)}p_I$
preserves every decomposable $r$-vector.
Then
\[
 s\text{ is coherent}
 \quad\Longleftrightarrow\quad
 s_j=0\text{ for every noncoherent bit }j.
 \tag{I.37}
\]
Thus the set of coherent seeds is a coordinate face of the
Boolean seed cube, rather than a collection of cancellations
between noncoherent code positions.

For any selected family $\mathcal S$ of coherent seeds, let
$T\subseteq[w]$ consist of the ground elements tested by its
active nonconstant bits.  By the binary classification of I-23,
each such bit is membership of one element or its complement,
so $|T|\le\ell$.  The entire family of weights depends on an
$r$-subset $I$ only through $I\cap T$.  In particular, if
$\ell\le w-2$, there are distinct $u,v\notin T$ such that
\[
 \rho_s(K\cup\{u\})=\rho_s(K\cup\{v\})
 \quad\text{for all }s\in\mathcal S
 \tag{I.38}
\]
whenever $|K|=r-1$ and $K\cap\{u,v\}=\varnothing$.
Consequently a short actual code cannot be repaired merely
by discarding its noncoherent seeds and using the surviving
weights as ground-variable cocharacters.

\paragraph{Proof of the exact filtering statement.}
For each Pl\"ucker exchange equality, write its deficit for
bit $j$ as
\[
 d_j=C(Kab)_j+C(Kcd)_j-C(Kac)_j-C(Kbd)_j.
\]
Each $d_j$ is an integer in $[-2,2]$.  Coherence of the seed
requires $\sum_j 6^{\ell-j}s_jd_j=0$ for this exchange,
and likewise for the other pairing.  If there were a largest
active power $6^h$ with nonzero coefficient, its absolute
contribution would be at least $6^h$, whereas all lower
powers together contribute at most
\[
 2\sum_{a=0}^{h-1}6^a
   =\frac25(6^h-1)<6^h.
\]
This is impossible.  Hence $s_jd_j=0$ for every $j$ and
every exchange.  An active bit must therefore be coherent.
Conversely a sum of coherent weight functions satisfies all
exchange equalities and hence preserves decomposability by
I-23.  This proves (I.37).  The common positive layer
multiplier in Hoza's formula does not affect the argument.
The deficit argument is an equality of integer exponents,
so it is valid in every coefficient-field characteristic.

For the family statement, each active coherent nonconstant
bit is a membership or complementary membership function of
one element, by I-23.  All other active bits are constant.
It follows directly that $\rho_s(I)$ is determined by
$I\cap T$.  Two untested elements exist when $|T|\le w-2$,
and exchanging them proves (I.38).  A suitable $K$ exists
since $r-1\le w-2$.

\paragraph{An explicit nonzero cancellation witness.}
Write the induced coherent weights as
$\rho_s(I)=c_s+\sum_{i\in I}a_{s,i}$.
Equation (I.38) forces $a_{s,u}=a_{s,v}$ for every seed;
this equality is independent of the common gauge ambiguity
in the choice of the $a_{s,i}$.  The nonzero polynomial
\[
 f(X_1,\ldots,X_w)=X_u-X_v
\]
therefore vanishes under every induced monomial substitution
$X_i\leftarrow Z^{a_{s,i}}$.  Adding the same scalar center
to every coordinate does not change that conclusion.
The assertion also holds for Laurent exponents; adding a
common exponent to make all exponents nonnegative preserves
the equality of the two relevant coordinates.

This is already an invertible-factor width-two example:
\[
 f=e_1^T(I_2+X_uE_{12})(I_2-X_vE_{12})e_2,
\]
with identity factors at the other positions.  Equivalently,
the decomposable vector
$e_K\wedge(X_u e_u+X_v e_v)$ has two Pl\"ucker coordinates,
and their oriented difference gives this same polynomial.
Thus the witness belongs to the algebraic setting that the
coherent decoder was intended to preserve, rather than merely
showing a combinatorial collision between unused codewords.
\hfill$\square$

\paragraph{Magnitude and scope.}
To distinguish all $r$-subsets by such a coherent family,
at least $w-1$ ground elements must be tested.  Consequently
at least $w-1$ distinct nonconstant bit positions must occur
in its seeds.  Among their powers of six, one is at least
$6^{w-2}$.  A seed activating that bit, and a codeword on
which the bit is one, give a raw weight at least $6^{w-2}$.
This last observation concerns the unmodified weights;
rescaling an individual cocharacter can change its magnitude,
so it is not asserted as an invariant degree lower bound
under arbitrary further decodings.  The exact face
description (I.37) and the short-code collision (I.38) remain
valid even if common factors of individual seeds are removed.

Longer codes followed by a separately justified small family,
or maps that leave the Grassmann cone, are not excluded.
Nor does this obstruct adding a dense SV direction: for
distinct nodes, $\lambda_u(Y)-\lambda_v(Y)$ is nonzero and
already destroys this particular linear collision.  I-26
tests that next dense repair by a different, determinant
counterexample, and leaves logarithmically many pairs open
for the general one-hidden-order rank-test class.


\subsection{I-28. Nonzero endpoint constants fit inside the Grassmannian support budget}
\label{I-28}

\paragraph{Statement.}
The two extra endpoint positions in I-25 are unnecessary when
the endpoint vectors are nonzero at the center.  Let
$1\le r\le w-1$, put $g=r(w-r)$, and consider
\[
 q(Y_0,\ldots,Y_{m+1})
 =\left\langle u(Y_0)
       \prod_{i=1}^m\bigwedge^r A_i(Y_i),\ v(Y_{m+1})\right\rangle.
 \tag{I.39}
\]
Here $A_i(0)$ is invertible, $u(Y_0)$ is a polynomial
decomposable exterior row with $u(0)\ne0$, and $v(Y_{m+1})$
is a polynomial exterior column with $v(0)\ne0$.
Only the left row is required to be decomposable.
All variables displayed are distinct; constant endpoints may
be included by omitting their variables.
If $q\ne0$, setting all but at most $g$ variable positions
to zero leaves a nonzero polynomial.  In particular
$\operatorname{SV}_{g}$ hits this class, and the count includes
both endpoint variables.

\paragraph{Proof.}
Extend the coefficient field to an algebraic closure; a zero
restriction over the original field stays zero after this
extension, so proving existence there suffices.  Use constant
changes of basis to make all $A_i(0)=I_w$, as in I-25.
They preserve nonzero endpoint constants and decomposability.
Restrict each internal parameter to the nonempty open set
where $A_i$ is invertible, and the left parameter to the open
set where $u$ is nonzero.  These opens contain zero.

Let $Z_{-1}$ be the projective point $[u(0)]$, and let $Z_0$
be the closure of the image of $[u(Y_0)]$.  For $i\ge1$,
let $Z_i$ be the closure of the projective image after applying
the first $i$ internal matrices.  Each $Z_i$ is an irreducible
closed subvariety of $\operatorname{Gr}(r,w)$, and
\[
 Z_{-1}\subseteq Z_0\subseteq\cdots\subseteq Z_m.
\]
For the first inclusion, use the available parameter value
$Y_0=0$; for the others, use $A_i(0)=I_w$.

Retain a parameter position exactly when this inclusion is
strict.  There are at most $\dim Z_m$ retained positions,
because a strict inclusion of irreducible closed varieties
strictly increases dimension, starting at a point.
Setting all other positions to zero preserves the final
image closure $Z_m$.  To check this rigorously by induction,
at a retained position the previous restricted image is
dense in the same $Z_{i-1}$; applying the regular invertible
actions and taking closure gives the same $Z_i$.  At a
nonretained position, zero substitution applies the identity
and leaves $Z_{i-1}=Z_i$.  The initial left curve obeys the
same statement: if its image closure is a point, setting its
parameter to zero preserves that point.

Write $\ell_0$ for the nonzero linear functional determined
by $v(0)$.  If $\ell_0$ does not vanish identically on $Z_m$,
set the right parameter to zero as well.  The retained image
is still dense in $Z_m$, so the resulting scalar polynomial
is nonzero, using at most $\dim Z_m\le g$ positions.

Otherwise $Z_m$ lies in the hyperplane $\ell_0=0$.
The Grassmannian is not contained in that hyperplane:
the coordinate decomposable wedges span the full exterior
space, and $\ell_0\ne0$.  Therefore
\[
 \dim Z_m\le \dim\bigl(\operatorname{Gr}(r,w)
                                  \cap\{\ell_0=0\}\bigr)\le g-1.
\]
As $q\ne0$, some value of the right parameter gives a
functional that does not vanish identically on $Z_m$.
The retained prefix image has the same closure, so keeping
this right parameter free gives a nonzero scalar polynomial
on at most $\dim Z_m+1\le g$ positions.  This proves the
restriction assertion in both cases.  A nonzero restriction
to at most $g$ coordinate positions is a specialization of
$\operatorname{SV}_g$, by setting its selectors to the
corresponding Lagrange nodes.  Thus the generator assertion
follows in every characteristic, with the usual distinct-node
field or extension convention.  No algorithm to find the
retained positions from the unknown matrices is assumed.
\hfill$\square$

\paragraph{Application to the width-three construction.}
In the rank-two stage of I-21, a selected nonzero minor of
a local rank test is a contraction of a second-compound
word with an invertible interior.  Its compound row and
column at either endpoint are nonzero at the initial common
center $x_i=t$.  Indeed, a nonzero selected contraction has
nonzero endpoint rows and columns as univariate polynomial
vectors; evaluating a nonzero univariate entry at the
transcendental element $t$ cannot make it zero.  Constant
endpoints are covered directly.  Single-factor tests are
already preserved by the same argument and need no pairs.

Consequently I-28 applies with $w=3,r=2$, giving support
budget $g=2$ for the entire rank-two test, including its
endpoints.  Two fresh SV pairs simultaneously preserve all
rank-two tests: each selected minor is a nonzero polynomial
after substitution, and the product of the finitely many
selected minors is nonzero.  Previously preserved rank-three
minors remain nonzero upon setting the new amplitudes to zero.
The rank-one stage of I-21 still uses its independently
proved six-pair budget.  With disjoint fresh parameter pairs
the two stages therefore give
\[
 \mathcal G_i^{(8)}
      =t+\sum_{j=1}^{8}z_j\lambda_i(y_j),
 \qquad f\ne0\Longrightarrow f(\mathcal G^{(8)})\ne0
 \tag{I.40}
\]
for every width-at-most-three one-hidden-order ROABP in the
scope of I-21.

All degree and interpolation arguments of I-21 remain
unchanged except for replacing ten pairs by eight.  In its
notation, this yields the explicit sample bound
\[
 (D+1)\bigl((n-1)D+1\bigr)^8\binom{D+8}{8},
 \tag{I.41}
\]
over a working field containing at least
$\max\{n,D+1,(n-1)D+1\}$ distinct available elements.
The earlier ten-pair theorem remains valid; (I.40) is a
strictly smaller proved sampler.  This refinement does not
address the singular reset profiles at variable width and
does not claim a polynomial dependence on that width.


\subsection{I-29. A polynomial pair budget for arbitrarily many resets at one singular rank}
\label{I-29}

\paragraph{Statement and scope.}
The rank-profile mechanism used in I-21 has a growing-width
positive specialization.  Fix $1\le r<w$, put $c=w-r$, and
suppose a constant-boundary width-$w$ ROABP representation
of $f$ in one hidden variable order has the following promise:
every nonempty interval of its internal matrix word has generic
rank either $r$ or $w$.  The boundary row and column are not
included in this promise.  There may be arbitrarily many
rank-$r$ factors and intervals; their matrices need not commute.
Then the generator
\[
 X_i=t+\sum_{j=1}^{K}z_j\lambda_i(y_j),
 \qquad K=2r(w-r)+w-1,
 \tag{I.42}
\]
hits every nonzero such $f$ in every characteristic, over a
field or extension supplying the requisite distinct nodes.
For total degree at most $D$, the explicit grid construction
of I-21 has size at most
\[
 (D+1)\bigl((n-1)D+1\bigr)^K\binom{D+K}{K}.
 \tag{I.43}
\]
Its field-size convention is the same as I-21.

This is a polynomial \emph{pair budget} in $w$, with
sampling cost $(nD)^{O(w^2)}$, not polynomial-size PIT for
variable width.  The statement permits unbounded total
corank and so is different from I-20.  It further specializes
the cited unit-rank mechanism; no historical novelty is claimed.

\paragraph{1. A two-rank coefficient-space lemma.}
Let $B_i(Y_i)$ be $w\times w$ polynomial matrices in distinct
variables over a field $L$.  Assume every interval has generic
rank equal to its rank at zero, and that these ranks belong
to $\{r,w\}$.  Then, for every constant row $b$,
\[
 \begin{gathered}
 \operatorname{Coeff}(bB_1\cdots B_m)
 =\operatorname{span}_L\{
       [Y^\alpha](bB_1\cdots B_m):
       |\operatorname{supp}\alpha|\le h\},\\
 h=r(w-r)+w-1=rc+w-1.
 \end{gathered}
 \tag{I.44}
\]
The assertion remains valid in a truncated product local ring
when the corresponding interval unit-rank assumptions hold.

If the full constant product is invertible, all constant
factors are invertible, and I-17's nested row-space argument
gives the smaller bound $w-1$.  Otherwise its rank is $r$.
Choose a persistent $r$-dimensional row space $U_0$ on which
the full constant product is injective and transport it to
$U_i$ at every cut.  Choose a complement $V_m$ of $U_m$;
the suffix pullback
\[
 V_i=\{z:zB_{i+1}(0)\cdots B_m(0)\in V_m\}
\]
has dimension $c$, complements $U_i$, and is carried into
$V_{i+1}$ by the next constant matrix.  In compatible bases
the constants therefore have the form
\[
 B_i(0)=\operatorname{diag}(I_r,Q_i),
 \tag{I.45}
\]
where $Q_i$ is either invertible or zero: the rank of the
whole factor is either $w$ or $r$.  These basis changes
preserve supports and interval ranks.

Let $\pi:L^{1\times w}\to L^{1\times r}$ be active-coordinate
projection.  Prefix coefficient spaces $S_i$ have nested
projections $\pi S_i$, by (I.45).
We need the following stable-segment assertion: on a segment
where every projected space is the same $T$, give basis lifts
of $T$ initial cost zero and a basis of the initial projection
kernel cost one.  Cost means initial cost plus the number of
nonconstant local coefficients used.  Then cost at most $c$
spans the final coefficient space.

First suppose all $Q_i$ in the segment are invertible.
After constant transport, the coefficient spaces are nested
inside $\pi^{-1}(T)$, whose dimension is $\dim T+c$.
If the initial kernel has dimension $k>0$, all initial rows
cost at most one and there can be at most $c-k$ subsequent
dimension increases.  At each increase one nonconstant
coefficient suffices, so the cost is at most
$1+c-k\le c$.  If $k=0$, the initial cost is zero and
there are at most $c$ increases.  Stable layers require only
constant propagation.  This proves the assertion in this case.

If a zero $Q_i$ occurs, take its last position $a$ and put
$P=B_a\cdots B_b$, where $b$ is the segment's final position.
The suffix has rank $r$ and constant matrix
$\operatorname{diag}(I_r,0)$: its rank is at most $r$ from
the factor at $a$, and its constant $r$-minor is nonzero.
In the power-series ring at zero, or in the truncated local
ring, write its top blocks as $[U\ V]$.  The matrix $U$
is invertible.  The vanishing $(r+1)$-minors give the row-graph
identity
\[
 (x,y)\in\operatorname{rowmod}(P)
       \quad\Longrightarrow\quad y=xU^{-1}V.
\]
Let $B$ be the matrix of the initial lifts propagated through
the constants before position $a$, and let $H=\pi B$ be the
chosen basis of $T$.  Stable projection implies
\[
 \pi(BP)=AH,\qquad \pi(sP)=\eta H\quad(s\in S_{a-1}),
 \qquad A(0)=I.
\]
Thus $A$ is invertible and the row graph gives
$sP=\eta A^{-1}BP$.  If $T=0$, it instead gives $sP=0$.
It follows, by taking coefficient rows, that the final space
is exactly $\operatorname{Coeff}(BP)$.  Indeed, the displayed
identity proves one containment, and the rows of $B$ already
belong to $S_{a-1}$ for the other.

Restart after position $a$ from
$\operatorname{Coeff}(BB_a)$.  It has constant lifts of $T$
of cost zero and a kernel basis of cost at most one, since
only that one variable position was used.  The remaining
constants are invertible.  The preceding dimension argument
then bounds the total cost by $c$, proving the stable-segment
assertion.  In the untruncated case the inverses may be formal
power series; each coefficient of their product with $BP$
is still a finite linear combination of coefficient rows of
$BP$.  No analytic convergence is used.

We now count support from the original single row $b$.
If $\pi b\ne0$, there are at most $r-1$ strict projection
growths.  The initial stable segment costs at most $c$;
each growth factor and the following stable segment add at
most $1+c$.  The bound is
\[
 c+(r-1)(c+1)=rc+r-1\le h.
\]
If $\pi b=0$, the first stable segment lies in the passive
$c$-space.  If it contains a zero $Q_i$, the row-graph
argument kills the entire current coefficient space and the
final word is zero.  Otherwise its constants are invertible,
and starting from one row there are at most $c-1$ dimension
increases.  This first segment costs at most $c-1$.
At most $r$ projection growths follow, giving the bound
\[
 (c-1)+r(c+1)=rc+r+c-1=rc+w-1=h.
\]
At every stage the retained rows are linear combinations of
actual coefficient rows of the original prefix with the
stated support.  At a restart the chosen lifts are propagated
by constants through skipped positions; selecting a local
nonconstant coefficient adds only its own variable position.
Thus the cost count proves the actual support assertion
(I.44), rather than just a bound on an abstract generating
module.  The same reasoning applies in the truncated ring.

\paragraph{2. Sharpness for every active and passive dimension.}
Write the active coordinates as $a_1,\ldots,a_r$ and the
passive coordinates as $p_1,\ldots,p_c$, and put
$P=\operatorname{diag}(I_r,0_c)$.
Use a fresh variable in every factor below.
First apply the $c-1$ unipotent chain factors
\[
 I+Y_1E_{p_1,p_2},\ I+Y_2E_{p_2,p_3},\ \ldots,
 I+Y_{c-1}E_{p_{c-1},p_c}.
\]
This initial list is empty if $c=1$.
For each $j=1,\ldots,r$, apply one reset
\[
 P+S_jE_{p_c,a_j},
\]
followed by $c$ unipotent chain factors for the edges
$a_j\to p_1\to p_2\to\cdots\to p_c$.
The total number of factors is
$(c-1)+r(c+1)=h$.

Each reset has generic rank exactly $r$: its only extra
row is a multiple of the existing active row at $a_j$.
Each chain factor is unipotent and invertible.  Every
interval without a reset has constant matrix $I_w$ and
generic rank $w$.  Every interval with a reset has constant
matrix $P$, has generic rank at most $r$ from the reset
factor, and has rank at least $r$ from its constant active
minor.  Thus every interval satisfies precisely the
two-rank promise, including rank preservation at zero.

Starting with $e_{p_1}^T$, the coefficient in the last
passive coordinate after the initial chain is the product
of all its $c-1$ variables.  At reset $j$, all passive
coordinates are killed except that last coefficient,
which is transferred, multiplied by $S_j$, into the
fresh active coordinate $a_j$.  Earlier active coordinates
are unchanged.  The following chain multiplies that
coefficient by each of its $c$ fresh variables before
placing it in $p_c$.  The earlier active coordinates
do not enter this chain.  Induction on $j$ proves that
the final contraction with $e_{p_c}$ is exactly the
product of all $h$ variables.

Consequently every coefficient row of support less than
$h$ is annihilated by $e_{p_c}$, whereas the full-support
coefficient is not.  This proves sharpness of (I.44) for
every $1\le r<w$.  In particular the corank-one case has
$h=2(w-1)$, and balanced active/passive dimensions require
quadratic support.  The scalar witness is a monomial and
is hit by one generic dense direction; this is not a
lower bound on dense-generator pair count.

\paragraph{3. Uniform preservation and the sampler.}
Start at the common transcendental center $X_i=t$.
All rank-$w$ intervals retain rank: each of their univariate
factors has a nonzero determinant, and the product of those
determinants remains nonzero at $t$.

For every rank-$r$ interval, use I-21's finite local-rank-test
construction.  Minimality makes each test's interior chunk
have rank greater than $r$, and the two-rank promise makes
that rank $w$.  Its selected nonzero $r$-minor is therefore
a contraction of an invertible compound chunk with at most
two endpoint factors.  Its endpoint compound row is
decomposable, and both endpoint vectors are nonzero at the
common transcendental center, exactly as checked in I-28.
That result gives a support budget $g=r(w-r)=rc$ for the
whole test, including the endpoints.

Add $g$ fresh SV pairs.  Each selected nonzero minor stays
nonzero, and their finite product is nonzero, so all such
tests and hence all rank-$r$ intervals are simultaneously
preserved over the parameter field.  Rank-$w$ intervals
remain preserved because the new amplitudes may be set to
zero.  This is an existence argument over one common
parameter field and requires neither the hidden order nor
the selected minors as input to the generator.

At the resulting center, translate all original variables
by independent residual variables.  Every interval now
has its generic rank at zero, in $\{r,w\}$.  The original
scalar remains a nonzero polynomial, since translation
over that field is invertible.  Apply (I.44) to its starting
row and then its final column.  Some scalar monomial has
support at most $h$; otherwise the final column would
annihilate the spanning coefficient rows and the whole
polynomial would be zero.  Adding $h$ further fresh SV pairs
therefore makes the scalar nonzero.  The two additive maps
combine into $K=g+h=2rc+w-1$ pairs, proving (I.42).

For sampling, the degree in $t$ is at most $D$, each selector
degree is at most $(n-1)D$, and the total amplitude degree
is at most $D$.  The same product grid and amplitude Newton
simplex as in I-21 give (I.43).  All arguments and the
interpolation convention are characteristic independent.
\hfill$\square$

\paragraph{What remains outside this positive case.}
I-24's exponentially long support witnesses use several
different singular ranks, rather than only one fixed rank
and full rank.  In that setting the passive blocks in
(I.45) may themselves have intermediate nonzero singular
ranks.  The last-zero restart used above then does not
replace the whole passive block, and recursive rank losses
can accumulate.  I-29 identifies a precisely proved regime
where this recursion is unnecessary.  It neither asserts
that recursion is necessary for dense nonvanishing nor
excludes rational-function or cocharacter constructions
for the remaining profiles.


\subsection{I-30. Absorbing endpoint vectors into unit factors; a six-pair width-three sampler}
\label{I-30}

\paragraph{An elementary endpoint factorization.}
Let $u(Y)$ be a polynomial row of length $W$ over a field $L$,
with $u_0=u(0)\ne0$.  Choose $j$ with $(u_0)_j\ne0$ and set
\[
 B(Y)=I_W+\frac{e_j}{(u_0)_j}\bigl(u(Y)-u_0\bigr).
 \tag{I.46}
\]
Then
\[
 u(Y)=u_0B(Y),\qquad B(0)=I_W,\qquad
 \det B(Y)=\frac{u_j(Y)}{(u_0)_j}\ne0.
\]
The first two identities follow by multiplication; the third
is the rank-one determinant identity
$\det(I+ab^T)=1+b^Ta$.
For a polynomial column $v(Z)$ with $v_0=v(0)\ne0$, choosing
$j$ with $(v_0)_j\ne0$ likewise gives
\[
 \begin{gathered}
 C(Z)=I_W+\bigl(v(Z)-v_0\bigr)\frac{e_j^T}{(v_0)_j},
 \qquad v(Z)=C(Z)v_0,\\
 C(0)=I_W,
 \qquad \det C(Z)=\frac{v_j(Z)}{(v_0)_j}\ne0.
 \end{gathered}
 \tag{I.47}
\]
Thus both new endpoint factors are generically invertible and
invertible in every truncated local ring at zero.  No
characteristic restriction is used.

Suppose the interior word has generic interval ranks equal
to its constant interval ranks, in some set of allowed ranks
containing $W$.  Then the enlarged word
$B(Y)A_1\cdots A_m C(Z)$ has the same property.  An interval
lying in the old word is unchanged.  Every other interval is
an old prefix, suffix, or whole word, multiplied on one or
both sides by generically invertible matrices whose constants
are identities; both generic and constant ranks are preserved.
An interval consisting only of a new endpoint factor has
rank $W$.  The same argument holds for interval unit ranks
in a truncated local ring, because multiplication by an
invertible matrix preserves its determinantal rank ideals.

Consequently every coefficient-span support bound proved for
constant endpoints and this rank profile also applies to
\[
 u(Y)A_1\cdots A_m v(Z)
      =u_0B(Y)A_1\cdots A_m C(Z)v_0,
\]
with both endpoint positions included in the original bound.
This costs no additional support positions and is an
existence argument about the representation, not a request
to recover unknown endpoint coefficients from the black box.

\paragraph{Application to all width-three rank tests.}
Use the rank-descent construction of I-21, starting at
$X_i=t$.  Full-rank intervals are preserved as before.
In a selected rank-two test, its second-compound interior
has width three and invertible constant factors.  Its
nonzero polynomial endpoint row and column have nonzero
constants at this center.  Absorb them by (I.46)--(I.47).
The enlarged compound word has all constant factors
invertible, so I-17's row-space argument gives support at
most $3-1=2$ for a nonzero scalar coefficient.  Equivalently
one may use I-28.  Two fresh SV pairs preserve all rank-two
tests simultaneously, since the finite product of their
nonzero substituted minors is nonzero.

The new center is
\[
 \alpha_i=t+\sum_{j=1}^2z_j\lambda_i(y_j).
\]
Every $\alpha_i$ remains transcendental over the original
coefficient field: it has coefficient one on the fresh
indeterminate $t$.  Therefore every nonzero univariate
polynomial endpoint row or column from an original factor
is nonzero at its center.  Constant endpoint vectors are
covered directly.  A selected nonzero contraction cannot
have an identically zero endpoint vector.

In a rank-one test, the interior interval ranks are two or
three and have already been preserved at the current center.
Translate to independent residual variables over the field
of the current parameters.  Absorb both endpoint vectors
using (I.46)--(I.47); every interval of the enlarged word
still has generic and constant rank in $\{2,3\}$.
I-29 with $(w,r)=(3,2)$ now gives support
\[
 h=2(3-2)+3-1=4
\]
for the entire test, including its endpoints.  Four further
fresh SV pairs preserve all the rank-one tests simultaneously.
The preceding rank-two and rank-three minors stay nonzero
because the new amplitudes may be set to zero.
I-21's local-rank induction therefore preserves the rank-one
augmented full word, whose single possibly nonzero entry is
the original scalar $f$.

It follows that the thirteen-parameter map
\[
 \mathcal G_i^{(6)}(t,\mathbf y,\mathbf z)
      =t+\sum_{j=1}^{6}z_j\lambda_i(y_j)
 \tag{I.48}
\]
is a hitting generator for every width-at-most-three ROABP
in one hidden order, in every characteristic.  This is the
same class and evaluation model as I-21; it assumes no
two-rank promise on the whole original word.  The two-rank
lemma is used only inside the rank-one local tests, where
the promise follows from the rank-descent construction.

\paragraph{Exact evaluation bound.}
For total degree at most $D$, the degree in $t$ is at most
$D$, each selector degree is at most $(n-1)D$, and the total
amplitude degree is at most $D$.  Applying the grid and
Newton-simplex argument of I-21 gives the hitting-set bound
\[
 (D+1)\bigl((n-1)D+1\bigr)^6\binom{D+6}{6}.
 \tag{I.49}
\]
The working field must supply at least
$\max\{n,D+1,(n-1)D+1\}$ distinct elements, and evaluations
over that working extension are required if it differs
from the coefficient field.  Earlier ten- and eight-pair
theorems remain valid; (I.48) is the strongest sampler
proved in this track.  Nothing here proves polynomial
dependence on arbitrary width.
\hfill$\square$


\subsection{I-31. Full matrix cones give a pair budget controlled by the number of rank levels}
\label{I-31}

\paragraph{Attribution and statement.}
The full-matrix cone method in Liu (2026), Section~3.3,
printed pp.~7--8, avoids treating compound matrices as
arbitrary vectors.  We prove an original-coordinate version
with a somewhat looser block count, so no first-block
function-field factorization is needed.  This is a further
checked specialization of that explicitly unrefereed source,
not a claim of historical novelty or a certification of its
affine-orbit theorem.

Fix an allowed rank set
$\mathcal R=\{s_1<\cdots<s_L=w\}\subseteq\{1,\ldots,w\}$,
where $w\ge2$.  Put
\[
 d_w=w^2-1,\qquad
 d_s=(2ws-1)\left(2\max_{t\in\mathcal R,\ t>s}d_t+1\right)
       \quad(s\in\mathcal R,\ s<w),
 \qquad K_{\mathcal R}=\sum_{s\in\mathcal R}d_s.
 \tag{I.50}
\]
For constant-boundary width-$w$ ROABPs in one hidden order
whose nonempty internal intervals have generic ranks in
$\mathcal R$, the map $t\mathbf1+\operatorname{SV}_{K_{\mathcal R}}$
is a hitting generator in every characteristic.  It satisfies
\[
 K_{\mathcal R}\le2w^2(4w^2)^{L-1}.
 \tag{I.51}
\]
Thus a fixed number of allowed rank levels gives a polynomial
pair budget in $w$, even with arbitrarily many noncommuting
singular factors.  The sample count is still
\[
 (D+1)\bigl((n-1)D+1\bigr)^{K_{\mathcal R}}
                  \binom{D+K_{\mathcal R}}{K_{\mathcal R}},
 \tag{I.52}
\]
for degree bound $D$, with the working-field convention of
I-21.  For the full rank set $\{1,\ldots,w\}$ this is only
$K_{\mathcal R}=2^{O(w\log w)}$, so it proves polynomial size
for each fixed width, not polynomial dependence on variable
width.  I-29 and I-30 give substantially smaller bounds in
their respective special cases.

\paragraph{1. Cones and composition.}
Work over an algebraic closure of the coefficient field.
For a polynomial matrix $M$, write $\mathcal C(M)$ for the
Zariski closure of all scalar multiples of its values.
This is an irreducible affine cone.  Equality of such cones
preserves every homogeneous polynomial identity in the matrix
entries.  It also composes under multiplication on disjoint
variable sets: replacing a factor by a restriction with the
same cone leaves the cone of the whole product unchanged.
Indeed, taking a dense set in either input of the regular
matrix-multiplication map does not change the closure of its
image, and scalar multiples are absorbed by the cone.
These observations remain true after an extension of the
coefficient field.

\paragraph{2. Deleting stable rank-$s$ blocks.}
Let $B_1,\ldots,B_m$ be polynomial $w\times w$ matrices in
disjoint variable blocks, each of generic rank $s$ and with
$B_i(0)=P_s=\operatorname{diag}(I_s,0)$.
For a constant $a\times w$ matrix $R$ with $\pi R\ne0$,
where $\pi$ selects the first $s$ columns, put
\[
 M_i=RB_1\cdots B_i,\qquad
 V_i=\mathcal C(\pi M_i),\qquad V_0=\mathcal C(\pi R).
\]
Then $V_{i-1}\subseteq V_i$, since the zero value of $B_i$
preserves the first $s$ columns.  The cones are irreducible,
start at a line, and lie in the $as$-dimensional matrix space.
There are at most $as-1$ strict dimension increases.

Suppose $V_{j-1}=V_j=V_{j+1}=V$.  We claim that setting
block $j$ to zero preserves the full cone of $M_{j+1}$.
Fix a value of $B_{j+1}$ on the dense open where its top-left
$s\times s$ block $E$ is invertible, and call its top-right
block $F$.  Rank $s$ gives the row-graph identity: any matrix
obtained by multiplying $B_{j+1}$ on the left is determined
by its first $s$ columns $Q$, and equals $Q[I_s,E^{-1}F]$.
After the restriction $B_j=P_s$, those first columns have
cone $VE$.  This is contained in $V$, since the restricted
values are among the original values contributing to
$V_{j+1}=V$.  But $VE$ is closed and irreducible, has
dimension $\dim V$, and is contained in the irreducible $V$;
therefore $VE=V$.

Every original output at this fixed value of $B_{j+1}$ has
its first columns in $V$, and hence belongs to the restricted
full cone by the row-graph identity.  Taking the closure over
the dense set of such values of $B_{j+1}$ proves the claim.
The opposite containment is automatic from restriction.
It also preserves all subsequent cones by composition.

Retain every strict-growth block and the last block of each
run of nongrowth blocks.  All other blocks can be deleted
one at a time using the claim, while the projected cones
and the full cone at the end of each affected pair remain
unchanged.  The number retained is at most
$2(as-1)+1=2as-1$.  In particular, with $R=I_w$, at most
$2ws-1$ blocks preserve the full product cone.

\paragraph{3. Coordinate compression of a preserved-rank word.}
We prove that if every interval of a univariate matrix word
has its generic rank at zero, all these ranks belong to
$\mathcal R$, and the full word has rank $s$, then setting
all but at most $d_s$ positions to zero preserves its full
matrix cone.

For $s=w$, all constant factors are invertible.  Gauge them
to identities.  The prefix cones are nested, starting with
the line through $I_w$ in dimension $w^2$.  Retain only the
strict-growth positions.  At a nongrowth position, replacing
that factor by its identity constant preserves the prefix
cone and then all later cones by composition.  At most
$w^2-1=d_w$ positions remain.

For $s<w$, greedily split the word at successive shortest
prefixes of rank $s$.  Write such a block as $C_jA_{e_j}$;
the preceding chunk $C_j$ is empty or has rank greater
than $s$, as do all of its subintervals.  When the remaining
tail no longer has rank $s$, absorb it into the last block;
that tail is empty or has rank greater than $s$ as well.
Every resulting block has generic and constant rank $s$.
Every product of consecutive blocks has constant rank $s$,
because it is an original interval and cannot have rank
below that of the full word.

Compatible constant bases at the block cuts normalize all
block constants to $P_s$.  To see this, transport an active
$s$-space injectively along the constant products.  At an
internal cut the left kernel of the next block complements
this active space, since the adjacent product retains rank
$s$.  Choosing those passive kernel bases makes each block
constant $\operatorname{diag}(I_s,0)$.
These changes multiply the entire word by fixed invertible
matrices at its ends and do not affect coordinate supports
or cone-preserving restrictions.

Part~2 retains at most $2ws-1$ blocks.  In each retained
block, the induction hypothesis compresses the higher-rank
prefix chunk $C_j$ to at most
$D_s=\max_{t>s,\ t\in\mathcal R}d_t$ positions; the final
tail, if present, costs at most another $D_s$.  Retain the
one factor $A_{e_j}$ between them.  By disjoint-block
composition these restrictions preserve the block cone and
then the full word cone.  Thus at most
$(2ws-1)(2D_s+1)=d_s$ positions suffice, completing the
induction.

\paragraph{4. Absorbing rectangular endpoints in the ground space.}
For an $r\times w$ polynomial matrix $L(Y)$ whose constant
$L_0$ has full row rank $r$, choose a constant right inverse
$R_0$ with $L_0R_0=I_r$.  Then
\[
 B(Y)=I_w+R_0(L(Y)-L_0),\qquad
 L(Y)=L_0B(Y),\qquad B(0)=I_w.
 \tag{I.53}
\]
The determinant of $B$ has constant term one, so $B$ is
generically invertible.  Similarly a $w\times r$ endpoint
$V(Z)$ with full-column-rank constant $V_0$ can be written
$V(Z)=C(Z)V_0$, where $C(0)=I_w$ and $C$ is generically
invertible, by choosing a constant left inverse of $V_0$.
Thus these endpoint factors can be absorbed in the original
$w$-dimensional matrix word, without passing to a compound
space.  Generic and constant interval ranks are unchanged,
just as in I-30.

\paragraph{5. Rank descent and final scalar detection.}
Start with the common center $X_i=t$, which preserves every
full-rank interval.  Process the remaining ranks in
$\mathcal R$ in decreasing order.  For a rank-$r$ local test
as in I-21, the interior chunk has rank greater than $r$,
and every one of its intervals is already preserved.
A selected nonzero $r$-minor of the test has the form
\[
 \det\bigl(L(Y_0)M(\mathbf Y)V(Y_{m+1})\bigr).
\]
The selected endpoint matrices have full row and column
rank $r$ at the current center.  Indeed, a nonzero selected
minor implies that each endpoint has some nonzero
univariate $r$-minor, and each current coordinate remains
individually transcendental over the original coefficient
field: it is $t$ plus terms independent of $t$.

Absorb these endpoints by (I.53).  The enlarged interior
word still has precisely the preserved higher-rank profile.
If its full rank is $s>r$, Part~3 gives a restriction to at
most $d_s\le D_r$ positions preserving its full matrix cone.
The fixed-endpoint determinant is homogeneous of degree
$r$ in the full matrix entries, so this restriction keeps
the selected minor nonzero.  It therefore has a monomial
of support at most $D_r$ and is hit by that many fresh SV
pairs.  All chosen test minors remain simultaneously nonzero
because their product is a nonzero polynomial.  Earlier
ranks remain preserved by setting the new amplitudes to zero.
The local-rank induction proves preservation of all internal
intervals at this rank.

After the smallest allowed rank is processed, every internal
interval is preserved.  Part~3 applied to the full word gives
a cone-preserving restriction on at most $d_s$ positions,
where $s$ is its generic rank.  The nonzero scalar output
is a homogeneous linear functional on the matrix entries,
so it remains nonzero on that restriction.  At most
$\max_{s\in\mathcal R}d_s=d_{s_1}$ additional SV pairs hit it.

The numbers $d_s$ strictly increase as the allowed rank
decreases.  At rank $s_i$ the preservation cost is therefore
$d_{s_{i+1}}$.  Summing those costs and the last scalar
detection cost gives exactly $K_{\mathcal R}$ in (I.50).
All pairs are fresh and their additive maps combine into
one SV map.  The map depends only on the bounds and the
chosen nodes, not on the unknown order or matrices.

Finally $d_s+1\le4w^2(D_s+1)$ and $d_w+1=w^2$.
The recurrence also gives $d_s>2D_s$ for $s<w$.
Hence the sum is at most twice its largest term, yielding
(I.51).  The degree and Newton-simplex sampling argument is
identical to I-21, proving (I.52).
\hfill$\square$

\paragraph{Remaining main gate.}
The proof replaces the exponential compound dimension by
matrix-cone dimensions, but a new block-selection factor
appears at each distinct singular rank level.  For bounded
$L$ this is polynomial in $w$ at the pair-count level;
for unrestricted $L$ the bound is still exponential in
$w\log w$.  Moreover even a polynomial pair budget would
still need an efficient sampling argument to yield a
polynomial-size hitting set.  No implication from Hoza's
path isolation to either of these missing steps has been
proved.  The dense lower bound I-26, the filtered-mask
obstruction I-27, and the known logarithmic every-order
result I-B08 continue to delimit the candidate repairs.


\subsection{I-32. Exponential coordinate-cone compression is necessary even at the common center}
\label{I-32}

\paragraph{Statement.}
For every $w\geq3$ and every coefficient field, put
\[
 h=3\cdot2^{w-2}-2,\qquad n=2h.
\]
There is an explicit width-$w$ word of $n$ univariate linear
matrices in distinct variables $x_1,y_1,\ldots,x_h,y_h$ with
the following properties. Image cones use the algebraic-closure
convention of I-31.

Every nonempty internal interval has generic rank in
$\{2,\ldots,w\}$, and that rank is preserved at every common
center $x_i=y_i=\tau$. In particular it is preserved at the
transcendental center $t\mathbf1$ used in I-31. Nevertheless,
after translation to residual variables at that center, a
restriction obtained by setting unselected coordinates to zero
preserves the full matrix image cone if and only if it retains
at least one coordinate from each pair $(x_i,y_i)$. The minimum
number of retained positions is therefore exactly $h$.

Thus a uniform polynomial-in-$w$ coordinate-cone compression
bound is false, even at a common transcendental center where
all internal interval ranks are preserved. In terms of the
number $L=w-1$ of allowed rank levels, the example needs
$3\cdot2^{L-1}-2$ positions. It precludes replacing the rank-level
recursion by a polynomial bound for this exact compression
criterion. It does not preclude changing the center or using
other dense maps: the distinguished nonzero scalar below is
detected by one centered SV pair.

\paragraph{1. The seed word.}
We recall the explicit matrices of I-24, including the property
needed here. Set $P_r=\operatorname{diag}(I_r,0_{w-r})$ and
\[
 Q_{r,a}(s,u)=(P_r+sE_{w,a})(P_r+uE_{a,w})
       \quad(2\leq r<w,\ a\leq r).
\]
With fresh variables at every occurrence, define
\[
 H_w=\text{empty word},\qquad
 H_k=H_{k+1}Q_{k,k}H_{k+1}\quad(3\leq k<w),
\]
and
\[
 W=H_3Q_{2,1}H_3Q_{2,2}H_3.
\]
If $L_k$ is the length of $H_k$, then
$L_w=0$, $L_k=2L_{k+1}+2$, and
$3L_3+4=h$. Label the successive variables of this word by
$Y_1,\ldots,Y_h$. Each factor is of one of the two forms
\[
 P_r+Y_iE_{w,a},\qquad P_r+Y_iE_{a,w}.
 \tag{I.54}
\]
They have generic rank $r$ and constant matrix $P_r$.

For completeness, if a row $q$ has $q_a=0$, then
\[
 qQ_{r,a}(s,u)
   =\sum_{j\leq r}q_j e_j^T
       +q_ws e_a^T+q_wsu e_w^T.
\]
Starting with coordinates $k,\ldots,w-1$ zero, $H_k$ leaves
the earlier coordinates unchanged and multiplies the last
coordinate by all its variables. This follows by induction:
the first $H_{k+1}$ leaves coordinate $k$ zero, $Q_{k,k}$
performs the displayed transfer and resets all higher active
coordinates, and the second $H_{k+1}$ applies again. In the
full word starting at $e_w^T$, the two rank-two gadgets use
successively the still-zero coordinates one and two and reset
the coordinates required by the intervening $H_3$ words.
Consequently
\[
 e_w^T W(Y_1,\ldots,Y_h)e_w=\prod_{i=1}^hY_i.
 \tag{I.55}
\]
This calculation is valid in every characteristic.

\paragraph{2. A paired lift with preserved common-center ranks.}
Replace every factor in (I.54) by two factors on fresh variables.
For $E=E_{w,a}$, use
\[
 (I_w+x_iE)\bigl((I_w-y_iE)P_r\bigr),
 \tag{I.56}
\]
and for $E=E_{a,w}$, use
\[
 \bigl(P_r(I_w+x_iE)\bigr)(I_w-y_iE).
 \tag{I.57}
\]
Since $E^2=0$, all matrices $I_w+zE$ are invertible with
inverse $I_w-zE$. In (I.56), $EP_r=E$; in (I.57),
$P_rE=E$. Each pair therefore equals
$P_r+(x_i-y_i)E$. The full lifted word $\widetilde W$ obeys
the matrix identity
\[
 \widetilde W(x_1,y_1,\ldots,x_h,y_h)
       =W(x_1-y_1,\ldots,x_h-y_h).
 \tag{I.58}
\]
Every factor is still univariate and linear. Exactly one half
of each pair is invertible; the other is an invertible matrix
times $P_r$, on the appropriate side, and has rank exactly $r$
for every parameter value.

We verify every partial interval, not only complete pairs.
At $x_i=y_i=\tau$, each complete pair is $P_r$.
A partial initial pair can only contribute its second half.
For (I.56) that half is $(I_w-\tau E)P_r$; for (I.57) it is
$I_w-\tau E$. A partial final pair can only contribute its
first half. For (I.56) it is $I_w+\tau E$; for (I.57) it is
$P_r(I_w+\tau E)$. All intervening pairs are complete.
Thus the constant interval product has the form
\[
 U_{\rm left}\left(\prod P_{r_i}\right)U_{\rm right},
 \tag{I.59}
\]
where the two boundary matrices are invertible, and the
projections are exactly those attached to singular halves
lying in the interval. The same description handles an
interval contained in one pair: a single half has its
displayed rank, and the complete pair is $P_r$. If no singular
half occurs, the interval is a product of invertible matrices.

The nested diagonal projections in (I.59) have product
$P_q$, where $q$ is the minimum rank of the singular halves
present; take $q=w$ if there are none. The common-center rank
is therefore exactly $q$. The generic interval rank is at
most $q$, because every singular half has that fixed rank
and products cannot increase it. It is at least $q$ because
the specialization (I.59) has rank $q$, so some $q$-minor is
a nonzero polynomial before specialization. These two bounds
prove equality. This argument applies to every $\tau$, in
particular to $\tau=t$ over the coefficient field extended
by a transcendental element.

\paragraph{3. The exact coordinate-compression threshold.}
Write $x_i=t+X_i$, $y_i=t+Y_i$. Formula (I.58) gives
\[
 \widetilde W(t+X_1,t+Y_1,\ldots,t+X_h,t+Y_h)
     =W(X_1-Y_1,\ldots,X_h-Y_h),
\]
and (I.55) becomes
\[
 e_w^T\widetilde W(t+\mathbf X,t+\mathbf Y)e_w
       =\prod_{i=1}^h(X_i-Y_i)\ne0.
 \tag{I.60}
\]
Let a coordinate restriction retain a set of residual positions
and set all other residual coordinates to zero. If it omits
both positions of any pair, the displayed scalar is identically
zero after restriction. The restricted matrix image cone is
then contained in the homogeneous hyperplane
$\{M:e_w^TMe_w=0\}$. The original cone is not contained in
that hyperplane by (I.60). Hence the two cones differ.

Conversely, if at least one position in every pair is retained,
the remaining differences $X_i-Y_i$ can still be assigned
independently and arbitrarily. Formula (I.58) shows that the
restricted full matrix image equals the original full matrix
image, not merely its cone closure. This proves the exact
criterion and the minimum $h$: retaining $X_1,\ldots,X_h$
alone attains it.

The counterexample lies in the actual final rank-one test
class. Append $L=e_1e_w^T$ and $R=e_we_1^T$ to the word.
Its full product is (I.60) times $e_1e_1^T$ and has generic
rank one. Every prefix beginning at $L$ is nonzero, since
otherwise the full product would vanish, and therefore has
rank one. Every interval wholly inside $\widetilde W$ has
rank at least two. Thus the shortest rank-one suffix of the
prefix just before $R$ begins at $L$, and adjoining $R$
gives the full word as a local final test. Its entire internal
higher-rank profile is already preserved at the common center.

\paragraph{4. One dense pair still detects the scalar.}
Choose distinct interpolation nodes for the $2h$ original
variables, and denote their Lagrange polynomials by
$\lambda_{x_i}$ and $\lambda_{y_i}$. The common-center
one-pair substitution gives
\[
 \prod_{i=1}^h(x_i-y_i)
   \longmapsto
 z^h\prod_{i=1}^h
       \bigl(\lambda_{x_i}(y)-\lambda_{y_i}(y)\bigr)\ne0.
 \tag{I.61}
\]
Each difference is nonzero because its value at the node
for $x_i$ is one. Their product is nonzero in the parameter
polynomial ring over every working field with these distinct
nodes. Numerical evaluation uses an adequate grid or extension,
as in I-21; no assertion is made that every node value detects
this product.

This separates the refuted exact coordinate-deletion criterion
from dense nonvanishing. In particular it rules out a polynomial
dimension charge for uniformly preserving the full matrix cone
by retaining original variable positions at every preserved
common center. It does not rule out a smaller bound after a
different dense shift, a weaker sufficient invariant than full
cone equality, or a more efficient sampling argument. The
constant-center difference cancellation is explicit, not an
assertion about every possible preserved center. $\square$

\paragraph{Derivation record.}
The paired lift and the common-center cone obstruction were
derived by the purity track from the independently proved
I-24 matrices. The result is a deduction within this report;
no literature-priority claim is made.


\subsection{I-33. Efficient prime cocharacters for two dense products, with a product-closure repair}
\label{I-33}

\paragraph{Scope and relation to known PIT.}
This construction targets efficient sampling rather than another
support bound.  Let $n\ge1$ and let $\delta\ge1$ be an integer;
the case $\delta=0$ consists only of constants.  Consider
\[
 f(\mathbf X)=\prod_{i=1}^n F_i(X_i)-\prod_{i=1}^n G_i(X_i),
 \qquad \deg F_i,\deg G_i\le\delta.
 \tag{I.62}
\]
It is a width-two diagonal ROABP, but its expanded support
may be exponential in $n$.  Polynomial PIT for this class
is already known, for example from I-B06.  Bisht, Gupta,
and Volkovich (2023), Theorem~1 (printed p.~9:4) and
Theorem~48 (printed p.~9:20), give a broader positive result
for sums of two products of read-once formulas.  We claim
no new PIT complexity for that established class and no
historical priority for the particular construction below.
The point is an explicit small family of compatible
integer cocharacters, a proof using divisors rather than
expanded monomials, and an exact test of its product closure.

Set
\[
 H=2\delta(n-1),\qquad M=H+2n+2,\qquad B=16M^2.
\]
Choose any $2n$ distinct integer primes
$p_1,\ldots,p_n,q_1,\ldots,q_n$ in $(H,B]$.
For $j=0,1,\ldots,n$, define positive integer weights
\[
 a_i^{(0)}=p_i,\qquad
 a_i^{(j)}=\begin{cases}q_j&i=j,\\p_i&i\ne j.\end{cases}
 \tag{I.63}
\]
These primes exist and can be found in polynomial time in
$n,\delta$ by an elementary sieve, as proved below.
Over \emph{every} field, every nonzero $f$ in (I.62) satisfies
\[
 f(T^{a_1^{(j)}},\ldots,T^{a_n^{(j)}})\ne0
 \quad\hbox{for at least one }0\le j\le n.
 \tag{I.64}
\]
No prime needs to be excluded when it equals the characteristic.
Each substituted polynomial has degree at most $n\delta B$.
Hence evaluating all $n+1$ curves on any $n\delta B+1$
distinct scalars gives a hitting set of size at most
\[
 S=(n+1)(n\delta B+1),
 \tag{I.65}
\]
with polynomial integer weights and polynomial query count.
As elsewhere, a working extension is required when the
coefficient field has too few scalar values.

\paragraph{1. The divisor separation lemma.}
Let $R_i(X)=F_i(X)/G_i(X)$ be nonzero rational functions,
with numerator and denominator degrees at most $\delta$,
and let $p_i$ be distinct primes greater than $H$.
If
\[
 \prod_i R_i(T^{p_i})\in\overline F^{\times},
 \tag{I.66}
\]
then each $R_i(X)$ is $c_iX^{b_i}$ for some nonzero
constant $c_i$ and integer $|b_i|\le\delta$.
Work over an algebraic closure and cancel numerator and
denominator common factors first.

If the characteristic equals one of the primes, say $p_k$,
then a nonzero root or pole of $R_k(T^{p_k})$ has nonzero
integer multiplicity of absolute value at least $p_k$.
At the same point, each other $R_i(T^{p_i})$ has multiplicity
of absolute value at most $\delta$, since the other power
maps are unramified away from zero.  Their combined
multiplicity has absolute value at most $(n-1)\delta<p_k$.
It therefore cannot cancel the first multiplicity, contrary
to (I.66).  Thus $R_k$ has no nonzero root or pole and is
a constant times a monomial.  There is at most one such
characteristic prime.

For any remaining $i$, a nonzero root or pole of $R_i$
pulls back to an orbit $O=\alpha\mu_{p_i}$ of $p_i$
distinct points.  Its multiplicity is nonzero at every
point of $O$.  For $j\ne i$ with $p_j$ different from the
characteristic, the nonzero divisor of $R_j(T^{p_j})$
is supported on at most $2\delta$ cosets of $\mu_{p_j}$.
Each such coset intersects $O$ in at most one point, since
$\gcd(p_i,p_j)=1$.  A characteristic-prime factor, if
present, was already shown to have no nonzero divisor.
Consequently all other factors together can cancel the
divisor on at most $2\delta(n-1)=H<p_i$ points of $O$.
This contradicts (I.66).

Every $R_i$ therefore has divisor supported only at zero
and infinity.  Factoring its numerator and denominator
over the algebraic closure gives $R_i(X)=c_iX^{b_i}$.
The degree bounds give $|b_i|\le\delta$.  This proves
the lemma, including the inseparable characteristic case.

\paragraph{2. The family separates the remaining monomial relation.}
If one product in (I.62) is zero and the other is nonzero,
every positive-weight curve keeps the nonzero product
nonzero, so (I.64) is immediate.  Otherwise both products
are nonzero and we may divide by the second.
If the base curve $j=0$ kills $f$, the divisor lemma gives
$R_i(X_i)=c_iX_i^{b_i}$, and the resulting Laurent identity
forces
\[
 \prod_i c_i=1,\qquad\sum_i b_ip_i=0.
\]
If every $b_i=0$, then $f=0$, contrary to the hypothesis.
Choose $j$ with $b_j\ne0$.  Changing only its prime changes
the exponent sum by the nonzero integer
$b_j(q_j-p_j)$.  Thus the $j$th curve does not kill $f$.
This is an equality of exponents of the formal indeterminate
$T$, so it is valid even when that integer is divisible by
the characteristic.  This proves (I.64).

The same argument shows that if the two component products
are linearly independent, at least one curve preserves
their two-dimensional coefficient span.  If their ratio is
not a monomial product, the base-curve ratio cannot be
constant by the divisor lemma.  If it is a nonconstant
monomial product, some weight vector has nonzero exponent
sum.  In either case the two univariate images are linearly
independent.  This is a concrete instance of the compression
condition I-14; no basis-isolating-weight claim is needed.

The weight family has full affine span: its differences
from the base vector are $(q_j-p_j)e_j$.  Encoding
$j\in\{0,\ldots,n\}$ by $\lceil\log_2(n+1)\rceil$ bits
gives an explicit nonlinear decoder using coordinate
indicator functions.  It is therefore outside the bounded
affine and fixed-degree decoder obstruction of I-15.
It is not obtained from Hoza's path-isolation theorem.

\paragraph{3. An elementary polynomial prime bound.}
For even $B$, the central binomial coefficient satisfies
\[
 \binom{B}{B/2}\ge\frac{2^B}{B+1}.
\]
For each prime $p\le B$, its valuation in this binomial
coefficient is at most $\lfloor\log_p B\rfloor$: each
term in the factorial-valuation difference is zero or one.
Hence its prime-power contribution is at most $B$, and
\[
 B^{\pi(B)}\ge\frac{2^B}{B+1},\qquad
 \pi(B)\ge\frac{B\log2-\log(B+1)}{\log B}.
\]
For $B=16M^2$ and $M\ge2$, the elementary inequalities
$\log(B+1)\le(B/2)\log2$ and $\log B\le4M\log2$
give $\pi(B)\ge2M$.  At most $H$ primes are at most $H$,
so more than $2n$ primes lie in $(H,B]$.
Sieve this polynomially bounded interval and take its first
$2n$ primes.  This avoids using an unproved prime-distribution
hypothesis or a bit-length-only prime bound.

\paragraph{4. The raw curve family is not closed under products.}
The preceding success cannot be extended to general width
two just by multiplying separate successful factors.
For the explicit construction with $n=64$ and $\delta=1$,
one has $B=16\cdot256^2$.  Split the variables into two
blocks of size $32$.  Within either block the $2^{32}$
subset sums of the base primes take at most $32B+1$
integer values, and $2^{32}>32B+1$.  Two different subsets
therefore have the same sum.  After canceling their common
indices, this gives a nonzero multilinear binomial on that
block which vanishes along the base-prime curve.
Let the two disjoint-block binomials be $A$ and $C$.
Their product $AC$ is nonzero and multilinear.

Changing one coordinate's prime affects at most one of
$A,C$; the other remains zero along that curve.  Thus
$AC$ vanishes along all $n+1$ curves in (I.63).
It is nevertheless a width-two ROABP in a suitable order:
each binomial is the difference of two monomial products,
computed by diagonal $2\times2$ factors, and their two
scalar computations concatenate through a constant
rank-one bridge.  Absorbing that bridge into an adjacent
univariate layer gives an ordinary width-two ROABP.
This is a counterexample to the specified raw family,
not to the known width-two hitting-set theorem I-B06.

\paragraph{5. A concrete product-closure repair.}
Let $\mathcal H$ be the list of at most $S$ points from
(I.65).  Interpolate a polynomial curve
$\Gamma(Z)\in E[Z]^n$ of coordinate degree at most $S-1$
through that list at $S$ distinct nodes, repeating points
if necessary.  Every nonzero factor of the form (I.62)
stays nonzero on $\Gamma$, because some listed point hits it.
Therefore any nonzero product of such factors also stays
nonzero, even when different factors share variables.
For a product of total degree at most $D\ge1$, its
composition has degree at most $D(S-1)$ and is hit on any
$D(S-1)+1$ distinct scalars.  A working field with at least
that many elements supplies both the interpolation nodes
and this final grid.

This is the scalar product-closure mechanism of AGKS,
Lemma~33, recorded in I-B06, applied to the new explicit
prime-curve family.  It repairs the precise four-term
counterexample above and gives a fully polynomial sampler
for products of the dense two-product factors.  It does
not handle arbitrary sums of three or more products,
general noncommuting matrix blocks, or their coefficient
matroids.  Those cancellation patterns remain outside
the divisor-separation argument.
\hfill$\square$


\subsection{I-34. Prime cocharacters with a bounded local nullity}
\label{I-34}

\paragraph{Statement and scope.}
Let $n\ge1$, let $\delta\ge1$ be an integer, and consider
\[
 f(\mathbf X)=\sum_{j=1}^{w}\prod_{i=1}^{n}F_{ij}(X_i),
 \qquad 0\ne F_{ij}\in F[X_i],\quad \deg F_{ij}\le\delta.
 \tag{I.67}
\]
Zero product terms can first be omitted. Scalars and signs can be
absorbed into the factors. Let $A_i$ be the $(\delta+1)\times w$
matrix whose $j$th column is the coefficient vector of $F_{ij}$.
We call $w-\operatorname{rank}A_i$ its local nullity.

There is an explicit family of polynomial-magnitude positive
integer cocharacters with the following properties, over every
characteristic.

If some $A_i$ has local nullity at most $c$, the family changing
at most $k=\min\{n,\max\{1,c\}\}$ coordinates of a base prime
weight vector hits every nonzero $f$ in (I.67). For fixed $c$
this gives polynomial sampling size even when $w$ grows. The
promise itself requires $w-c\le\delta+1$; the assertion does not
silently treat $c$ as constant for arbitrary representations.

Without a local-nullity promise, sums of at most three separated
products require only one changed coordinate. More generally,
sums of $w\ge3$ products are hit using at most
$\min\{n,w-2\}$ changed coordinates. The last bound has an
exponent depending on $w$ and is not a polynomial-in-width PIT
result. Constant-top-fan-in PIT is already known, for example
from I-B08. The purpose here is to test and repair the specific
prime-cocharacter construction, not to claim a new general
complexity bound for an established constant-width class.

One-coordinate variation cannot extend this result to all sums
of four products: an explicit tensor-rank-four multilinear
example below vanishes on every entire one-coordinate cylinder
around an admissible prime curve. Thus increasing the number of
alternate primes while still varying only one coordinate does
not remove that obstruction.

\paragraph{1. The explicit family and its sampling cost.}
Set
\[
 H=2\delta(n-1),\qquad
 M=H+n(\delta+1)+2,\qquad B=16M^2.
\]
Choose $n(\delta+1)$ distinct primes in $(H,B]$, organized as
$p_i$ and $q_{i,a}$ for $i\in[n]$, $a\in[\delta]$.
The elementary prime count in I-33 applies verbatim: it gives
$\pi(B)\ge2M$, and at most $H$ primes lie below the threshold.
All these primes can be found by a polynomial-time sieve.

Let $\mathcal W_k$ contain the base vector $(p_1,\ldots,p_n)$
and every vector obtained by replacing $p_i$ by one of
$q_{i,1},\ldots,q_{i,\delta}$ on at most $k$ coordinates.
Then
\[
 |\mathcal W_k|=
 \sum_{a=0}^{k}\binom na\delta^a.
 \tag{I.68}
\]
Each curve $X_i=T^{\omega_i}$ has substituted degree at most
$n\delta B$. A scalar grid of $n\delta B+1$ distinct values
therefore gives a hitting set of size at most
\[
 (n\delta B+1)\sum_{a=0}^{k}\binom na\delta^a.
 \tag{I.69}
\]
An adequate working extension is understood if the coefficient
field is too small. These are degree bounds in the formal
indeterminate $T$, so no characteristic restriction on the
integer weights is needed.

The reason to use $\delta$ alternate primes, rather than just
one, is exact interpolation over $F(T)$. For $S\subseteq[n]$
write
\[
 f_S(\mathbf Z_S,T)=
 f\bigl((Z_i)_{i\in S},(T^{p_i})_{i\notin S}\bigr).
 \tag{I.70}
\]
This is the coordinate cylinder with the coordinates in $S$
free. It has individual degree at most $\delta$ in each $Z_i$.
The $\delta+1$ elements
$T^{p_i},T^{q_{i,1}},\ldots,T^{q_{i,\delta}}$ are distinct
in $F(T)$. Consequently, if every curve from $\mathcal W_k$
kills $f$, tensor-product interpolation over $F(T)$ shows that
$f_S=0$ for every $|S|\le k$. Conversely, if all these cylinders
are zero, every such curve is zero. This equivalence uses only
formal roots and is valid in all characteristics.

\paragraph{2. Nullity one: the circuit argument.}
We prove that if one $A_i$ has nullity at most one, then a
nonzero $f$ has a nonzero one-coordinate cylinder. Suppose
instead that every such cylinder is zero. Put
\[
 B_j(T)=\prod_{\ell\ne i}F_{\ell j}(T^{p_\ell}).
\]
Every $B_j$ is nonzero. The identity $f_{\{i\}}=0$ says
$A_i(B_1,\ldots,B_w)^T=0$ over $F(T)$.
If $A_i$ has full column rank this is impossible. Otherwise
its kernel is one-dimensional and defined over $F$, so
\[
 B_j(T)=h(T)c_j\quad(1\le j\le w),
 \qquad h\ne0,\quad c_j\in F^\times.
 \tag{I.71}
\]
The nonvanishing of every $c_j$ follows from that of every
$B_j$. Thus the columns of $A_i$ form a circuit: their unique
relation, up to a scalar, has full support.

For each $j$, the ratio $B_j/B_1$ is constant. Apply the
divisor separation lemma I-33 to the $n-1$ remaining variables.
Its threshold is no larger than our $H$; empty products when
$n=1$ cause no problem. We obtain
\[
 \frac{F_{\ell j}(X)}{F_{\ell1}(X)}
       =d_{\ell j}X^{b_{\ell j}}
 \quad(\ell\ne i),
 \qquad d_{\ell j}\in F^\times,
 \quad b_{\ell j}\in\mathbb Z.
 \tag{I.72}
\]
The constants lie in $F$: the ratio has coefficients in $F$
and equals a constant times a monomial. Negative exponents
are allowed in this rational-function identity. Set
$b_{\ell1}=0$.

Fix $\ell\ne i$ and use the zero cylinder with $X_\ell=Z$.
Dividing by nonzero rational functions and using (I.71)--(I.72)
reduces it to the Laurent identity
\[
 \sum_{j=1}^w c_j F_{ij}(T^{p_i})
       (Z/T^{p_\ell})^{b_{\ell j}}=0.
 \tag{I.73}
\]
For every nonempty group $J_b=\{j:b_{\ell j}=b\}$, its
coefficient gives
$\sum_{j\in J_b}c_jF_{ij}(T^{p_i})=0$.
The substitution $X\mapsto T^{p_i}$ is injective on $F[X]$,
so $\sum_{j\in J_b}c_jF_{ij}(X)=0$.
The vector equal to $c_j$ on $J_b$ and zero elsewhere is
therefore in $\ker A_i$. Since that kernel is generated by
the full-support vector $c$, every nonempty $J_b$ must be
all of $[w]$. In particular all $b_{\ell j}$ equal
$b_{\ell1}=0$.

It follows that every off-$i$ factor is a constant multiple
of its first-column factor. Equation (I.71) then implies that
$f$ is the product of these common off-$i$ factors with the
zero polynomial $\sum_jc_jF_{ij}(X_i)$, up to a common
nonzero scalar. Thus $f=0$, a contradiction. This proves
the nullity-one assertion. The proof includes the case where
a proposed one-dimensional kernel has a zero coordinate:
then (I.71) cannot occur, and the $i$th cylinder is already
nonzero.

\paragraph{3. Three terms and the general local-nullity reduction.}
For three nonzero product terms, either some $A_i$ has rank
at least two, when the preceding argument applies, or every
$A_i$ has rank one. In the second case every factor list is
proportional, and $f$ is a scalar multiple of one separated
product. If nonzero, it survives the base curve. The cases
of one or two terms follow directly as well. Hence every
nonzero sum of at most three products has a nonzero cylinder
with at most one free coordinate.

For later use, we prove by induction that a sum of at most
$s\ge3$ separated products has a nonzero cylinder with at
most $s-2$ free coordinates (or all $n$ coordinates if fewer
are available). The base $s=3$ was just proved. If every
local factor matrix has rank one, the assertion again
reduces to one product. Otherwise choose a coordinate $i$
whose local rank is $r\ge2$, and choose $r$ of its factors
as a basis. Expanding the other factors in this basis gives
\[
 f=\sum_{a=1}^r F_{i,j_a}(X_i)Q_a(\mathbf X_{\ne i}),
 \tag{I.74}
\]
where each $Q_a$ is a sum of at most $s-r+1\le s-1$
separated products: the original pivot term and at most one
term for each nonpivot column. If every cylinder of $f$
with at most $s-2$ free coordinates were zero, independence
of the basis factors would imply that every $Q_a$ vanished
on cylinders with at most $s-3$ free coordinates. The
induction hypothesis would force every $Q_a=0$, contrary
to $f\ne0$. The prime threshold inherited by a subproblem
on fewer variables is still sufficient.

Now suppose a specified $A_i$ has nullity $c\ge2$.
Expansion (I.74) makes each $Q_a$ a sum of at most $c+1$
products. The result just proved detects each nonzero $Q_a$
with at most $c-1$ free coordinates. Adding $i$ gives at
most $c$ free coordinates for $f$. Nullity zero is handled
by freeing $i$ alone, and nullity one by paragraph~2. This
proves all the stated cylinder bounds and, by paragraph~1,
the explicit cocharacter and sampling bounds.

The same finite family jointly preserves the coefficient
span of the vector of product terms under the local-nullity
promise. More precisely, if $V\subseteq F^w$ is that
coefficient span and the sampling grid lies in a working
extension $E$, the sampled vectors span $V\otimes_F E$
over $E$. Indeed, every $E$-linear combination is still
such a sum over $E$ after scaling or deleting columns of
$A_i$; column deletion cannot increase its nullity, and
scalar extension preserves the original matrix rank.
If an $E$-linear functional annihilates every sampled vector,
the scalar hitting theorem applied over $E$ forces it to
annihilate $V\otimes_F E$. The reverse inclusion follows
because every evaluated vector is a combination of the
original coefficient vectors. In particular no original-field
functional nonzero on $V$ is lost. This is the exact
joint-span criterion I-14. We do not assert that one member
of the family is always a BIWA or preserves the entire span
individually.

\paragraph{4. An actual additive cancellation of three dense products.}
The multiplicative conclusion in I-33 does not extend to an
additive relation among three products. Let $n=6$, $\delta=12$,
and assign the following distinct prime weights to variables
$A,B,C,Q,R,S$, respectively:
\[
 (1663,523,139,127,1093,331).
\]
All exceed $H=120$ and are below the polynomial bound above.
The polynomial
\[
 \begin{aligned}
 F={}&(1-A)(1-S^2)-(1-R^2)(1-C)\\
     &{}-C(1-B)(1-Q^{12})
 \end{aligned}
\]
is nonzero: the monomial $AS^2$ has coefficient one.
Its three displayed terms are separated products, and no
pairwise ratio is a monomial in the original variables.
Nevertheless this prime curve kills it identically in every
characteristic. To verify this, put $a=1663$, $b=523$,
$c=139$. The remaining weights satisfy
$2\cdot331=b+c$, $2\cdot1093=a+b$, and
$12\cdot127=a-c$. The elementary identity
\[
 (1-U)(1-VW)-(1-UV)(1-W)=(1-V)(W-U)
\]
with $U=T^a$, $V=T^b$, $W=T^c$ proves the assertion.

The repaired family detects this example. Freeing $A$ and
holding the other variables on their prime curve gives
$(T^{1663}-A)(1-T^{662})\ne0$. This example refutes the
specific proposed extension of the divisor lemma; it is
not a lower bound for sparse PIT or for the repaired family.

\paragraph{5. Nullity two already defeats every single-coordinate repair.}
Let $n=8$, $\delta=1$, and take the first eight primes
greater than $H=14$:
\[
 (17,19,23,29,31,37,41,43).
\]
Set
\[
 F=(X_1X_8-X_2X_7)(X_3X_6-X_4X_5).
\]
Both binomials vanish on the base curve, because
$17+43=19+41=60$ and $23+37=29+31=60$.
On any cylinder freeing one coordinate, at least one of
the two disjoint binomials remains unchanged and zero.
Thus every one-coordinate cylinder is identically zero,
including all substitutions that change that coordinate
to any alternate prime power or any rational function.

The polynomial has tensor rank exactly four in the separated
product model. Expansion gives the upper bound four. For
the lower bound, partition the variables as
$\{1,8,3,6\}$ and $\{2,7,4,5\}$.
Each binomial has a rank-two coefficient matrix across its
corresponding partition, and the product on disjoint
variables gives their Kronecker product, of rank four.
Every separated product has coefficient-matrix rank one,
so at least four terms are necessary. This argument works
in every characteristic. In the displayed four-term
representation each local factor matrix has rank two and
nullity two.

Freeing, for example, $X_1$ and $X_3$ makes both binomials
nonzero, so the two-coordinate repair proved above removes
this precise obstruction. The next unresolved efficiency
question is to avoid a sampling exponent proportional to
unbounded local nullity. Polynomially many alternate primes
alone cannot do that while only one coordinate changes.
No conversion of this repair into a Hoza path-isolation
construction has been proved.\hfill$\square$

\paragraph{Derivation record.}
The full proof is a deduction in this track from the I-33
divisor lemma and elementary coefficient elimination.
The specific prime numbers, the three-term identity, and
the tensor-rank-four obstruction were checked directly.
No literature-priority claim is made.

\subsection{I-35. A smaller contraction invariant, its exact composition rule, and the rank-one lift obstruction}
\label{I-35}

\paragraph{Scope.}
This result concerns the rank-$r$ outer matrix products left after
factoring higher-rank resets. It separates three assertions:
preserving all constant scalar boundary contractions of one bridge,
composing bridge substitutions with independent parameters, and
synchronizing those parameters. The first is strictly weaker than
preserving coefficient span or generic matrix rank. The second has
an exact positive proof. The third does not follow, and explicit
counterexamples below identify the missing quantifier. No general
ROABP hitting set is claimed.

\paragraph{1. The exact rank-one-annihilator criterion.}
Let $M(\mathbf X)$ be an $r\times r$ polynomial matrix over a field
$K$, and let $U\subseteq\operatorname{Mat}_r(K)$ be the span of its
coefficient matrices. After a polynomial substitution $G$, write
$U_G$ for the coefficient-matrix span of $M(G)$; then $U_G\subseteq U$.
All statements can instead be made over a specified working extension.
For rational substitutions, clear a common nonzero polynomial
denominator first. The span is independent of this clearing choice:
a constant functional annihilates the cleared vector if and only if
it annihilates the original rational vector.

Use the pairing $\langle C,A\rangle=\operatorname{tr}(C^{\mathsf T}A)$,
and let $\Sigma_1$ be the set of matrices of rank at most one. Then
\[
 \begin{gathered}
 uM(\mathbf X)v\ne0\ \Longrightarrow\ uM(G)v\ne0
       \quad\hbox{for every constant row }u\hbox{ and column }v
 \\
 \Longleftrightarrow\qquad
 U_G^\perp\cap\Sigma_1=U^\perp\cap\Sigma_1.
 \end{gathered}
 \tag{I.75}
\]
Indeed, $uAv=\langle u^{\mathsf T}v^{\mathsf T},A\rangle$, and every
nonzero rank-one matrix has this form. A scalar contraction is the
zero polynomial exactly when its functional annihilates the entire
coefficient space. Since $U_G\subseteq U$, one inclusion of the two
intersections is automatic and the other is precisely nonvanishing
preservation. For preservation after every scalar extension, impose
this equality over an algebraic closure. Thus the target is equality
of base loci on the Segre variety of boundary pairs, rather than
equality of the entire linear spaces or matrix image cones.

\paragraph{2. An explicit optimal test space of dimension $2r-1$.}
Index rows and columns by $0,\ldots,r-1$. Let $H_k$ have entry one
when $i+j=k$ and zero otherwise, for $0\le k\le2r-2$. Their disjoint
nonempty supports make them linearly independent. Put
\[
 a(T)=(1,T,\ldots,T^{r-1})^{\mathsf T},\qquad
 H(T)=a(T)a(T)^{\mathsf T}=\sum_{k=0}^{2r-2}T^kH_k.
 \tag{I.76}
\]
For any nonzero constant $u,v$,
\[
 uH(T)v=
   \left(\sum_{i=0}^{r-1}u_iT^i\right)
   \left(\sum_{j=0}^{r-1}v_jT^j\right)\ne0.
 \tag{I.77}
\]
This is a product of two nonzero polynomials over a domain, in every
characteristic. Consequently the $(2r-1)$-dimensional Hankel space
$\mathcal H_r=\operatorname{span}\{H_k\}$ has no nonzero rank-one
annihilator. It satisfies (I.75) with $U=\operatorname{Mat}_r(K)$,
although $\mathcal H_r$ is a proper subspace when $r\ge2$.
Evaluating $H(T)$ at $2r-1$ distinct field values spans this space by
Vandermonde interpolation; an extension may supply the values.

The dimension $2r-1$ is optimal over an algebraically closed field.
This is known completely-entangled-subspace background, not a new
claim: Parthasarathy (2004), Propositions~1.3--1.4 and Theorem~1.5
(P-B16), gives this two-factor construction and optimal codimension
over $\mathbb C$. The algebraic-closure formulation follows from
the same Segre dimension argument; its checked proof is retained
in the separate background audit. The lower bound need not hold
when only rational boundary pairs over a nonclosed field are tested.

This is an actual matrix-bridge substitution. Take
$M(\mathbf X)=(X_{ij})$, whose coefficient span is all of
$\operatorname{Mat}_r(K)$, and substitute $X_{ij}=T^{i+j}$.
The generic matrix rank drops from $r$ to one, while (I.77) preserves
every nonzero constant boundary contraction. Moreover this generic
matrix is realized by an invertible width-$2r$ interior word:
\[
 \prod_{i,j}\bigl(I_{2r}+X_{ij}E_{i,r+j}\bigr)
     =\begin{pmatrix}I_r&M(\mathbf X)\\0&I_r\end{pmatrix}.
 \tag{I.78}
\]
Here one may relabel matrix-unit indices from one instead of zero.
All products of two off-diagonal generators vanish, so the identity
holds in every order. Selecting the upper-right block uses constant
full-rank rectangular endpoints. Thus generic rank preservation is
strictly stronger than scalar boundary preservation even for an
actual bridge with invertible interior factors.

More generally, after scalar extension to an algebraic closure,
every coefficient space $U$ has a subspace $W\subseteq U$ of
dimension at most $2r-1$ with the same rank-one annihilators.
To prove this, remove the base locus from the Segre
variety and map a boundary pair to its nonzero evaluation functional
in $\mathbb P(U^*)$. The closure $Y$ of this image has dimension at
most $2r-2$. If $\dim U\le2r-1$, take $W=U$. Otherwise set
$D=\dim U$ and $c=2r-1$. The Grassmannian of $(D-c)$-dimensional
vector subspaces of $U^*$ has dimension $c(D-c)$. Requiring such
a subspace to contain one fixed projective point imposes $c$
conditions: its fiber is a Grassmannian of dimension
$c(D-c-1)$. The incidence variety of subspaces meeting $Y$
therefore has dimension at most
$\dim Y+c(D-c-1)<c(D-c)$. Its image is closed because $Y$ is
projective, so a nonempty open set of subspaces misses $Y$.
The annihilator of any one of them defines the desired $W$.
This proves existence of a small contraction test space, not
an explicit coefficient-oblivious substitution realizing it for
an arbitrary bridge. The Hankel construction realizes it for the
full generic matrix example only. No field-of-definition or
effective-construction bound is asserted for the general existence
statement.

\paragraph{3. The valid composition statement and its independent-seed cost.}
Let $B_j(\mathbf X_j)$, $1\le j\le s$, be square matrices with
pairwise disjoint variable sets. Suppose each substitution
$\mathbf X_j=G_j(\mathbf S_j)$ preserves every constant scalar
boundary contraction of $B_j$, and does so after every field extension.
If the seed sets $\mathbf S_j$ are algebraically independent of one
another and of all remaining variables, then
\[
 uB_1(\mathbf X_1)\cdots B_s(\mathbf X_s)v\ne0
 \ \Longrightarrow\
 uB_1(G_1)\cdots B_s(G_s)v\ne0.
 \tag{I.79}
\]
Replace the matrices successively. At step $j$, take as coefficient
field the rational function field in all unprocessed variables and
all previously introduced seeds. The product on the left of $B_j$
is a constant row over this field, and the product on its right is
a constant column. The current scalar polynomial is nonzero, so
the hypothesis over that extension field preserves its nonzeroness
when the fresh seed set $\mathbf S_j$ is introduced. Induction proves
(I.79). This argument requires no equality of coefficient spaces.

If adjacent bridges share cut variables $\mathbf Z$, first work
over the field $K(\mathbf Z)$ and leave those variables unassigned.
Assume the bridge preservation hypothesis over this field and all
its extensions. The same proof then applies to substitutions of
their disjoint interior variables. A later specialization of $\mathbf Z$ must still preserve
the resulting nonzero numerator and any relevant denominators.
Nothing in (I.79) supplies that specialization automatically.

For a prescribed decomposition into $s$ blocks, suppose each map
is polynomial over $K$ in its seeds and the retained cut variables,
uses at most $d$ seeds, and has total coordinate degree at most $L$.
If $z$ shared cut variables remain free, then the resulting
original scalar polynomial of total degree at most $D$ has at most
$sd+z$ parameters and total degree at most $D\max\{1,L\}$.
A full grid therefore costs at most
\[
 \bigl(D\max\{1,L\}+1\bigr)^{sd+z}.
 \tag{I.80}
\]
This displays the cost rather than concealing it in a genericity
assertion. The argument presupposes the decomposition; it does not
produce a uniform unknown-order sampler or identify the reset
locations from a scalar black box. For maps with arbitrary
$K(\mathbf Z)$ coefficients instead, numerator and denominator
degrees require separate bounds; (I.80) is not asserted for them.

\paragraph{4. Why constant contraction preservation does not synchronize.}
Even preserving each entire local coefficient space does not justify
using the same seed in different blocks. For example, with
$u=e_1^{\mathsf T}$ and $v=e_2$,
\[
 u(I+XE_{12})(I-YE_{12})v=X-Y\ne0.
 \tag{I.81}
\]
Each substitution $X=T$ or $Y=T$ separately preserves the full
coefficient space of its invertible matrix factor. Simultaneously
making both substitutions kills the scalar contraction.

The stronger Hankel example makes the boundary quantifier explicit.
For $r=2$, the map $H(T)$ in (I.76) preserves every nonzero constant
boundary contraction, but
\[
 (-T,1)H(T)e_1=0.
 \tag{I.82}
\]
With an independent boundary parameter $S$, the corresponding
expression is $T-S\ne0$. Diagonal identification $S=T$ creates the
cancellation. Treating all constant boundary pairs over extensions
is not a license to identify a new seed with a preexisting element
of the coefficient field.

In fact an unrestricted same-seed rational-boundary condition would
be impossible for $r\ge2$. Given any matrix $A(\mathbf S)$ over
$K(\mathbf S)$, there are nonzero $u(\mathbf S),v(\mathbf S)$ with
$uAv=0$: choose a nonzero row and then a nonzero vector in the
kernel of its resulting row functional. For the generic matrix
$M(\mathbf X)$, however, $uM(\mathbf X)v$ is nonzero over
$K(\mathbf S)$ because at least one product $u_iv_j$ is nonzero.
Thus a viable synchronization theorem must use restrictions on the
boundary functions actually induced by neighboring bridges. Asking
for all rational functions of the same seed is too strong.

\paragraph{5. A rank-one multiplicative lift cannot retain an entry for $r\ge2$.}
There is also a precise obstruction to turning the higher-rank
outer product into scalar product closure by a multiplicative lift.
Let $r\ge2$, let $w\ge r$, and embed $r\times r$ matrices as the
upper-left corner $\iota(B)$ of $w\times w$ matrices. Put
$P=\iota(I_r)$. Suppose a multiplicative map
\[
 \rho:\operatorname{Mat}_w(K)\longrightarrow\operatorname{Mat}_N(E)
\]
is given over some extension $E$, and suppose $\rho(P)$ has rank
one. There cannot be a fixed $E$-linear readout $\Lambda$ satisfying
\[
 \Lambda(\rho(\iota(B)))=B_{11}
 \quad\hbox{for every }B\in\operatorname{GL}_r(K).
 \tag{I.83}
\]
No polynomiality or linearity of $\rho$ is needed for this obstruction;
multiplicativity and the fixed linear readout suffice.

Indeed $Q=\rho(P)$ is a rank-one idempotent. Since
$P\iota(B)P=\iota(B)$, every lifted corner matrix lies in
$Q\operatorname{Mat}_N(E)Q=EQ$ and has the form $\sigma(B)Q$.
The identity case of (I.83) gives $\Lambda(Q)=1$, so
$\sigma(B)=B_{11}$. Multiplicativity would then make $B\mapsto B_{11}$
a multiplicative function on $\operatorname{GL}_r(K)$. But for
\[
 B=I_r+E_{12},\qquad C=I_r+E_{21},
\]
one has $B_{11}=C_{11}=1$ and $(BC)_{11}=2\ne1$, in every
characteristic. This is a contradiction. Multiplying the proposed
readout by any integral power of $\det B$ does not evade it, because
$B,C,BC$ all have determinant one.

Exterior powers are consistent with this obstruction: at a rank-$r$
reset, the $r$th exterior power has rank one and retains maximal-minor
data, which acts by the determinant on an invertible $r$-dimensional
corner. It cannot also recover an arbitrary fixed entry there by
such a readout. This excludes the stated functorial scalarization,
not all nonfunctorial or coefficient-dependent constructions.

\paragraph{Remaining concrete gate.}
The relevant smaller target is the Segre-annihilator condition for
the actual contractions, together with a synchronization theorem
for the restricted boundary functions produced by adjacent reset
bridges. The independent-seed statement and its cost are proved
above. The Hankel bridge proves that this target can preserve scalar
information while destroying generic matrix rank. Neither fact
establishes an efficient synchronized map for general rank-$r$
outer words. That last assertion remains open in this report.
\hfill$\square$

\paragraph{Background and provenance.}
The optimal Hankel test-space statement is the known two-factor
case of Parthasarathy (2004), Propositions~1.3--1.4 and Theorem~1.5,
preprint pp.~4--6, already recorded in P-B16. The complete local
proofs were independently reopened and read for this audit.
The projective facts used here are the Segre embedding, Stacks
Project Authors (2026), Lemma~27.13.6, tag~01WD, and the projective
intersection dimension theorem, Lemma~33.34.3, tag~0B2R. Their
primary statements and displayed proofs were read. The Hankel
identity, actual unipotent bridge realization, composition rule,
synchronization counterexamples, and multiplicative-corner
obstruction are proved directly above. No historical priority
claim or general derandomization consequence is made.

\subsection{I-36. A pairwise dominant center preserves maximal minors; five pairs for singular width-three words}
\label{I-36}

\paragraph{Statement and scope.}
Fix distinct Lagrange nodes $a_1,\ldots,a_n$ in a working field
$E$, and let $\lambda_i(Y)$ be their degree-at-most-$(n-1)$
Lagrange polynomials. The coefficient field is a subfield of $E$.
For $n\ge1$, define
\[
 \alpha_i=t+z\lambda_i(y),\qquad
 \mathcal G_i^{(5)}=t+\sum_{j=1}^{5}z_j\lambda_i(y_j).
 \tag{I.84}
\]
All displayed parameters are algebraically independent.

The map $\boldsymbol\alpha$ preserves every nonzero maximal-size
minor of every interval in a univariate matrix word whose individual
factors have generic rank at most that size. More explicitly, if
for integers $1\le r\le w$, the $A_i(X_i)$ are $w\times w$
polynomial matrices of generic rank at
most $r$, with distinct variables at the different positions, then
every nonzero $r\times r$ minor of every interval product remains
nonzero after $\mathbf X=\boldsymbol\alpha$. Identically zero minors
remain zero. The map depends on neither the matrices nor their order.

Consequently $\mathcal G^{(5)}$ hits every nonzero width-at-most-three
ROABP admitting a representation in which every variable-dependent
matrix factor has generic rank at most two. Constant invertible
width-three factors may be absorbed into neighboring factors or
the boundaries. Arbitrarily many singular factors and arbitrary
drops in the ranks of longer intervals are allowed. The variable
order is unknown to the sampler. This improves I-30's six-pair bound
on this specified singular-factor subclass; it does not prove that
two centered pairs suffice, and it does not extend this improvement
to arbitrary variable-dependent invertible width-three factors.

\paragraph{1. Dominance on every coordinate pair.}
For $i\ne j$, the polynomial $\lambda_i(y)-\lambda_j(y)$ is nonzero:
its value at $y=a_i$ is one. Over $E(y)$, the linear substitution
\[
 \begin{pmatrix}X_i\\X_j\end{pmatrix}
 =\begin{pmatrix}1&\lambda_i(y)\\1&\lambda_j(y)\end{pmatrix}
   \begin{pmatrix}t\\z\end{pmatrix}
 \tag{I.85}
\]
is invertible. It therefore induces an injection
$E[X_i,X_j]\hookrightarrow E(y)[t,z]$. Every nonzero polynomial
depending on at most two specified coordinates stays nonzero.
The singleton case follows directly from the coefficient one on
$t$; it covers $n=1$. The same argument holds after any extension
of the coefficient field independent of the new parameters.

This is a formal polynomial assertion. It does not require a
single numerical parameter choice to work over a small base field;
the later interpolation step supplies the needed finite evaluations.

\paragraph{2. Rank-one words factor into adjacent bivariate contractions.}
Let $B_i(X_i)$ be nonzero polynomial matrices of generic rank one.
Choose any nonzero entry $b_i$ of $B_i$, let $p_i$ be its column,
and let $q_i$ be its row. The vanishing of the two-by-two minors gives
\[
 B_i=\frac{p_iq_i}{b_i}
 \quad\hbox{over the rational function field.}
\]
For every interval $[s,t]$ and constant boundary row $u$ and column
$v$ of the appropriate sizes, clearing the nonzero univariate
pivots gives the exact polynomial identity
\[
 \left(\prod_{i=s}^{t}b_i\right)
       uB_sB_{s+1}\cdots B_tv
 = (up_s)(q_tv)
       \prod_{i=s}^{t-1}(q_ip_{i+1}).
 \tag{I.86}
\]
The two endpoint factors are univariate. Each middle factor depends
on only the two distinct adjacent variables $X_i,X_{i+1}$.
This remains valid for a one-factor interval, with an empty middle
product. The variable sets of different factors on the right may
overlap; their nonzero product is still nonzero in a polynomial
domain.

If the original scalar contraction is nonzero, every factor on the
right of (I.86) is nonzero. Part~1 preserves all those factors and
all the pivots under the same substitution $\boldsymbol\alpha$.
The substituted scalar contraction is therefore nonzero. In
particular every nonzero entry of every rank-one interval word is
preserved. This proof uses scalar product closure after an exact
rank-one factorization, not an unsupported composition of arbitrary
constant-boundary tests with shared parameters.

\paragraph{3. Applying the factorization to maximal minors.}
For a matrix $A$, write $C_r(A)$ for its $r$th compound matrix.
Cauchy--Binet gives
\[
 C_r(A_s\cdots A_t)=C_r(A_s)\cdots C_r(A_t).
 \tag{I.87}
\]
If a factor has generic rank below $r$, its compound is zero and
every interval containing it has all $r$-minors identically zero.
Otherwise all factors in the interval have generic rank exactly
$r$, so each $C_r(A_i)$ has generic rank one. Part~2 applies to
this compound word and preserves each nonzero entry of (I.87),
which is precisely the corresponding $r$-minor of the original
interval. This proves the first assertion of the statement.

Compound matrices occur only in the proof. The sampler does not
construct them, recover their pivot entries, inspect the input
representation, or enumerate the unknown variable order. All
intervals and minors are preserved by the same formal map: the
argument establishes a nonzero image for each one individually,
and any finite product of these nonzero images is nonzero.

\paragraph{4. Completing the rank-one stage in width three.}
Remove constant invertible factors from the width-three word by
absorbing them into neighboring matrices or the constant boundary
vectors. This does not increase the generic rank of any remaining
factor. Unused variables cause no problem for (I.85). A representation
of smaller width may be padded with zero rows and columns to width
three, so its factors also have generic rank at most two.

Apply parts~1--3 with $r=2$. Every nonzero two-by-two interval
minor survives at the center $\boldsymbol\alpha$. There are no
generic rank-three intervals in the remaining word. The two
constant rank-one sentinels used in I-21 to encode the scalar
output introduce no rank-three interval or additional variable.

Now use the local rank-test construction proved in I-21 and
the endpoint absorption of I-30. In each selected nonzero
rank-one test, the interior intervals have generic ranks at
least two. Here their ranks are exactly two, and they have
already been preserved at $\boldsymbol\alpha$. Every nonzero
univariate polynomial endpoint vector is nonzero at its center:
each $\alpha_i$ is transcendental over the original coefficient
field, since its coefficient on $t$ is one. Constant nonzero
endpoints are covered directly.

Pass to the field of the first-stage parameters and translate
to independent residual variables. I-30 absorbs the endpoints
into unit factors, producing a word with matching generic and
constant interval ranks in $\{2,3\}$. The complete support
lemma I-29 with $(w,r)=(3,2)$ supplies a nonzero coefficient
whose monomial has support at most
\[
 2(3-2)+3-1=4,
 \tag{I.88}
\]
including both endpoint positions. Four fresh SV pairs preserve
every selected rank-one test, because a polynomial with such a
nonzero low-support coefficient cannot vanish under $\mathrm{SV}_4$.
The new pairs are independent of the first-stage parameter field.
The earlier two-by-two minors remain nonzero, as is seen by setting
the new amplitudes to zero. The local-rank induction of I-21 now
preserves the nonzero scalar output. Adding the first pair and
the four residual pairs gives exactly $\mathcal G^{(5)}$ in (I.84).

This last argument is the already checked rank-one completion
from I-30, with its rank-two-preserving center supplied by the
new rank-one-compound factorization. It does not assert that
the first-stage center alone preserves rank-one intervals.

\paragraph{5. Exact finite sampling and field requirements.}
For a total-degree-at-most-$D$ input, the substituted polynomial
has degree at most $D$ in $t$, at most $(n-1)D$ in each of the
five selectors, and total degree at most $D$ in the five
amplitudes. The tensor grid in $t$ and the selectors, together
with the Newton simplex in the amplitudes used in I-21, yields
at most
\[
 (D+1)\bigl((n-1)D+1\bigr)^5\binom{D+5}{5}
 \tag{I.89}
\]
evaluations. The working field must contain at least
$\max\{n,D+1,(n-1)D+1\}$ distinct elements. If it is a proper
extension of the coefficient field, black-box evaluation over
that extension is part of the evaluation model. All arguments,
including the simplex interpolation, hold in every characteristic.
The constant-polynomial case $D=0$ needs only one evaluation;
the displayed formula also equals one. No polynomial dependence
on arbitrary width is inferred from this fixed-width result.
\hfill$\square$

\paragraph{Provenance and the remaining gate.}
The rank-one pivot factorization and scalar product closure are
the established mechanism recorded in I-B06 and P-14. The new
deduction here is that pairwise dominance supplies an entire
maximal-minor preservation stage with one centered interpolation
pair, followed by the previously proved four-pair completion.
No historical priority claim is made. The centered two-pair
question for rank-two reset words remains open; the explicit
uncentered failure P-16 does not decide it.

\subsection{I-37. Two centered pairs for affine singular width-three words}
\label{I-37}

\paragraph{Statement and exact promise.}
Let $E$ be any field, let $a_1,\ldots,a_n\in E$ be distinct, and let
$\lambda_i(Y)$ be their Lagrange polynomials. A nonzero polynomial
$f\in E[X_1,\ldots,X_n]$ is assumed to admit a width-at-most-three
ROABP whose variable-dependent factors are affine-linear in their
own variables and have generic rank at most two. Constant factors
are allowed. The order is arbitrary and unknown to the map. Then
\[
 f\left(t+z_1\lambda_1(y_1)+z_2\lambda_1(y_2),\ldots,
        t+z_1\lambda_n(y_1)+z_2\lambda_n(y_2)\right)\ne0.
 \tag{I.90}
\]
This holds in every characteristic. It improves I-36's five-pair
bound on the stated affine subclass. It does not cover a mixture
containing variable-dependent generic-rank-three factors, or
replace the affine hypothesis by arbitrary individual degree.

We prove the slightly more convenient rational statement. For
arbitrary $c_1,\ldots,c_n\in E^\times$, set
\[
 R_{k,i}=t c_i+\sum_{h=1}^{k}\frac{z_h}{a_i-y_h}.
 \tag{I.91}
\]
Every input in the promised class satisfies $f(\mathbf R_2)\ne0$.
All parameters here are algebraically independent; rational
identities are understood in their rational function field.

\paragraph{1. Pairwise dominance and differential consequences.}
Write $\alpha_i=t c_i+z/(a_i-y)$. For $i\ne j$, the determinant
\[
 \frac{c_i}{a_j-y}-\frac{c_j}{a_i-y}
 =\frac{c_i(a_i-y)-c_j(a_j-y)}{(a_i-y)(a_j-y)}
 \tag{I.92}
\]
is nonzero. If its numerator were zero, comparison of its linear
and constant coefficients would give $c_i=c_j$ and
$c_i(a_i-a_j)=0$, a contradiction. The substitution in the two
coordinates is thus an invertible linear change of $(t,z)$ over
$E(y)$. It preserves every nonzero polynomial in at most two
specified coordinates, over any fixed extension of $E$ as well.

Suppose, toward a contradiction, that $f(\mathbf R_2)=0$. Setting
the second amplitude to zero gives $f(\boldsymbol\alpha)=0$.
The coefficient of a new amplitude $\epsilon$ in
$f(\alpha_i+\epsilon/(a_i-u):1\le i\le n)=0$ is
\[
 \sum_i\frac{f_i(\boldsymbol\alpha)}{a_i-u}=0,
 \qquad f_i=\partial f/\partial X_i.
\]
The distinct simple poles force $f_i(\boldsymbol\alpha)=0$ for
every $i$: multiply by $\prod_j(a_j-u)$ and set $u=a_i$.
Since $\boldsymbol\alpha+s\mathbf c$ is obtained by replacing
$t$ by $t+s$, these identities also give
\[
 \sum_j c_j f_{ij}(\boldsymbol\alpha)=0
 \quad\hbox{for every }i,
 \qquad H(\boldsymbol\alpha)\mathbf c=0,
 \tag{I.93}
\]
where $f_{ij}$ denotes the second partial derivative. This uses
only the coefficient of $s$, so it is valid in every characteristic.
For multilinear $f$, every diagonal entry $f_{ii}$ is zero.

\paragraph{2. Irreducible factors inherit the representation promise.}
We may extend the coefficient field to its algebraic closure.
This is a proof device, not a requirement that the final sampler
evaluate over the algebraic closure. A nonzero substitution over
an extension is a nonzero substitution over its original field.
The kernel of the rational substitution $\mathbf X=\mathbf R_2$
is prime, because the target is a field. Hence an assumed nonzero
annihilator has a nonconstant irreducible factor $g$ that is also
an annihilator.

The affine-factor representation makes $f$ multilinear. For a
product of nonzero polynomials, degree in each fixed variable is
additive. Thus the nonconstant irreducible factors of a multilinear
polynomial depend on disjoint variable sets. Specialize all
variables outside the set used by $g$ so that the product of the
other factors is a nonzero scalar. Such a specialization exists
over the infinite algebraically closed field. It realizes a
nonzero scalar multiple of $g$ as a specialization of the original
matrix word. Remove the now constant factors by absorbing their
products into a neighboring factor or a boundary. Multiplication
by constant matrices preserves affine dependence and cannot
increase generic rank. The surviving factors therefore still
have generic rank at most two. A width-two or width-one word may
be padded with zero coordinates to width three.

Consequently it suffices to treat an irreducible $f$ that depends
on every one of its remaining variables, with those variables
listed in their representation order. The corresponding nodes
remain distinct and the corresponding $c_i$ remain nonzero.
The proof does not enumerate or discover this order. If at most
two variables remain, part~1 already contradicts annihilation.
Assume henceforth that $n\ge3$.

\paragraph{3. Every strict internal interval has generic rank two.}
We use an elementary polynomial rank-one factorization. If a
nonzero matrix $B$ over a multivariate polynomial ring over a
field has rank one over its fraction field, then $B=pq$ for a
polynomial column $p$ and polynomial row $q$. Indeed, take a
nonzero polynomial column, divide its entries by their greatest
common divisor, and call the primitive resulting column $p$.
Every column of $B$ is a rational multiple $h p$. In a reduced
fraction for $h$, each irreducible denominator factor would have
to divide every nonzero entry of $p$. Primitivity excludes it,
so $h$ is polynomial. This proves the asserted factorization.
In particular, when $B$ is multilinear, the polynomial factors
$p$ and $q$ depend on disjoint variable sets: the maximum entry
degrees satisfy
$\deg_{X_i}B=\deg_{X_i}p+\deg_{X_i}q\le1$.

Write the irreducible scalar as
$f=u A_1(X_1)\cdots A_n(X_n)v$. For a strict internal interval
$2\le s\le t\le n-1$, the product $B=A_s\cdots A_t$ cannot
be zero, because then $f=0$. If it had rank one, the preceding
factorization would give
\[
 f=(uA_1\cdots A_{s-1}p)(qA_{t+1}\cdots A_nv).
 \tag{I.94}
\]
The first factor depends on $X_1$ and the second on $X_n$,
because $f$ does and each of these variables appears on only
its indicated side. Both factors are therefore nonconstant,
contradicting irreducibility. Since each original factor has
rank at most two, every strict internal interval has rank
exactly two.

The pairwise-dominant map of part~1 preserves every nonzero
two-by-two minor of every such interval. To see this with the
same quantifiers as I-36, take second compounds: each
$C_2(A_i)$ has generic rank one, and Cauchy--Binet turns an
interval minor into an entry of their product. The pivot-cleared
rank-one identity (I.86) writes that entry, times univariate
nonzero pivots, as a product of univariate endpoint polynomials
and adjacent bivariate contractions. Part~1 preserves every
factor. Thus every internal interval has rank two at
$\mathbf X=\boldsymbol\alpha$, simultaneously.

\paragraph{4. Compatible constants and nonzero active endpoints.}
Work over $K=\overline E(t,z,y)$, and translate to independent
variables $Y_i=X_i-\alpha_i$. Absorb the first and last factors
into the affine endpoint row $u(X_1)=uA_1(X_1)$ and column
$v(X_n)=A_n(X_n)v$. The interior consists of positions
$2,\ldots,n-1$.

The constant row $u(\alpha_1)A_2(\alpha_2)$ is nonzero. Indeed,
the polynomial row $u(X_1)A_2(X_2)$ is nonzero, since otherwise
$f$ would be zero, and one of its nonzero bivariate entries
survives by part~1. Similarly
$A_{n-1}(\alpha_{n-1})v(\alpha_n)\ne0$. Because every internal
constant interval has rank two, multiplication by the remaining
constant interior factors is injective on these two-dimensional
active spaces. In particular both endpoint projections to the
persistent active space are nonzero.

Here is an explicit justification for the compatible gauges.
View the interior constants $C_i=A_i(\alpha_i)$ as linear maps
from the state space at cut $i$ to the state space at cut $i-1$.
At each internal cut, the image of the next constant factor is
two-dimensional and is complementary to the kernel of the
preceding constant factor, because the product of those two
factors has rank two. Use that image as the active space and
that kernel as its one-dimensional complement. Choose an
active complement to the kernel at the final cut, and transport
its basis leftwards through the constants; the full interval
rank condition makes all transported active bases invertible.
At the initial cut choose any complementary line. In these
bases every interior constant is
\[
 P=\operatorname{diag}(1,1,0),\qquad
 A_i(\alpha_i+Y_i)=P+Y_iD_i.
 \tag{I.95}
\]
The nonzero active endpoint row $a$ and column $b$ satisfy
$ab=f(\boldsymbol\alpha)=0$. Choose an active basis with
$a=e_1^{\mathsf T}$ and $b=e_2$; this is possible by taking
the second active basis vector to be $b$ and choosing the first
to have pairing one with $a$. Use this same active basis change
at every cut, keeping all copies of $P$ fixed.

The passive endpoint components can also be removed while
leaving those constants fixed. At the initial cut use
$L=\left(\begin{smallmatrix}I_2&h\\0&1\end{smallmatrix}\right)$:
$L^{-1}P=P$ and the row changes from $(a,\rho)$ to
$(a,ah+\rho)$. Choose $ah=-\rho$. At the final cut use
$R=\left(\begin{smallmatrix}I_2&0\\k^{\mathsf T}&1\end{smallmatrix}\right)$:
$PR=P$ and the column becomes
$(b,\sigma-k^{\mathsf T}b)^{\mathsf T}$. Choose
$k^{\mathsf T}b=\sigma$. These independent endpoint changes
also work when there is only one interior factor. They preserve
all matrix ranks and affine dependence. We have obtained
\[
 f(\boldsymbol\alpha+\mathbf Y)
 =(e_1^{\mathsf T}+Y_1p)
   \prod_{i=2}^{n-1}(P+Y_iD_i)
   (e_2+Y_nq).
 \tag{I.96}
\]
Here $p$ is a row and $q$ a column. All gauges are constant with
respect to the residual variables $\mathbf Y$; their dependence
on the center parameters is allowed.

Since every affine interior pencil has determinant identically
zero, the coefficient of $Y_i$ in
$\det(P+Y_iD_i)$ gives $(D_i)_{33}=0$. The first-derivative
identities from part~1 give
\[
 p_2=0,\qquad q_1=0,\qquad (D_i)_{12}=0
 \quad(2\le i\le n-1).
 \tag{I.97}
\]
Thus each active two-by-two derivative block is lower triangular,
and the passive state has no self-loop in either $P$ or $D_i$.

\paragraph{5. A tridiagonal Hessian forces every output path to vanish.}
If $|i-j|>1$, at least one constant factor $P$ separates the
two differentiated positions in (I.96). It kills the passive
coordinate. The active derivative blocks are lower triangular,
the active part of $p$ is supported on state one, and the active
part of $q$ on state two. Therefore $f_{ij}(\boldsymbol\alpha)=0$
for $|i-j|>1$. The diagonal is zero by multilinearity.

Consequently $H(\boldsymbol\alpha)$ is symmetric tridiagonal
with zero diagonal. Equation (I.93), and $c_i\ne0$ for every
$i$, force all its entries to be zero. The first row gives
$H_{12}=0$; the second row then gives $H_{23}=0$; continuing
along the path proves $H_{i,i+1}=0$ for all $i$. This argument
uses no division by two.

The adjacent entries have particularly simple formulas:
\[
 \begin{aligned}
 H_{12}&=p_3(D_2)_{32},\\
 H_{i,i+1}&=(D_i)_{13}(D_{i+1})_{32}
                 &&(2\le i\le n-2),\\
 H_{n-1,n}&=(D_{n-1})_{13}q_3.
 \end{aligned}
 \tag{I.98}
\]
All alternative active contributions are zero by (I.97).

Expand (I.96) as a sum over state paths. An active path cannot
move directly from state one to state two, because every
$(1,2)$ entry is zero. The passive state three has no self-loop.
Hence any path that first reaches state two after starting in
state one must contain a transition $1\to3$ immediately
followed by $3\to2$. Its two entries have product equal to
one of the middle expressions in (I.98), times the two residual
variables, and the product is zero. If the path starts in the
passive endpoint component, its first possible move into state
two is covered by the first expression in (I.98); if it first
returns to state one, the preceding argument applies later.
Finally, an output through the passive component of the final
column is covered by the last expression in (I.98), unless
state two was already reached, in which case an earlier zero
transition has already occurred. A path staying in state one
cannot exit through the active final column because $q_1=0$.
Thus every contributing state path contains a zero factor.
The polynomial (I.96) is identically zero over $K$, contrary
to the nonzero input. This proves $f(\mathbf R_2)\ne0$.

\paragraph{6. Conversion to ordinary centered SV and exact sampling.}
Put $A(T)=\prod_i(T-a_i)$ and $d_i=A'(a_i)\ne0$. For the
diagonally rescaled polynomial
$F(W_1,\ldots,W_n)=f(W_1/d_1,\ldots,W_n/d_n)$, the affine
singular-factor promise is unchanged. Take $c_i=d_i$ in (I.91).
The Lagrange identity
\[
 d_i\left(t+\sum_{h=1}^{2}z_h\lambda_i(y_h)\right)
 =d_i t-\sum_{h=1}^{2}\frac{z_h A(y_h)}{a_i-y_h}
 \tag{I.99}
\]
shows that the ordinary centered map is the rational map just
proved, after the invertible change of amplitudes
$z_h\mapsto-z_hA(y_h)$ in the rational parameter field.
Therefore its polynomial substitution is nonzero, proving (I.90).

For total degree at most $D$ (here $D\le n$), that substitution
has degree at most $D$ in $t$, at most $(n-1)D$ in each selector,
and total degree at most $D$ in the two amplitudes. The Newton
simplex interpolation argument already proved in I-21 yields
a hitting set of at most
\[
 (D+1)\bigl((n-1)D+1\bigr)^2\binom{D+2}{2}.
 \tag{I.100}
\]
The working field must supply at least
$\max\{n,D+1,(n-1)D+1\}$ distinct elements. A proper working
extension uses the same extension-evaluation convention as
I-21 and I-36. The proof holds in all characteristics, and
$D=0$ requires just one evaluation. The selected order and
all coefficient-dependent factorizations occur only in the
proof, not in the sampler.\hfill$\square$

\paragraph{Provenance and the remaining obstruction.}
This deduction uses the pairwise dominance and rank-one compound
factorization proved in I-36. The first-order pole-extraction
observation was supplied independently in this run by the purity
agent while investigating the Hu--van Melkebeek--Morgan
vanishing-ideal framework; the full argument above proves the
needed arbitrary-slope and arbitrary-characteristic version
directly. No historical priority claim is made.

The new closing step is specific to affine singular pencils:
after rank preservation their passive constant entry and passive
linear coefficient both vanish. A variable-dependent invertible
factor can propagate the passive state, so nonadjacent Hessian
entries need not vanish and (I.93) need not force all adjacent
transition products to vanish. Higher-degree singular factors
also need not have a zero passive entry at every positive degree.
Thus the mixed rank-two/rank-three centered-two-pair gate remains
open, as does an unrestricted polynomial-in-width construction.

\subsection{I-38. Two centered pairs with isolated invertible affine factors}
\label{I-38}

\paragraph{Statement.}
The conclusion and sample bound of I-37 also hold for a width-at-most-three
affine univariate matrix word with the following larger promise:
after constant factors have been absorbed, no two consecutive
variable-dependent factors both have generic rank three. Thus
arbitrarily many generic-rank-three affine factors may occur,
provided a singular factor separates each two of them. The order
remains unknown, the nodes are any distinct working-field elements,
and the conclusion holds in every characteristic.

Equivalently, for any nonzero slopes $c_i$ and distinct nodes $a_i$,
the rational map $R_{2,i}=tc_i+z_1/(a_i-y_1)+z_2/(a_i-y_2)$
hits the stated class. The diagonal rescaling (I.99) then proves
ordinary common-center $t\mathbf1+\mathrm{SV}_2$ nonvanishing.
We prove the rational formulation.

\paragraph{1. The one-pole center dominates every three coordinates.}
Let $\alpha_i=tc_i+z/(a_i-y)$. Pairwise dominance is (I.92).
For distinct $i,j,k$, solve for $t,z$ from $X_i,X_j$ over $E(y)$.
The third coordinate has the form
\[
 X_k=A(y)X_i+B(y)X_j,
 \qquad A,B\in E(y).
 \tag{I.101}
\]
The two coefficients cannot both be constant. If they were,
comparison of the $z$ coefficients would give
$1/(a_k-y)=A/(a_i-y)+B/(a_j-y)$, which the distinct poles
exclude. Hence the right side of (I.101) is a nonconstant
rational function of $y$ over $E(X_i,X_j)$, and is transcendental
over that field. Thus the three coordinates are algebraically
independent. This proves dominance in every characteristic,
without a Jacobian argument. It also proves the claim over every
fixed extension of the coefficient field.

\paragraph{2. The representation promise survives irreducible reduction.}
Assume a nonzero annihilator and apply I-37's prime-kernel and
disjoint-variable factor argument. Specializing the variables
outside an irreducible factor makes the intervening matrix
factors constant. Absorb each constant block into the next
remaining factor, and absorb the final block into the terminal
column. All remaining factors are affine and their ranks cannot
increase.

The no-adjacent-rank-three promise also persists. If two adjacent
remaining factors both had rank three, both corresponding
original factors would have had rank three. The original promise
puts a singular original factor between them. It has now been
specialized into the intervening constant block, so that block
has rank at most two. Absorbing it into the later factor makes
that factor singular, a contradiction. This reasoning does not
require any specialized singular factor to keep its original
rank. A zero specialization cannot occur in the realized nonzero
scalar word.

We may therefore assume that the input is irreducible and depends
on every remaining variable. As in I-37, a strict internal
rank-one interval has a polynomial outer factorization and
splits the scalar into two nonconstant factors. Every strict
internal interval consequently has generic rank two or three.
If there are at most three remaining variables, part~1 already
hits the scalar. Assume $n\ge4$ below.

\paragraph{3. One pole preserves every internal interval rank.}
Each generic-rank-three factor is invertible at its univariate
center, because its nonzero determinant is univariate and the
center is transcendental. Its second compound is likewise
invertible there.

Consider a generic-rank-two interval. In its second compound
word, a singular original factor becomes a rank-one factor,
and every intervening rank-three factor becomes invertible.
There is at most one such intervening factor between consecutive
rank-one factors. Write a nonzero rank-one factor as $p_iq_i/b_i$
using a nonzero univariate pivot. The connecting scalar between
two successive rank-one factors is therefore either
\[
 q_i(X_i)p_j(X_j)
 \quad\hbox{or}\quad
 q_i(X_i)C_2(A_k(X_k))p_j(X_j).
 \tag{I.102}
\]
It depends on at most three distinct coordinates. Since the
original compound interval product is nonzero, every connecting
scalar is nonzero. Part~1 preserves all these scalars and all
univariate pivots. The remaining rank-one endpoint vectors are
nonzero at their centers, and any full-rank prefix or suffix is
invertible there. The compound interval product stays nonzero,
so the original interval still has rank two. A generic-rank-three
interval consists entirely of invertible factors and retains
rank three. Thus all strict internal interval ranks are preserved
simultaneously. The proof requires preservation of their ranks;
it does not assert preservation of every individual two-minor.

Absorb $A_1,A_n$ into the affine endpoint row and column. There
is an internal rank-two factor: otherwise the at-least-two
internal positions would be consecutive rank-three factors.
The first internal rank-two factor is at position two or three.
The endpoint row times the prefix through that factor is a
nonzero polynomial row in at most three variables, and remains
nonzero at the center by part~1. All subsequent constant
intervals have rank at least two, so propagation of its active
two-dimensional row space is injective. Hence the left endpoint
has a nonzero active projection for the full constant interior
word. The reversed argument at the last internal rank-two
factor gives a nonzero active projection of the right endpoint.

\paragraph{4. Compatible $I/P$ constants and endpoint shears.}
Over the field $K=\overline E(t,z,y)$, the constant interior word
has full-product rank two, and every subinterval rank two or
three. We give the compatible-space argument explicitly.
Use row-vector orientation. Choose a two-dimensional initial
space $U_0$ complementary to the left kernel of the full
constant product, and propagate it through the constants to
spaces $U_i$. Each propagation is an isomorphism on $U_i$.
Choose a complementary line $V_m$ at the final cut. For the
suffix product $T_i$ from cut $i$ to the end, put
\[
 V_i=\{v:vT_i\in V_m\}.
\]
If $T_i$ has rank two, its row image equals $U_m$, so
$V_i=\ker T_i$ is a line. If it has rank three, the inverse
image of $V_m$ is a line. In either case $U_i\cap V_i=0$,
and each constant factor sends the corresponding passive line
into the next one. Active bases transported through the word
make every active block $I_2$. The passive block is zero at
rank-two factors and a nonzero scalar at rank-three factors;
rescale the successive passive bases to make the latter scalar
one. Thus all interior factors have the form
\[
 A_i(\alpha_i+Y_i)=C_i+Y_iD_i,
 \qquad C_i\in\{P,I_3\},\quad P=\operatorname{diag}(I_2,0),
 \tag{I.103}
\]
with $C_i=P$ precisely at the rank-two positions. No consecutive
interior constants are both $I_3$.

The nonzero active endpoint vectors pair to
$f(\boldsymbol\alpha)=0$. Normalize them to $e_1^{\mathsf T}$
and $e_2$ by one common active change of basis. The endpoint
passive components must also be removed. Use the upper shear
$L$ from I-37 at every cut of the initial string of $I_3$
constants, up to the first $P$. Then
$L^{-1}I_3L=I_3$ along that string and $L^{-1}P=P$ at its end,
while the initial passive row component is killed. Similarly
use the lower shear $R$ from I-37 at every cut of the terminal
string of $I_3$ constants after the last $P$:
$PR=P$ and $R^{-1}I_3R=I_3$. This kills the passive column
component. If the first and last $P$ are the same factor,
$L^{-1}PR=P$ permits both shears simultaneously. Thus all
constants in (I.103) remain as stated, and
\[
 f(\boldsymbol\alpha+\mathbf Y)
 =(e_1^{\mathsf T}+Y_1p)
   \prod_{i=2}^{n-1}(C_i+Y_iD_i)
   (e_2+Y_nq).
 \tag{I.104}
\]

The pole extraction of I-37 gives $f_i(\boldsymbol\alpha)=0$
for every $i$. Therefore $p_2=q_1=(D_i)_{12}=0$. At every
singular factor, the coefficient of $Y_i$ in
$\det(P+Y_iD_i)=0$ also gives $(D_i)_{33}=0$. An invertible
factor may have a nonzero passive derivative entry; it is not
silently discarded.

\paragraph{5. Two tangent directions kill the bandwidth-two Hessian.}
Set $h_i=1/(a_i-y)$. Differentiate the identities
$f_i(\boldsymbol\alpha)=0$ with respect to $z$ and $y$.
Equivalently, extract their first coefficients under independent
parameter increments. Since $z\ne0$ in the rational function
field, this gives in every characteristic
\[
 H(\boldsymbol\alpha)\mathbf h=0,
 \qquad H(\boldsymbol\alpha)\mathbf h^{\circ2}=0,
 \tag{I.105}
\]
where $\mathbf h^{\circ2}=(h_1^2,\ldots,h_n^2)$. The diagonal
of $H$ is zero by multilinearity.

If $i<j$ and a singular interior position lies strictly between
them, the intervening constants contain $P$. Their product
is $P$, which removes the passive coordinate, while the active
derivative blocks are lower triangular. Hence $H_{ij}=0$.
If there is no such position, all intervening constants are
$I_3$, and $H_{ij}$ is exactly the product of the state-one-to-
passive entry at position $i$ and the passive-to-state-two
entry at position $j$, with the endpoint versions $p_3$ and
$q_3$ when appropriate. There are no alternative active
contributions because all $(1,2)$ entries vanish.

Since no two consecutive internal positions are invertible,
every gap of length at least three contains a singular
intermediate position. Thus $H_{ij}=0$ whenever $|i-j|>2$.
In the first row, at most two future entries remain. For two
indices $j\ne k$, the corresponding coefficient determinant
in (I.105) is
\[
 \det\begin{pmatrix}h_j&h_k\\h_j^2&h_k^2\end{pmatrix}
 =h_jh_k(h_k-h_j)\ne0.
 \tag{I.106}
\]
Both entries are therefore zero. If only one entry remains,
$h_j\ne0$ suffices. By symmetry the earlier entries in the
next row are now zero; apply the same argument to its at most
two future entries. Induction gives $H=0$.

Finally expand (I.104) into state paths, as in I-37. A path
cannot move directly from active state one to state two.
Before its first arrival at state two, take its last departure
from state one to the passive state, or its initial passive
endpoint entry if it has not yet visited state one. Every
intermediate passive stay must occur at an invertible factor:
a singular factor has both passive constant and passive
derivative entries zero. The product of the departure and
arrival entries is the corresponding Hessian entry $H_{ij}$,
and is zero. Intermediate passive stays only multiply this
zero product. The same argument covers an output through the
passive terminal component; an output through active state one
is excluded by $q_1=0$. Therefore every output path is zero,
contradicting the assumed nonzero scalar. This proves the
rational two-pair assertion and hence the ordinary centered
two-pair assertion.\hfill$\square$

\paragraph{Sampling, provenance and the remaining gate.}
The exact field and interpolation requirements are unchanged
from I-37: for total degree at most $D\le n$, the sample count
is (I.100), over a working field with at least
$\max\{n,D+1,(n-1)D+1\}$ distinct elements, or an extension
with the explicitly stated evaluation convention. No input
order or representation is reconstructed by the sampler.

The extra proof ingredients are the function-field proof of
triple dominance, the preserved adjacency promise under
irreducible specialization, the rank-one compound connectors
through one invertible factor, and the two tangent kernel
identities. Purity independently suggested the use of
$\mathbf h,\mathbf h^{\circ2}$ instead of slope-dependent
minors, avoiding an invalid Jacobian shortcut in characteristic
two. No historical priority claim is made.

Two or more consecutive invertible factors remain outside
this proof. Their compound connectors involve at least four
coordinates, so three-coordinate dominance alone does not
preserve their rank; even at a rank-preserved center, the
Hessian can have more than two future entries in a row. The
larger semiseparable Hessian structure and the consequences
of the assumed two-pair identity for possible one-pole rank
loss remain the specific next gate.

\subsection{I-39. Two exact failed repairs for unrestricted mixed affine words}
\label{I-39}

\paragraph{Scope.}
We prove two obstructions to the next proposed extensions of I-38.
First, the ordinary common-center one-pair map can lower a strict
internal interval from generic rank two to rank one even in an
irreducible multilinear width-three scalar word whose every generic
interval rank is two or three. Second, the symmetric semiseparable
Hessian pattern, all three tangent kernel identities, and the whole
additional-pole quadratic identity do not by themselves force the
Hessian to be zero. Neither construction is a counterexample to
the centered two-pair sampler. The first scalar is hit already
by one pair; the second is an abstract family of quadratic jets,
with a stated failure of integration to a fixed polynomial.
Both constructions are over characteristic zero.

\paragraph{1. A rational identity and its residue forms.}
We recall, and verify the needed parts of, the explicit construction
of P-16. Put $a=\sqrt{3/2}$ and, for an integer $d\ge2$, set
\[
 \begin{aligned}
 P_1(T)&=T^{2d}-1,&P_2(T)&=T^d-a,\\
 P_3(T)&=-2(T^d+a),&P_4(T)&=T^{2d}-2,\\
 R_1&=P_2/P_1,&R_2&=1/P_2,\\
 R_3&=1/P_3,&R_4&=P_3/P_4.
 \end{aligned}
 \tag{I.107}
\]
Work in a splitting field of these polynomials. Their four pole
sets are disjoint and simple, of sizes $2d,d,d,2d$ respectively;
none of the displayed numerators cancels a pole. For a pole
$b$ of $R_j$, let $\rho_b\ne0$ be its residue. Partial fractions
and decay $R_j(T)=O(T^{-d})$ give
\[
 R_j(T)=\sum_{b\in\mathcal A_j}\frac{\rho_b}{T-b},
 \qquad \sum_{b\in\mathcal A_j}\rho_b b^e=0
                    \quad(0\le e\le d-2).
 \tag{I.108}
\]
Moreover $P_2P_3=-P_1-P_4$, so
$R_1R_2+R_1R_4+R_3R_4=0$.
For the disjointly supported nonzero linear forms
$L_j=\sum_{b\in\mathcal A_j}\rho_b X_b$, define
\[
 Q=L_1L_2+L_1L_4+L_3L_4\ne0.
 \tag{I.109}
\]
Nonvanishing follows, for example, from any coefficient using
one variable from group one and one from group two. If $g$ has
degree at most $d-1$, division by $T-u$ and (I.108) give
\[
 L_j\left(\frac{g(b)}{b-u}:b\in\bigcup\mathcal A_j\right)
 =-g(u)R_j(u).
 \tag{I.110}
\]
Therefore $Q$ vanishes on all these numerator-at-most-$(d-1)$,
one-pole evaluations. These are the only P-16 facts used below.

\paragraph{2. Ordinary one-pair rank loss on an irreducible affine word.}
Take $d=2$, giving twelve core nodes. They come in opposite pairs,
and every $R_j$ is an even rational function. Add the five nodes
$0,3,-3,4,-4$, giving seventeen distinct nodes in total. Assign
the four nonzero new nodes to variables $s,u,v,t$, respectively;
the coordinate at zero is unused by the input polynomial.
Let
\[
 A(T)=\prod_{b\text{ among all seventeen nodes}}(T-b),
 \qquad d_b=A'(b),\qquad
 L_j^*=\sum_{b\in\mathcal A_j}\rho_b d_bX_b.
 \tag{I.111}
\]
The polynomial $A$ is odd and $A'$ is even. The residues of an
even rational function at opposite poles are negatives of one
another. It follows that $\sum_{b\in\mathcal A_j}\rho_b d_b=0$
for every group. Hence under the ordinary common-center map
$X_b=\tau+z\lambda_b(y)$ on these seventeen nodes,
\[
 L_j^*=zA(y)R_j(y),\qquad
 Q^*=L_1^*L_2^*+L_1^*L_4^*+L_3^*L_4^*=0.
 \tag{I.112}
\]
Indeed $d_b\lambda_b(y)=A(y)/(y-b)$, so this is an exact
polynomial identity after clearing the displayed rational notation.
All coefficients $\rho_bd_b$ are nonzero, and $Q^*$ is a
nonzero quadratic before substitution.

Put $A_0=L_1^*+L_3^*$, $B_0=L_2^*+L_4^*$ and $C_0=L_2^*L_3^*$.
The middle word consists of the affine unipotent factors
$I-\rho_bd_bX_bE_{21}$ in groups one and three, and
$I-\rho_bd_bX_bE_{32}$ in groups two and four, in group order.
Within each group their product is $I-L_j^*E$, since $E^2=0$.
Thus its full product is exactly
\[
 M=I-A_0E_{21}-B_0E_{32}+C_0E_{31}.
 \tag{I.113}
\]
Let $P_L=\operatorname{diag}(0,1,1)$ and
$P_R=\operatorname{diag}(1,1,0)$. Their corner is
\[
 P_LMP_R=
 \begin{pmatrix}0&0&0\\-A_0&1&0\\C_0&-B_0&0\end{pmatrix},
 \qquad
 \det\begin{pmatrix}-A_0&1\\C_0&-B_0\end{pmatrix}=Q^*.
 \tag{I.114}
\]
Use genuinely variable-dependent singular anchor factors
\[
 J_L(u)=(I+uE_{32})P_L,\qquad
 J_R(v)=P_R(I+vE_{21}).
\]
They have rank two at every value. The strict internal interval
$J_L(u)M J_R(v)$ has generic rank two because $Q^*\ne0$.
It has rank exactly one at the center (I.112): its indicated
two-by-two determinant is zero, while its $(2,2)$ entry remains
one. The outer unipotent multipliers in the anchors preserve
this rank calculation.

Now take the actual scalar word
\[
 F=e_3^{\mathsf T}(I+sE_{32})J_L(u)M J_R(v)
                      (I+tE_{21})e_1
   =(s+u-B_0)(t+v-A_0)-Q^*.
 \tag{I.115}
\]
Every variable is read once; every matrix factor is affine.
The two anchors are the only singular factors. An interval
containing no anchor is invertible, one containing exactly one
anchor has rank two, and one containing both has rank two
because (I.114) has nonzero determinant. Thus all generic
interval ranks are in $\{2,3\}$.

The scalar $F$ is irreducible. The invertible polynomial change
of variables $S=s+u-B_0$, $T=t+v-A_0$, keeping the remaining
variables fixed, converts it to $ST-Q^*$. As a polynomial in
$T$ over the UFD of the other variables, its coefficients $S$
and $-Q^*$ are coprime. It is primitive and linear, hence
irreducible over that UFD by Gauss's lemma. This also proves
irreducibility after extending to the algebraic closure.

Nevertheless $F$ is detected by the very center that loses
the internal rank. By (I.112), its image is the product of
the two factors in (I.115). Each has coefficient two on
$\tau$, while $A_0,B_0$ have no $\tau$ term. The image has
coefficient four on $\tau^2$ and is nonzero. This demonstrates
that irreducibility alone does not repair one-pole interval-rank
preservation, and that requiring preservation of every interval
rank is stronger than the desired scalar nonvanishing.

\paragraph{3. The longer-block Hessian conditions alone also fail.}
Take $d=4$ in part~1, with its twenty-four core nodes, and let
$H_0$ be the symmetric zero-diagonal Hessian of $Q$ in (I.109).
In the group order, define vectors $\alpha,\beta$ by putting
$\alpha_b=\rho_b$ in groups one and three and zero otherwise,
and $\beta_b=\rho_b$ in groups two and four and zero otherwise.
For $i<j$ the upper entries satisfy
$(H_0)_{ij}=\alpha_i\beta_j$. This is exactly the single-block
upper rank-one semiseparable pattern arising from an arbitrary
invertible stretch in the I-38 Hessian calculation.

Let $h_i(y)=1/(a_i-y)$, where $a_i$ now lists the twenty-four
core nodes, and set
\[
 D(y)=\operatorname{diag}((a_i-y)^2),\qquad
 H(y)=D(y)H_0D(y).
 \tag{I.116}
\]
It is nonzero, symmetric and zero-diagonal, and its upper
entries still factor as
$[\alpha_i(a_i-y)^2][\beta_j(a_j-y)^2]$.
Every $L_j$ kills polynomial evaluations of degrees at most
two by (I.108). Therefore $H_0$ annihilates each such evaluation
vector, and
\[
 H(y)\mathbf1=H(y)\mathbf h(y)
              =H(y)\mathbf h(y)^{\circ2}=0.
 \tag{I.117}
\]
The three vectors entering $H_0$ here are respectively
$((a_i-y)^2)_i$, $((a_i-y))_i$, and $\mathbf1$.

It also satisfies the complete extra-pole quadratic identity:
for an independent pole $u$,
\[
 \mathbf h(u)^{\mathsf T}H(y)\mathbf h(u)
 =2Q\left(\frac{(a_i-y)^2}{a_i-u}:1\le i\le24\right)=0.
 \tag{I.118}
\]
This is (I.110) with numerator $(T-y)^2$ of degree two.
Consequently the tangent kernels and the additional-pole
quadratic restriction do not suffice, even with the exact
semiseparable pattern, to force $H=0$.

This family is not itself the Hessian pullback of a fixed
polynomial at $\tau\mathbf1+z\mathbf h(y)$. If it were,
specializing $z=0$ would make that pullback independent of $y$,
whereas every nonzero entry of (I.116) is a nonconstant
degree-four polynomial in $y$. Thus the construction explicitly
exhibits the missing coordinate-ring and Hessian-integrability
constraints; it does not silently assert they have been met.
Multiplying or clearing denominators to enter the center's
coordinate ring would still require checking those constraints,
and the original width-three affine representation, anew.
\hfill$\square$

\paragraph{Provenance and next gate.}
The rational identity and residue forms are from the fully
checked P-16 construction; all extra matrix, parity, irreducibility
and Hessian deductions needed here are written out. No historical
priority claim is made. The remaining two-pair target must use
the assumed scalar annihilation, rather than irreducibility
alone, to control rank loss; or retain the integrability and
affine compound constraints beyond the displayed Hessian
identities. Neither obstruction settles that target or the
growing-width polynomial-size problem.


\subsection{I-40. A determinant-weight budget for translation-invariant matrix words}
\label{I-40}

\paragraph{Statement.}
Let $E$ be any field and let $0\ne f\in E[X_1,\ldots,X_n]$ be
homogeneous and multilinear of degree $D$, where $n\ge1$.
Assume the formal common-translation identity
\[
 f(X_1+s,\ldots,X_n+s)=f(X_1,\ldots,X_n).
 \tag{I.119}
\]
Suppose that the coefficient-cut ranks of $f$ in some order are
at most $w$. Put $m=n-2D$. Then $m\ge0$.
There is a minimal affine univariate matrix representation of
$f$ in that order, with bond sizes equal to its coefficient-cut
ranks, having the following property: every consecutive run of
generic-rank-$w$ factors in its consistent $w\times w$ padding
has length at most $m$. The same bound survives absorption of
constant factors. This assertion is valid in every characteristic;
the translation identity is a polynomial identity, not merely an
identity on the rational points of a finite field.

For $w\le3$ and $m\le1$, I-37 and I-38 consequently give
\[
 f\bigl(t\mathbf1+\mathrm{SV}_2(\mathbf y,\mathbf z)\bigr)\ne0,
 \qquad
 f\left(\left(tc_i+\frac{z_1}{a_i-y_1}
                         +\frac{z_2}{a_i-y_2}\right)_{i=1}^n\right)\ne0
 \tag{I.120}
\]
for arbitrary distinct nodes $a_i$ and arbitrary nonzero slopes
$c_i$ in the working field. No order is supplied to either map.
For this corollary use padding to width three even if the actual
width is smaller; if no rank-three core occurs, I-37 applies.

Together with P-17, the result also proves that balanced rational
two-pair evaluation hits every nonzero multilinear polynomial
of ROABP width at most three on at most nine variables, in every
characteristic and with unknown order. This last assertion is
about the balanced rational map. The proof does not identify it
with ordinary unweighted centered $\mathrm{SV}_2$ for every
nine-variable input.

\paragraph{1. The established highest-weight input and conventions.}
Write $V=E e_0\oplus E e_1$ and use the standard tensor action of
$G=\mathrm{SL}_2$. Let
\[
 u_+(s)=\begin{pmatrix}1&s\\0&1\end{pmatrix},
 \qquad h(t)=\operatorname{diag}(t,t^{-1}).
\]
We use the standard characteristic-free fact that a nonzero
$u_+$-fixed vector in a rational $G$-module, on which $h(t)$ acts
as $t^d$, has $d\ge0$. Jantzen (n.d.), Lecture~I, p.~2
(Frobenius reciprocity) and p.~3 (the explicit
$\mathrm{SL}_2$ induction computation), supply this exact input
in every characteristic; see I-B09. Upper- versus lower-Borel
conventions are changed by dualizing and conjugating by the Weyl
element. The split-torus weight decomposition is the established
integer-character grading of G-B02. These facts concern rational
group representations, and do not invoke complete reducibility
or Lie algebra eigenvalues. The convention check and an
elementary Bruhat verification are retained in the separate
\texttt{I40\_highest\_weight\_audit.md}.

Multihomogenize $f$ as the tensor
\[
 F=\sum_{S\subseteq[n]}c_S
       \bigotimes_{i=1}^n e_{\mathbf1_{i\in S}}
   \ \equiv\ \left(\prod_i U_i\right)
                f(V_1/U_1,\ldots,V_n/U_n).
 \tag{I.121}
\]
The upper unipotent fixes $e_0$ and sends $e_1$ to $e_1+s e_0$.
Thus (I.119) says exactly that $u_+(s)F=F$. Homogeneity gives
$h(t)F=t^{n-2D}F=t^mF$. The cited highest-weight fact gives
$m\ge0$. Scalar extension to an algebraic closure is harmless
here: it preserves polynomial identities and every coefficient
rank. The support spaces and matrix cores below may all be
defined over the original field $E$.

\paragraph{2. Both sides of a coefficient cut bound its determinant weight.}
Let $L_i\subseteq V^{\otimes i}$ and
$R_i\subseteq V^{\otimes(n-i)}$ be the left and right support
spaces of $F$ at the cut after position $i$. Concretely, $L_i$
is the image of contraction with all linear functionals on the
right tensor factors, and $R_i$ is defined symmetrically.
Their common dimension $r_i$ is the coefficient-cut rank;
$r_0=r_n=1$. Since $F$ is a semi-invariant of the upper Borel,
both support spaces are stable under that Borel. Indeed applying
a Borel element to a contraction is, up to its nonzero scalar
on $F$, another contraction after applying the inverse action
to the functional. Applying the inverse element gives equality
of the resulting spaces.

Let the determinant characters of the torus on these spaces be
$t^{\delta_i}$ and $t^{\delta_i'}$, respectively. They are integer
characters. The upper unipotent has determinant one on either
space: its representation matrix $G(s)$ has entries in $E[s]$
and inverse $G(-s)$, so its determinant is a unit in $E[s]$;
its value at zero is one. Consequently the determinant line of
$L_i$ is an upper-unipotent-fixed line in the rational
$G$-module $\bigwedge^{r_i}(V^{\otimes i})$. The highest-weight
fact gives $\delta_i\ge0$. The same argument on the right gives
$\delta_i'\ge0$.

In bases of $L_i$ and $R_i$, the tensor $F$ has a nonsingular
$r_i\times r_i$ coefficient matrix $C$. Torus semi-invariance
therefore gives
\[
 G_L(t) C G_R(t)^{\mathsf T}=t^m C.
\]
Taking determinants yields the exact budget
\[
 \delta_i+\delta_i'=r_i m,
 \qquad 0\le\delta_i\le r_i m.
 \tag{I.122}
\]
This is a statement in the character lattice $\mathbb Z$.
In positive characteristic, neither the inequality nor the
equality is reduced modulo the characteristic.

\paragraph{3. Minimal affine cores and their exact covariance.}
There is a natural inclusion
\[
 \iota_i:L_i\longrightarrow L_{i-1}\otimes V.
\]
To verify it, write $F$ as an element of
$L_{i-1}\otimes V\otimes V^{\otimes(n-i)}$ and contract its last
$n-i$ factors. These inclusions are Borel-equivariant. Choose
bases of the $L_i$, with the basis of $L_0=E$ equal to $1$ and
that of $L_n=EF$ equal to $F$. Write
\[
 \iota_i(b_{i,b})
   =\sum_a b_{i-1,a}\otimes
       \bigl(A_i^0[a,b]e_0+A_i^1[a,b]e_1\bigr),
 \qquad A_i(X_i)=A_i^0+X_iA_i^1.
 \tag{I.123}
\]
Repeatedly applying the inclusions reconstructs every tensor
coefficient, so the scalar matrix product
$A_1(X_1)\cdots A_n(X_n)$ equals $f$. The matrix shape is
$r_{i-1}\times r_i$, and these sizes are minimal because they
are the actual coefficient-cut ranks. This construction uses
only support-space inclusions and linear algebra, over any
field.

Let $G_i(s)$ be the upper-unipotent action matrix on $L_i$.
Equivariance of (I.123), comparing the coefficients of $e_0$
and $e_1$, gives
\[
 A_i(X_i)G_i(s)=G_{i-1}(s)A_i(X_i+s).
 \tag{I.124}
\]
In particular this formula includes both endpoints, with
$G_0=G_n=1$. It is an identity of matrices over $E[X_i,s]$;
there are no coefficient-dependent rational denominators.

Suppose $r_{i-1}=r_i=r$ and $A_i$ has generic rank $r$.
The determinant-one property of $G_{i-1}(s)$ and $G_i(s)$
implies
\[
 \det A_i(X_i+s)=\det A_i(X_i).
\]
Setting $X_i=0$ proves that this nonzero determinant is a
nonzero constant. Thus $A_i^0$ is invertible. For the torus,
comparison of the $e_0$ coefficient of (I.123) gives
\[
 A_i^0G_i(t)=t\,G_{i-1}(t)A_i^0.
\]
Taking determinants and cancelling $\det A_i^0\ne0$ proves
\[
 \delta_i=\delta_{i-1}+r.
 \tag{I.125}
\]
This is the key local restriction: a square full-rank core
increases the determinant character by its entire bond size.

\paragraph{4. The run bound and the width-three corollary.}
Pad each core consistently to $w\times w$, embedding each
$E^{r_i}$ as the first $r_i$ coordinates of $E^w$. Use the first
coordinate row and column as the scalar endpoints. This
preserves the product exactly. A padded factor can have generic
rank $w$ only if $r_{i-1}=r_i=w$ and its unpadded core is full
rank. On a run of $\ell$ such factors, (I.125) telescopes to
\[
 w\ell=\delta_{\mathrm{end}}-\delta_{\mathrm{start}}
           \le wm,
 \qquad \ell\le m,
 \tag{I.126}
\]
by (I.122) at the two bounding cuts. This proves the asserted
budget.

To absorb constant factors, multiply each intervening constant
block into the following variable-dependent factor, and absorb
the terminal block into the final column. Ranks cannot increase.
If two or more remaining full-rank factors form a consecutive
run, every constant factor removed between them must itself
have had full rank; otherwise its product would make the later
factor singular. Thus those remaining factors belong to one
full-rank run before absorption. The length of the remaining
run is no larger. Initial constant blocks and the final boundary
cannot create an extra adjacent variable factor. The same run
bound therefore holds after absorption.

For width at most three, use padding to three. If $m=0$, every
factor is singular, so I-37 applies. If $m=1$, no adjacent
variable-dependent factors have generic rank three after
absorption, so I-38 applies. This proves (I.120), with its exact
node and slope quantifiers. Its ordinary polynomial sampler
has the already proved size
\[
 (D+1)\bigl((n-1)D+1\bigr)^2\binom{D+2}{2}
 \tag{I.127}
\]
over a working field with at least
$\max\{n,D+1,(n-1)D+1\}$ elements. For $D=0$ one evaluation
suffices. A proper extension uses the explicit extension-input
convention of I-21 and I-37; the proof does not assume access to
extension evaluations when the black box only accepts base-field
points.

\paragraph{5. The nine-variable balanced two-pair theorem.}
Fix distinct nodes and let $\mathcal I_2$ denote the vanishing
ideal of balanced rational evaluation with numerator and
denominator degrees at most two, as in P-17. P-17 has already
proved the entire multilinear case on at most eight variables
and the homogeneous cubic case on at most ten variables, in
every order and characteristic. We use that stable result
without changing its statement.

Suppose that a nonzero multilinear $f\in\mathcal I_2$ on nine
variables has width at most three in some order. The established
homogeneous-degree decomposition of this ideal and the
lowest-component closure for ROABP width (P-B20, Theorem~9.13
and Proposition~8.3) give a nonzero homogeneous annihilator of
width at most three, of degree between three and six. The cubic
case of P-17 excludes degree three. Multilinear reversal
preserves both the balanced ideal and all coefficient-cut ranks,
so it excludes degree six and reduces degree five to degree
four. We may therefore suppose that $f$ is a homogeneous
quartic.

Balanced rational evaluation is closed under adding a constant.
Thus every coefficient of $f(X+s\mathbf1)$ in $s$ lies in
$\mathcal I_2$. Each coefficient of $s^j$ has degree $4-j$,
and this ideal has no nonzero multilinear component of degree
at most two. Hence
\[
 f(X+s\mathbf1)=f(X)+s g(X),
 \qquad g\in\mathcal I_2\text{ homogeneous cubic}.
 \tag{I.128}
\]
If $g\ne0$, it has the same cut-rank upper bound as $f$.
Indeed common translation is an invertible product of local
linear changes on the coefficient tensor, preserving cut ranks
over $E(s)$. With $q=s^{-1}$, the coefficient matrix of
$s^{-1}f(X+s\mathbf1)=qf(X)+g(X)$ is regular at $q=0$.
Every four-by-four minor vanishes before specialization and
therefore at $q=0$. Thus every coefficient-cut matrix of $g$
has rank at most three. P-17's nine-variable cubic theorem
forces $g=0$.

We have proved the translation identity (I.119) for this
quartic. Its weight is $m=9-2\cdot4=1$. Part~4 therefore gives
nonzero evaluation by the rational map with all slopes $c_i=1$.
That map is balanced of degrees at most two, contradicting
$f\in\mathcal I_2$. This proves the full nine-variable result.
All coefficient extractions and rank specializations in this
argument are valid in every characteristic.

\paragraph{6. An explicit finite balanced sampler.}
For any multilinear polynomial $f$ of total degree at most $D$,
write
\[
 R_i=t+\frac{z_1}{a_i-y_1}+\frac{z_2}{a_i-y_2},
 \qquad
 Q(\mathbf y)=\prod_{i=1}^n(a_i-y_1)(a_i-y_2).
 \tag{I.129}
\]
The rational map $\mathbf R$ has the same vanishing ideal as
balanced degree-two rational evaluation. One way to check the
scope is to pass to an algebraic closure: a generic quadratic
denominator has two distinct roots, and partial fractions give
exactly a constant plus two simple poles. The denominators with
distinct roots form a dense open subset in every characteristic.
Thus the parametrizations have the same formal polynomial
identities; no assertion about every rational function splitting
over the original finite field is needed.

For a nonzero pullback, $P=Qf(\mathbf R)$ is a nonzero polynomial.
Multilinearity is used here: expanding one monomial of $f$,
each index contributes either $(a_i-y_1)(a_i-y_2)$ or
\[
 t(a_i-y_1)(a_i-y_2)+z_1(a_i-y_2)+z_2(a_i-y_1).
\]
Consequently $P$ has degree at most $n$ in each selector $y_h$
and total degree at most $D$ in $(t,z_1,z_2)$. Choose $n+1$
distinct selector values, all outside $\{a_1,\ldots,a_n\}$,
and $D+1$ distinct amplitude values $b_0,\ldots,b_D$.
Sample the three amplitudes at the Newton simplex
\[
 (t,z_1,z_2)=(b_j,b_k,b_\ell),\qquad
               j,k,\ell\ge0,\quad j+k+\ell\le D.
\]
Univariate interpolation in the two selectors and the lower-set
Newton interpolation in the three amplitudes detect every
nonzero polynomial of these degrees. Every selected denominator
is nonzero. Thus the nine-variable theorem has a finite hitting
set of at most
\[
 (n+1)^2\binom{D+3}{3}
 \tag{I.130}
\]
points, for $n\le9$, over a working field with at least
$\max\{2n+1,D+1\}$ elements. The nodes and the sets of sample
values are fixed in advance and are independent of the hidden
order and all input coefficients. Extension evaluations, if
needed, obey the same explicitly stated black-box convention
as above.\hfill$\square$

\paragraph{Scope and the next gate.}
The determinant budget is not a polynomial-in-width PIT
construction. It controls runs of square full-width cores of a
homogeneous translation-invariant tensor; it does not constrain
every lower-rank transition or arbitrary nonhomogeneous words.
At weight two it permits two consecutive invertible cores, so
I-38 alone no longer closes the scalar question. The
translation-invariant ten-variable quartic is a concrete next
case, and the unrestricted mixed width-three problem and the
growing-width isolation transfer remain open here.

The highest-weight input is established representation theory.
The determinant-character budget, its core covariance
application, and the nine-variable deduction are proved above
as deductions in this research record, without a historical
priority claim. The standard dominance calculation and the
independent all-characteristic checks are recorded separately.


\subsection{I-41. Weight-two connectors and the balanced ten-variable theorem}
\label{I-41}

\paragraph{Statement and scope.}
Let $E$ be any field, and fix arbitrary distinct nodes
$a_1,\ldots,a_n$ in $E$. Every nonzero homogeneous multilinear
polynomial $f$ of width at most three in some variable order,
satisfying
\[
 f(X+s\mathbf1)=f(X),\qquad 0\le n-2\deg f\le2,
 \tag{I.131}
\]
has nonzero pullback under the balanced rational map
\[
 R_i=t+\frac{z_1}{a_i-y_1}+\frac{z_2}{a_i-y_2}.
 \tag{I.132}
\]
The order is arbitrary and unknown to the map, and the result
holds in every characteristic. Together with P-17 and I-40,
this proves the same balanced two-pair conclusion for every
nonzero multilinear width-at-most-three polynomial on at most
ten variables, without a homogeneity or translation-invariance
promise.

For total degree at most $D$, both assertions admit the finite
sample bound
\[
 (n+1)^2\binom{D+3}{3}
 \tag{I.133}
\]
over a working field of size at least
$\max\{2n+1,D+1\}$, with the extension-input convention of I-40.
The new weight-two argument uses the common rational center
$t\mathbf1$. It does not claim arbitrary nonzero center slopes
or the ordinary unweighted centered-$\mathrm{SV}_2$ conclusion
for this entire larger class.

\paragraph{1. Reduction to an irreducible tensor of weight two.}
Work over an algebraic closure as a proof device. Suppose a
nonzero input in (I.131) vanishes under (I.132). Since the
substitution target is a field, some nonconstant irreducible
factor $g$ also vanishes. Irreducible factors of a multilinear
polynomial use disjoint variable sets; specializing the other
sets to nonzero scalar factor values realizes $g$, up to scalar,
as a width-at-most-three polynomial in the induced order, as
proved in I-37.

These factors are homogeneous. Indeed the minimum and maximum
total degrees in a product add, and equality of those degrees
for a homogeneous product forces equality for every factor.
Each factor is also invariant under common translation of its
own variables. To see this, divide the translated product
identity by the original product over $E(s)$. The resulting
factors are rational functions in disjoint sets of variables,
so each ratio belongs to $E(s)^\times$. Comparing its highest
homogeneous part with the original homogeneous factor forces
that ratio to be one. I-40 therefore gives a nonnegative weight
for every factor. Their weights sum to $n_{\rm used}-2\deg f$,
which is at most the original $n-2\deg f$; unused coordinates
only increase the latter. Thus the chosen factor has weight
at most two.

Weights zero and one are already covered by I-40, as are all
inputs on at most nine variables. We may henceforth assume
that $f$ is irreducible, depends on every remaining variable,
has weight $m=n-2D=2$, and has $n\ge10$. Construct its actual
minimal affine cores by I-40's support inclusions, and pad
consistently to width three. We use those new minimal cores,
not a possibly nonminimal representation inherited from the
original factorization. Write the unpadded bond ranks as $r_i$.
Every interior cut has $r_i\ge2$: rank one would split the
scalar into a product in two nonempty disjoint variable sets.
Moreover
\[
 r_1=r_{n-1}=2,
 \tag{I.134}
\]
since a single multilinear variable supplies at most two
coefficient vectors. In particular the first and last internal
factors, $A_2$ and $A_{n-1}$, are singular after padding.

As in I-37, every strict internal interval has generic rank at
least two. A rank-one interval has a polynomial outer
factorization over the multivariate UFD and splits the scalar
into two nonconstant factors, because the first and last
variables occur on opposite sides. Each such interval therefore
has rank two or three. I-40 bounds every run of rank-three
factors by two. The only new rank-preservation issue is a pair
of consecutive rank-three cores between rank-two anchors.

\paragraph{2. A saturated run has invariant bounding support spaces.}
Suppose $A_{j+1},A_{j+2}$ are both rank three. They are a
maximal such run. The determinant characters of I-40 satisfy
\[
 \delta_j=0,\qquad \delta_{j+1}=3,
 \qquad\delta_{j+2}=6,
 \qquad\delta'_{j+2}=0.
 \tag{I.135}
\]
Indeed both increments are three and the allowed interval is
$0\le\delta_i\le3m=6$ at these cuts.

A highest-weight line of weight zero is fixed by the whole
$\mathrm{SL}_2$. For completeness, the Bruhat identity in the
I-40 audit shows that its opposite-unipotent orbit is a
polynomial of degree at most zero; it is therefore constant.
The two unipotent subgroups generate the group. Apply this to
the determinant lines in (I.135). A subspace is determined by
its nonzero decomposable exterior vector, equivalently by the
vectors whose exterior product with that vector is zero.
Thus $L_j$ and $R_{j+2}$ themselves are $\mathrm{SL}_2$-stable
three-dimensional spaces.

Taking all single-coordinate contractions propagates this
stability to $L_{j-1}$ and to $R_{j+3}$, respectively. These
are the actual support spaces: contracting $L_j$ by all
functionals on its final physical factor spans $L_{j-1}$,
and the analogous statement holds on the right. Every torus
weight in $L_j\subseteq V^{\otimes j}$ is congruent to $j$
modulo two. Since its determinant weight is zero and its
dimension is three, $j$ is even. Also $n$ is even because
$n-2D=2$. The adjacent support spaces lie in odd tensor powers.
Weyl pairing of their nonzero integer weights makes their
dimensions even; the width bound and nonzeroness therefore give
\[
 \dim L_{j-1}=\dim R_{j+3}=2.
 \tag{I.136}
\]
This also shows that both adjacent anchors exist away from the
scalar endpoints.

\paragraph{3. The canonical anchors, without a semisimplicity assumption.}
Let $W=L_{j-1}$. Its two torus weights are $q,-q$ for an odd
positive integer $q$. If $q>1$, all four weights of $W\otimes V$
are nonzero, namely $q+1,q-1,-q+1,-q-1$. Every
$\mathrm{SL}_2$-stable subspace of this tensor product then has
even dimension, again by Weyl pairing. This contradicts its
three-dimensional stable subspace $L_j$. Therefore $q=1$.
This elementary argument excludes nontrivial Frobenius twists
as well as every other possible larger character.

A two-dimensional rational $\mathrm{SL}_2$-module with weights
$1,-1$ is the standard module in every characteristic. One can
check the convention directly: torus conjugation forces the
upper and lower unipotent entries to be $b s$ and $c s$;
neither is zero, since the two unipotents generate the group.
The Gauss identity for their product forces $bc=1$, and a
basis rescaling gives the standard matrices. This is also the
degree-one case of the established I-B09 input.

Now identify $W$ with $V$. A stable three-plane in $V\otimes V$
must contain the one-dimensional weight-two space: otherwise
Weyl symmetry would remove weight minus two as well, leaving
only the two-dimensional weight-zero space. Applying the
opposite unipotent to $e_0\otimes e_0$ shows that the three-plane
contains, and therefore equals,
\[
 \Gamma^2 V=\operatorname{span}\{
 e_0e_0,\ e_0e_1+e_1e_0,\ e_1e_1\}.
 \tag{I.137}
\]
Here juxtaposition denotes tensor product. The three displayed
vectors are independent in characteristic two as well. We do
not identify this submodule with its dual or invoke a
semisimple decomposition of $V\otimes V$.

Using its displayed weight basis, the left anchor is exactly
\[
 C(u)=\begin{pmatrix}1&u&0\\0&1&u\end{pmatrix},
 \qquad u=X_j.
 \tag{I.138}
\]
The same argument on the right identifies $R_{j+2}$ with the
three-plane (I.137) in $V\otimes R_{j+3}$, where the latter
two-dimensional residual space is standard.

The two full-rank cores have nonzero constant determinants by
I-40. Make their constant matrices the identity by constant
changes of the two intermediate bases, starting from the chosen
left weight basis. The torus actions on those bases are shifted
successively by weights one and two. Their linear coefficient
matrices consequently have entries only one step above the
diagonal: their row weight minus column weight must equal two.
Thus their product has the form
\[
 M(x,y)=
 \begin{pmatrix}1&A&B\\0&1&C\\0&0&1\end{pmatrix},
 \tag{I.139}
\]
where $A,C$ are homogeneous linear forms in $x=X_{j+1}$ and
$y=X_{j+2}$, and $B$ is homogeneous quadratic (in fact a scalar
multiple of $xy$). All these changes of basis are over the
coefficient field and are independent of the evaluation-center
parameters.

The left weights at the end of the run are $(4,2,0)$, while
the right weights in $R_{j+2}$ are $(2,0,-2)$. The tensor has
weight two and full rank three across that cut. Its coefficient
pairing is therefore antidiagonal with three nonzero coefficients
$d_1,d_2,d_3$. Expand the canonical right three-plane (I.137),
and order the next left residual basis oppositely to the right
standard basis. The right anchor is
\[
 D(v)=\begin{pmatrix}d_1v&0\\d_2&d_2v\\0&d_3\end{pmatrix},
 \qquad v=X_{j+3},\qquad d_1d_2d_3\ne0.
 \tag{I.140}
\]
This uses the actual dual pairing. In particular it remains
valid when $\Gamma^2V$ is not self-dual in characteristic two.
The full-rank constant gauges identify torus actions only;
they do not assert an isomorphism of the neighboring
upper-unipotent modules.

\paragraph{4. The four-coordinate connector cannot vanish at one pole.}
The connecting matrix between the two-dimensional cuts is
$C(u)M(x,y)D(v)$. Its determinant is
\[
 \begin{split}
 \det(CMD)={}&d_1d_2v^2+d_1d_3v(C+u)\\
 &+d_2d_3\bigl(u^2+Au+AC-B\bigr).
 \end{split}
 \tag{I.141}
\]
This follows by multiplying the two-by-three and three-by-two
matrices; equivalently its $u^2$ coefficient is the bottom
two-row minor of $MD$. It is a homogeneous quadratic in the
four distinct coordinate variables, and its $u^2$ coefficient
is the constant $d_2d_3\ne0$.

Set $\alpha_i=t+z/(a_i-y_0)$, where $t,z,y_0$ are independent.
At $y_0=a_j$, only $u=\alpha_j$ has a pole among these four
coordinates. The term $d_2d_3u^2$ has a double pole with nonzero
leading coefficient $d_2d_3z^2$, and every other term in
(I.141) has pole order at most one there. Hence the connector
determinant survives this substitution in every characteristic.

We now verify all strict internal interval ranks at this same
one-pole center. Every individual full-rank core remains
invertible; its determinant is a nonzero constant. Every
rank-two core keeps rank two, because some nonzero univariate
minor is evaluated at a transcendental coordinate. In the
second compound word, these are respectively full-rank and
rank-one factors. A rank-two interval has a nonzero compound
product. Factoring each rank-one compound by a nonzero
univariate pivot writes it as nonzero endpoint vectors times
the scalar connectors between consecutive rank-one factors.

If no or one full-rank factor lies between two such factors,
the connector depends on at most three coordinates and survives
by I-38's three-coordinate dominance. If two lie between them,
they form exactly the saturated run just analyzed. At its
adjacent two-dimensional cuts the compound connector is the
determinant (I.141), up to nonzero univariate pivot factors,
so it too survives. Thus every compound interval product that
was nonzero stays nonzero. A rank-three interval consists
entirely of invertible factors and stays invertible. All
strict internal interval ranks are therefore preserved
simultaneously. This asserts preservation of ranks, not of
every individual minor.

\paragraph{5. Endpoints and compatible constant gauges.}
Absorb the first and last factors into the affine row and column.
By (I.134), the first and last internal factors are rank two.
The polynomial row $A_1(X_1)A_2(X_2)$ is nonzero, since otherwise
the full scalar would be zero. Pairwise dominance preserves a
nonzero entry at the center. The same holds for
$A_{n-1}(X_{n-1})A_n(X_n)$. Propagation through the remaining
rank-preserved constant intervals is injective on the persistent
two-dimensional active spaces, so both endpoint active
projections are nonzero. No four-coordinate endpoint assertion
is being assumed.

The full constant interior product has rank two, and every
subinterval has rank two or three. The compatible-space proof
of I-38, part~4, applies to this rank profile without any bound
on run length: propagate a two-dimensional complement of the
full-product row kernel, and take the suffix preimages of a
final complementary line. Each passive line is either the
rank-two suffix kernel or the inverse image under an invertible
suffix. In the transported bases the active blocks are
$I_2$, the singular passive blocks are zero, and the invertible
passive blocks can be scaled to one.

The endpoint active vectors pair to $f(\boldsymbol\alpha)=0$.
Normalize them to $e_1^{\mathsf T}$ and $e_2$, and remove the
passive endpoint components by the upper and lower shears of
I-37. They leave the first and last singular constants fixed;
if those coincided the same identities would permit both
shears. Over $K=\overline E(t,z,y_0)$, with independent residual
variables $Y_i=X_i-\alpha_i$, the resulting scalar is
\[
 f(\boldsymbol\alpha+\mathbf Y)=
 (e_1^{\mathsf T}+Y_1p)
 \prod_{i=2}^{n-1}(C_i+Y_iN_i)
 (e_2+Y_nq),
 \quad C_i\in\{P,I_3\},\quad P=\operatorname{diag}(1,1,0).
 \tag{I.142}
\]
The $I_3$ runs have length at most two, and $C_2=C_{n-1}=P$.
Pole extraction from the assumed two-pole identity, exactly
as in I-37, gives $f_i(\boldsymbol\alpha)=0$ for every $i$.
It follows that
\[
 p_2=q_1=(N_i)_{12}=0.
 \tag{I.143}
\]
At each singular factor the coefficient of $Y_i$ in
$\det(P+Y_iN_i)=0$ also gives $(N_i)_{33}=0$.
No such assertion is made at invertible factors.

\paragraph{6. Three tangent directions kill the bandwidth-three Hessian.}
Let $H=(f_{ij}(\boldsymbol\alpha))$, and put
$h_i=1/(a_i-y_0)$. Differentiate the already zero vector
$(f_i(\boldsymbol\alpha))$ with respect to $t,z,y_0$.
Equivalently take first coefficients under separate parameter
increments. Since $z\ne0$ in $K$, this gives
\[
 H\mathbf1=H\mathbf h=H\mathbf h^{\circ2}=0.
 \tag{I.144}
\]
All three identities hold in every characteristic. The diagonal
of $H$ is zero by multilinearity.

If a singular position lies strictly between $i$ and $j$, the
intervening constant product contains $P$. It removes the
passive state, while the active derivative blocks in (I.143)
are lower triangular, so $H_{ij}=0$. Otherwise all those
constants are identities, and $H_{ij}$ is exactly the product
of the state-one-to-passive derivative entry at $i$ and the
passive-to-state-two derivative entry at $j$. The endpoint
versions are $p_3$ and $q_3$. These statements follow directly
from (I.142); the other active contributions contain a zero
$(1,2)$ entry. Since no invertible run has length three,
\[
 H_{ij}=0\qquad\text{when }|i-j|>3.
 \tag{I.145}
\]

The first row of $H$ has at most three remaining entries.
For three distinct future indices $j,k,\ell$, the three
equations (I.144) have coefficient determinant
\[
 \det\begin{pmatrix}
 1&1&1\\h_j&h_k&h_\ell\\h_j^2&h_k^2&h_\ell^2
 \end{pmatrix}
 =(h_k-h_j)(h_\ell-h_j)(h_\ell-h_k)\ne0.
 \tag{I.146}
\]
The one- and two-entry cases use the corresponding initial
rows of this Vandermonde matrix. Thus the first row is zero.
By symmetry, the earlier entries of the next row are now zero,
leaving at most three future entries to which the same argument
applies. Induction proves $H=0$.

Finally expand (I.142) in state paths. Direct passage from
active state one to active state two is excluded by (I.143).
Before a path first reaches active state two, take its last
departure from active state one into the passive state; if it
starts passively, take its passive endpoint entry instead.
Any intervening passive stays must occur at invertible factors,
because a singular factor has both its passive constant and
linear entry zero. The product of the departure and arrival
entries is precisely the corresponding $H_{ij}=0$.
Additional passive stays only multiply this zero product.
The same reasoning applies to a final output through $q_3$;
an output through active state one is excluded by $q_1=0$.
Therefore every output path is zero, contradicting the
nonzero polynomial $f$. This proves (I.132) under (I.131).

\paragraph{7. All multilinear inputs on at most ten variables.}
I-40 already covers at most nine variables. For ten variables,
the established balanced-ideal degree interval is $3\le d\le7$.
The cubic theorem of P-17 excludes degree three, and reversal
excludes degree seven. Consider a homogeneous quartic
annihilator of width at most three. Its common translation has
the form $f+s g$ because the ideal has no components of degree
at most two. A nonzero $g$ would be a cubic annihilator with
the same width bound: as in I-40, it is the leading coefficient
of a product of invertible local translations, and the
coefficient-cut four-minors remain zero at the leading limit.
P-17 forces $g=0$. The quartic is translation-invariant of
weight $10-2\cdot4=2$, so the theorem just proved excludes it.
Reversal also excludes degree six.

For a homogeneous quintic annihilator, every positive-power
translation coefficient has degree at most four and remains
in the balanced ideal. If any were nonzero, the coefficient
of the highest occurring power would retain width at most
three by the same limit argument. The already excluded cubic
and quartic cases, and the absence of lower degrees, rule this
out. The quintic is therefore translation-invariant of weight
zero, and I-40 excludes it as well. This covers every allowed
homogeneous degree. Passing to a lowest nonzero homogeneous
component, which preserves the ROABP width in the same order,
proves the general multilinear ten-variable assertion.

The finite balanced sampler is exactly I-40, part~6: multiply
the pullback by $\prod_i(a_i-y_1)(a_i-y_2)$, interpolate its two
selector degrees at most $n$ away from every node, and its
three amplitude parameters of total degree at most $D$ on
the Newton simplex. This proves (I.133) and its field scope
for both classes.\hfill$\square$

\paragraph{Provenance and the next obstruction.}
The proof uses the established representation input I-B09 and
the fully proved I-37--I-40 normalization and rank arguments.
The purity agent independently checked the canonical
$\Gamma^2$ anchors, the actual dual pairing in characteristic
two, the connector determinant and the endpoint simplification.
Those checks do not assume that $\Gamma^2V$ is self-dual or
that constant full-rank gauges intertwine the entire Borel
action. The present result makes no historical priority claim.

At weight three, I-40 permits an invertible run of length three.
The three tangent directions in (I.144) do not by themselves
kill a bandwidth-four symmetric matrix. Shorter nonsaturated
runs also need not have both invariant bounding support spaces.
The general mixed affine two-pair gate, the unrestricted
eleven-variable case, and polynomial dependence on a growing
ROABP width remain separate questions.


\subsection{I-42. Reset Hessians, sharp block thresholds, and saturated runs}
\label{I-42}

\paragraph{Scope.}
The abstract Hessian examples below are not ROABP annihilators.
The positive polynomial consequence retains its explicit
saturated-run promise.

\paragraph{1. An exact reset-Hessian lemma.}
Let $K$ be a field, let $h_1,\ldots,h_n$ be distinct elements,
and write $v_i=(1,h_i,h_i^2)^{\mathsf T}$. Let $S\subseteq[n]$
be a set of reset positions. Suppose a symmetric zero-diagonal
matrix $H$ has the following form for $i<j$:
\[
 H_{ij}=\begin{cases}
 a_i b_j,&\text{if no element of $S$ lies strictly between $i,j$},\\
 0,&\text{otherwise},
 \end{cases}
 \qquad a_i,b_i\in K.
 \tag{I.147}
\]
Let $t(i)$ be the first reset strictly after $i$, or $n$ if
there is none. Assume
\[
 H\mathbf1=H\mathbf h=H\mathbf h^{\circ2}=0.
 \tag{I.148}
\]
If $t(i)-i\le5$ for every $i<n$, then $H=0$ in every
characteristic. If $\operatorname{char}K\ne2$, the weaker
bound $t(i)-i\le6$ still implies $H=0$.

\paragraph{Proof.}
Suppose $H\ne0$, and let $i$ be its first nonzero row.
All earlier rows and, by symmetry, earlier columns are zero.
Put $t=t(i)$. Some future entry in this row is nonzero,
so $a_i\ne0$. Let $k$ be the first index in $(i,t]$ with
$b_k\ne0$. The three row equations imply
\[
 \sum_{j=k}^t b_jv_j=0.
 \tag{I.149}
\]
Any at most three of the $v_j$ are linearly independent.
Consequently at least four of these $b_j$ are nonzero, and
$t-k\ge3$. Under the first bound $t-i\le5$, this implies
$k-i\le2$.

Since $k\le t-3$, it is not the terminal reset. Every future
entry in its row is still in the same block. Earlier columns
before $i$ are zero, so its row equation is
\[
 b_k\sum_{j=i}^{k-1}a_jv_j+
 a_k\sum_{j=k+1}^t b_jv_j=0.
\]
Use (I.149) and divide by $b_k\ne0$ to obtain
\[
 \sum_{j=i}^{k-1}a_jv_j-a_kv_k=0.
 \tag{I.150}
\]
There are at most three distinct vectors here, so their
independence forces $a_i=0$, a contradiction.

Now suppose $\operatorname{char}K\ne2$ and only $t-i\le6$
is assumed. The same argument applies unless $k=i+3$.
In that remaining case $t=i+6$, and all four coefficients
$b_k,b_{k+1},b_{k+2},b_{k+3}$ are nonzero. Equation (I.150)
is a nonzero relation among four Vandermonde columns; every
coefficient is nonzero, since any three are independent.
In particular $a_k\ne0$, and
\[
 \sum_{j=i}^{k-1}a_jv_j=a_kv_k.
\]
The row at $\ell=k+1<t$ has no contribution beyond this same
block. Substitute the preceding identity and (I.149) into its
equation. This gives
\[
 (2a_kb_\ell-a_\ell b_k)v_k-a_\ell b_\ell v_\ell=0.
 \tag{I.151}
\]
The two columns are independent and $b_\ell\ne0$, so first
$a_\ell=0$ and then $2a_kb_\ell=0$, a contradiction outside
characteristic two. This proves both bounds.

\paragraph{2. Sharpness for abstract Hessians.}
Both thresholds are sharp for matrices of the stated form,
provided the field contains the displayed number of distinct
nodes. In characteristic two, take seven distinct $h_i$.
Choose nonzero vectors $a$ and $b$, supported respectively on
$\{1,2,3,4\}$ and $\{4,5,6,7\}$, with
\[
 \sum_{i=1}^4a_iv_i=0,
 \qquad\sum_{i=4}^7b_iv_i=0.
 \tag{I.152}
\]
Each kernel is one-dimensional and all its displayed entries
are nonzero, by the three-column Vandermonde property. Define
$H_{ij}=a_ib_j$ for $i<j$, impose symmetry, and put $H_{ii}=0$.
The first three rows satisfy (I.148) by the second relation;
the last three do so by the first. The fourth row is
$-2a_4b_4v_4=0$. Thus $H\ne0$ has all three kernels. It is
one reset block of seven positions, with $S=\{1,7\}$, corresponding to five
intermediate invertible positions in the matrix-word setting.

In every characteristic, use eight distinct nodes and take
$a$ supported on $\{1,2,3,4\}$ and $b$ on
$\{5,6,7,8\}$, each annihilating its four-column Vandermonde
matrix. Then
\[
 H=ab^{\mathsf T}+ba^{\mathsf T}
 \tag{I.153}
\]
is a nonzero symmetric zero-diagonal matrix, with $S=\{1,8\}$, of the same
semiseparable form and satisfies (I.148). This proves the
outside-characteristic-two threshold is also exact. These
examples concern an abstract evaluated Hessian; no fixed
polynomial integrability or ROABP realization is asserted.

\paragraph{3. Consequence for an actual rank-preserved width-three word.}
Suppose a nonzero affine multilinear scalar word has at least
one internal rank-two factor and has, at the common
one-pole center $\alpha_i=t+z/(a_i-y)$, all the internal rank
and nonzero active-endpoint properties used in I-41. Assume
its nonzero generic internal interval ranks are two or three,
and that each invertible run has length at most four, or at
most five if the characteristic is not two. Then its nonzero
scalar cannot annihilate the balanced two-pole map.

Indeed, an assumed annihilation gives the same normalized
$I/P$ word, zero first derivatives and passive-zero entries
at singular factors as (I.142)--(I.143). If an endpoint is
followed or preceded by an invertible constant string, transport
its endpoint shear unchanged through that $I_3$ string to the
first or last $P$, exactly as in I-38; this leaves every
constant $I_3$ and $P$ intact. Its Hessian has exactly
the form (I.147), with resets at singular factors and the
endpoint versions of $a_i,b_i$. The distance to the next reset
or final endpoint is at most the invertible-run length plus
one. The three tangent identities give (I.148), so part~1
forces $H=0$. The full state-path argument of I-41 then kills
every possible passive excursion and every output path,
contradicting nonzero scalar output. Thus no additional
Hessian condition remains under these run bounds once the
specified ranks and endpoints have been preserved.

\paragraph{4. A positive saturated-run class beyond weight two.}
Here is one explicit structural application that supplies those
rank hypotheses rather than assumes them. Let $f$ be an
irreducible homogeneous multilinear polynomial, depending on
all its $n$ variables, of width at most three in some order,
and invariant under common translation. Let
$m=n-2\deg f$. Suppose either $m\in\{3,4\}$ in any
characteristic, or $m=5$ in characteristic different from two.
Construct the minimal affine cores from the actual coefficient
support spaces, as in I-40. Assume that every maximal run of
rank-three cores is either isolated or has length exactly $m$.
Then balanced two-pole evaluation hits $f$, in unknown order.
Its finite grid is the same
$(n+1)^2\binom{D+3}{3}$ grid as I-40, with the same field
and extension-input conditions.

To prove this, generic strict internal ranks are at least two
by irreducibility, and the first and last internal factors
are singular by minimality, exactly as in I-41. Isolated
invertible cores give connectors in at most three coordinates,
so I-38's dominance preserves them. Consider a length-$m$
run beginning after cut $j$. Every full-rank core increments
the determinant weight by three, and its entire allowed
budget is $3m$. Hence
\[
 \delta_j=0,\qquad\delta_{j+m}=3m,
 \qquad\delta'_{j+m}=0.
 \tag{I.154}
\]
Its two bounding support spaces are actual three-dimensional
$\mathrm{SL}_2$ submodules. Since $\delta_j=0$, the index
$j$ is even; since $n-m=2D$, the right bounding suffix also
has even length. The adjacent support spaces therefore have
dimension two and odd weights. The argument of I-41,
part~3, excludes all weights other than $\pm1$ and identifies
both bounding three-planes with the actual submodule
$\Gamma^2V\subset V\otimes V$.

Normalize the $m$ full-rank constant cores to the identity.
Each linear coefficient has nonzero entries only on the first
superdiagonal by torus covariance. Their product is still
\[
 M=\begin{pmatrix}1&A&B\\0&1&C\\0&0&1\end{pmatrix},
 \tag{I.155}
\]
where $A,C$ are linear and $B$ quadratic in all $m$ run
variables: a product of three strictly upper triangular
three-by-three matrices is zero. The final left weights are
$(m+2,m,m-2)$ and the right three-plane weights are
$(2,0,-2)$. Semi-invariance of weight $m$ gives the same
nonzero antidiagonal pairing and the same right anchor
as (I.140). Thus the complete connector still has determinant
(I.141), with nonzero constant coefficient $d_2d_3$ on the
square of the left-anchor variable.

Its unique double pole preserves this connector, independently
of how many coordinates occur in $A,B,C$. The other connectors
and univariate pivots survive as above. Therefore every strict
internal interval rank is preserved at the one-pole center.
The minimal endpoint ranks again reduce their active
nonvanishing to pairwise dominance. The maximum run is at
most four in the unrestricted-characteristic cases and at
most five otherwise. Part~3 now proves the asserted scalar
nonvanishing and the claimed finite sampler.

\paragraph{Exact remaining rank gate.}
At weight three, the only missing run length for this
structural theorem is two; at weight four it is two or three.
The reset-Hessian lemma shows that these cases need only a
rank-preservation repair, not an additional bandwidth argument.
The theorem does not claim that the saturated-run promise
survives arbitrary factorization or specialization. It is
stated for an irreducible input and its actual minimal cores.
Products of inputs already known to be hit also remain hit
by the same formal map, since its target is a field, but no
unproved structural promise for their factors is inferred.

The overlapping seven-node characteristic-two example and the
disjoint eight-node example are barriers to extending this
specific Hessian-kernel argument without further information.
They do not disprove a balanced two-pair sampler or a general
polynomial-in-width hitting set. No historical priority claim
is made.


\subsection{I-43. Grassmann lines repair weight-three connectors}
\label{I-43}

\paragraph{Statement.}
Let $E$ be any field, and let $a_1,\ldots,a_n$ be distinct
nodes in $E$. Every nonzero homogeneous multilinear polynomial
$f$, having ROABP width at most three in some variable order
and satisfying
\[
 f(X+s\mathbf1)=f(X),\qquad 0\le n-2\deg f\le3,
 \tag{I.156}
\]
has nonzero pullback under the balanced two-pole map
\[
 R_i=t+\frac{z_1}{a_i-y_1}+\frac{z_2}{a_i-y_2}.
 \tag{I.157}
\]
The order need not be known, and the assertion holds in every
characteristic. Together with the arbitrary-variable cubic
theorem P-19, this proves the same conclusion for every
nonzero multilinear width-at-most-three polynomial on at most
eleven variables, without a homogeneity or translation promise.

For total degree at most $D$, the finite hitting set in both
assertions has at most
\[
 (n+1)^2\binom{D+3}{3}
 \tag{I.158}
\]
points over a working field of size at least
$\max\{2n+1,D+1\}$, with the extension-input convention of
I-40. The new conclusion is for the balanced rational map
with common center $t\mathbf1$. It does not assert the same
result for arbitrary center slopes or for the ordinary
unweighted centered-$\mathrm{SV}_2$ map.

\paragraph{Background input I-B10.}
We use the established description of a projective line in a
Grassmannian: its three-planes all contain a fixed two-plane
and lie in a fixed four-plane. The exact exterior-algebra
input is Ghorpade, Patil, and Pillai (2009), Lemma~3, printed
p.~4 of the author/arXiv version: two independent decomposable
$\ell$-vectors have decomposable sum exactly when their
associated $\ell$-planes intersect in dimension $\ell-1$.
That statement and proof are over an arbitrary field. The
highest-weight-zero and highest-weight-one inputs are the
degree-zero and degree-one cases of I-B09, with the conventions
and elementary Bruhat audit already recorded in I-40.

\paragraph{1. The precise connector problem.}
Work over an algebraic closure for the proof. The factor and
minimal-core reduction of I-41 applies unchanged: homogeneous
multilinear factors use disjoint variable sets, inherit the
width by specializing complementary factors, and each is
translation-invariant. Their nonnegative weights add to the
weight on the used coordinates. Thus a counterexample to
(I.156) would have an irreducible factor of weight at most
three. I-41 covers weights at most two. We may assume that
$f$ is irreducible, depends on all its variables, and has
weight $m=n-2D=3$.

Construct its actual minimal affine cores from the coefficient
support inclusions, as in I-40, and pad consistently to width
three. Every proper bond rank is at least two, and the first
and last are exactly two. Every strict internal interval has
generic rank two or three; a rank-one polynomial outer
factorization would factor the scalar into two nonconstant
polynomials, as in I-37. In particular the first and last
internal factors are singular. I-40 bounds every run of
rank-three factors by three.

It suffices first to preserve the rank-two contractions
between consecutive singular factors at the common one-pole
center
\[
 \alpha_i=t+\frac{z}{a_i-y}.
 \tag{I.159}
\]
Here the two singular factors themselves are included in the
connector. Consecutive singular factors with no intervening
full-rank factor, or with one, involve at most three distinct
coordinates; I-38's all-characteristic triple dominance
preserves their nonzero compound matrices. A run of length
three is saturated and is covered by I-42, part~4.

Consider therefore two full-rank cores $A_{j+1},A_{j+2}$,
between singular anchors $A_j,A_{j+3}$. Let $\delta_j$ be the
determinant weight of the left support $L_j$. The two increments
of three and the total budget nine give
\[
 \delta_j\in\{0,1,2,3\},\qquad
 \delta'_{j+2}=3-\delta_j.
 \tag{I.160}
\]
Reversing the witnessing variable order exchanges left and
right support spaces and transposes this connector. Thus it
is enough to handle $\delta_j=0$ and $\delta_j=1$. This is
order reversal, not reciprocal multilinear reversal, and
does not change the generator.

\paragraph{2. A single invariant anchor is enough.}
Suppose $\delta_j=0$. The argument of I-41 identifies the
left support with the actual invariant submodule
$\Gamma^2V\subset V\otimes V$, its preceding support with
the standard two-dimensional module, and its anchor with
\[
 C(u)=\begin{pmatrix}1&u&0\\0&1&u\end{pmatrix},
 \qquad u=X_j.
 \tag{I.161}
\]
Normalize the two full-rank constant cores to the identity by
constant basis changes. Torus covariance gives
\[
 M(x_1,x_2)=
 \begin{pmatrix}1&A&B\\0&1&C\\0&0&1\end{pmatrix},
 \tag{I.162}
\]
where $A,C$ are linear in the two internal variables and
$B$ is quadratic. Write the arbitrary opposite affine anchor
as $D(v)$, with rows $D_1(v),D_2(v),D_3(v)$, where
$v=X_{j+3}$. Its generic matrix rank is two. Its three rows
are linearly independent as coefficient tensors in
$E[v]^{r_{j+3}}$: a constant relation between them would be
a relation between the three actual suffix coefficient
functions at the preceding minimal bond. The constant gauges
above preserve this independence.

The wedge of the two rows of $C(u)M D(v)$ is the row of its
second compound. Its coefficient on $u^2$ is exactly
\[
 D_2(v)\wedge D_3(v).
 \tag{I.163}
\]
If this polynomial vector is nonzero, substitute (I.159) and
consider the pole at $y=a_j$. Every other coordinate is
regular there. The coefficient of the double pole is a
nonzero scalar multiple of
$z^2[D_2(v_0)\wedge D_3(v_0)]$, where
$v_0=t+z/(a_{j+3}-a_j)$. This vector is nonzero because
$v_0$ is transcendental over the coefficient field. Hence
this connector retains rank two.

It remains to treat $D_2\wedge D_3=0$. A primitive rank-one
factorization over the univariate polynomial ring gives
\[
 D_2(v)=p_2(v)q(v),\qquad D_3(v)=p_3(v)q(v).
\]
Both rows are nonzero and linearly independent over $E$.
They are affine. If the primitive vector $q$ had positive
degree, the degree bound would force both $p_2,p_3$ to be
constants, a contradiction. Thus $q$ is a nonzero constant
row, $p_2,p_3$ are independent affine polynomials, and
$D_1(v)\wedge q$ is nonzero because $D$ has generic rank two.
Direct multiplication now gives the complete exterior row
\[
 \bigl[p_2(v)+(C+u)p_3(v)\bigr]
       \bigl[D_1(v)\wedge q\bigr].
 \tag{I.164}
\]
At $y=a_j$ its first factor has a nonzero simple-pole
coefficient proportional to $z p_3(v_0)$, whereas the second
is regular and nonzero. The terms $A,B$ cancel from this
exterior expression, so no other internal dependence can
cancel that pole. Again the connector retains rank two.
This proves the case $\delta_j=0$, without assuming that the
opposite anchor is invariant or self-dual.

\paragraph{3. Determinant weight one forces two equal positive weights.}
Suppose now that $\delta_j=1$. The nonzero Pl\"ucker vector
of $L_j$ is fixed by the upper unipotent and has torus weight
one. Its $\mathrm{SL}_2$ orbit spans the standard
two-dimensional module, by the highest-weight-one input.
Its projective orbit is therefore a line, all of whose
points are decomposable three-vectors. I-B10 gives fixed
spaces
\[
 K\subset L_j\subset T,\qquad \dim K=2,\quad \dim T=4.
 \tag{I.165}
\]
They are the intersection and span of all planes in that
orbit line, so both are $\mathrm{SL}_2$-stable.

We give the needed small-width consequence explicitly; it
does not assume semisimplicity in positive characteristic.
The parity of tensor weights and $\delta_j=1$ imply that
$j$ is odd. Thus the two torus weights of the stable space
$K\subset V^{\otimes j}$ are $q,-q$ for an odd positive
integer $q$, by Weyl symmetry. Take all contractions of its
last physical factor. Their span $W_K$ is an
$\mathrm{SL}_2$-stable subspace of the preceding actual
support $L_{j-1}$, and hence has dimension at most three.

Assume $q>1$, and choose weight vectors $k_+,k_-$ of $K$.
Their coefficient vectors on the last factor can have only
the four distinct weights
\[
 q-1,\ q+1,\ -q-1,\ -q+1.
 \tag{I.166}
\]
Weyl symmetry pairs the first with the last and the second
with the third. Nonzero coefficient vectors of distinct
weights are independent. Since $\dim W_K\le3$, exactly one
of those pairs can occur. Up to exchanging $e_0,e_1$, this
means
\[
 k_+=p\otimes e_0,\qquad k_-=r\otimes e_1,
 \tag{I.167}
\]
where $p,r$ are independent, or the same display with the
two final basis vectors exchanged.

Apply the unipotent that sends $e_0$ to $e_0+s e_1$ in
(I.167), and write its action on $p$ as $\rho(s)p$.
Since $K$ is stable, the resulting tensor
$\rho(s)p\otimes e_0+s\rho(s)p\otimes e_1$ must remain in
the span of the two displayed tensors. Its $e_0$ coefficient
forces $\rho(s)p$ into $E(s)p$; its $e_1$ coefficient then
forces the same nonzero vector into $E(s)r$, a contradiction.
For the exchanged case use the opposite unipotent. This
argument works over the rational function field in $s$ and
in every characteristic. Hence $q=1$.

The determinant character of the stable space $K$ is zero.
The quotient line $L_j/K$ therefore has weight one. The
stable two-dimensional quotient $T/K$ has weights $1,-1$,
because a rational $\mathrm{SL}_2$-module has trivial
determinant character. Consequently the weights on $L_j$ are
\[
 (1,1,-1).
 \tag{I.168}
\]
No splitting of either short exact sequence in (I.165) was
used, and this deduction excludes Frobenius-twisted larger
weights directly.

\paragraph{4. A unique internal pole detects the remaining connector.}
Choose a weight basis as in (I.168) and again normalize the
two full-rank constant cores to the identity. Torus covariance
allows their linear matrices only in the common
negative-to-positive weight step. They thus have the form
\[
 N_\ell=v_\ell e_3^{\mathsf T},\qquad
 (v_\ell)_3=0,\qquad N_\ell N_k=0.
 \tag{I.169}
\]
Their product is
$M=I+x_1N_1+x_2N_2$. Its second compound is affine in
$x_1,x_2$: the quadratic exterior term of a rank-one update
vanishes. Multiplying by the compounds of the two arbitrary
affine anchors gives
\[
 C_2\bigl(A_j(u)M(x_1,x_2)A_{j+3}(v)\bigr)
   =P_0(u,v)+x_1P_1(u,v)+x_2P_2(u,v).
 \tag{I.170}
\]
Here the $P_\ell$ are polynomial matrices and the displayed
matrix is not identically zero, because the original
connector has generic rank two.

If some $P_\ell$ with $\ell=1,2$ is nonzero, take the pole
at the node belonging to $x_\ell$. Every other summand is
regular. The coefficient of its simple pole is a nonzero
scalar multiple of $zP_\ell(u_0,v_0)$, where
\[
 u_0=t+\frac{z}{a_j-a_{j+\ell}},\qquad
 v_0=t+\frac{z}{a_{j+3}-a_{j+\ell}}.
 \tag{I.171}
\]
The two coefficients of $z$ are distinct, so $(u_0,v_0)$
is an invertible linear change of $(t,z)$. A nonzero
polynomial matrix $P_\ell$ therefore stays nonzero. If both
$P_1,P_2$ vanish, the nonzero matrix $P_0(u,v)$ survives by
pairwise dominance. This proves preservation in the
$\delta_j=1$ case and, by order reversal, in every case of
(I.160).

\paragraph{5. Simultaneous ranks, endpoints, and the scalar output.}
Every univariate singular-factor pivot is nonzero after
(I.159), since each $\alpha_i$ is transcendental. The
preceding arguments preserve the nonzero rank-two compound
of every connector between consecutive singular factors.
All these assertions hold as rational-function identities
in the same independent parameters $t,z,y$. They therefore
hold simultaneously; no independent seed is assigned to an
individual connector.

For clarity, use the same pivot-cleared rank factorizations
as I-36: a rank-two anchor is $U_iV_i$, with both factors
of full respective rank. A connector of rank two says that
the $2\times2$ contraction between its adjacent anchors is
invertible. Products of these invertible contractions
preserve every interval containing anchors; intervals made
only of full-rank cores remain invertible. Thus all strict
internal interval ranks are preserved. This deduction uses
the common evaluated pivots, not preservation of only
constant boundary functionals.

Minimality gives singular first and last internal factors.
Pairwise dominance preserves the nonzero endpoint/first-anchor
and last-anchor/endpoint contractions. The preserved internal
rank-two identifications propagate their nonzero active
projections, exactly as in I-41. Hence all the hypotheses of
I-42, part~3, are now proved at the same one-pole center.
The maximum invertible run is three, which is inside its
all-characteristic bound four. An assumed annihilation of
(I.157) gives the reset Hessian with its three tangent kernels;
I-42 kills that Hessian and then every output path. This
contradicts $f\ne0$. The factor reduction in part~1 proves
the full assertion (I.156).

\paragraph{6. All multilinear inputs on at most eleven variables.}
I-41 already handles at most ten variables. Suppose a
homogeneous width-three annihilator on eleven variables
exists. P-19 and the known minimum balanced degree leave
only degrees $4,5,6,7$. For degree four, every nonconstant
coefficient of common translation has degree at most three.
If any were nonzero, the coefficient of the highest power
of the translation parameter would retain width at most
three by the coefficient-cut leading-limit argument of I-40.
It would contradict P-19 or the absence of smaller degrees.
Thus the quartic is translation-invariant, has weight
$11-8=3$, and is excluded by the theorem just proved.
Multilinear reversal excludes degree seven.

For degree five, the same highest-translation-coefficient
argument would produce a width-three annihilator of degree
at most four, all of which have now been excluded. The
quintic is therefore translation-invariant of weight one,
and I-40 excludes it. Reversal also excludes degree six.
Passing to a lowest nonzero homogeneous component preserves
the witnessing width and proves the assertion for arbitrary
multilinear inputs on at most eleven variables.

Finally apply the denominator-safe interpolation in I-40:
after multiplication by
$\prod_i(a_i-y_1)(a_i-y_2)$, each selector degree is at most
$n$, and the total degree in the three amplitudes is at most
$D$. The selector grid outside all nodes and the amplitude
Newton simplex give exactly (I.158), with the stated field
and extension scope.\hfill$\square$

\paragraph{Provenance and remaining scope.}
The Grassmann-line fact is an established input, not a new
claim about Grassmannians. The new use here is its interaction
with the actual adjacent coefficient support of dimension at
most three and the resulting pole calculation for constrained
ROABP connectors. All endpoint, interval, and simultaneous
substitution assertions above concern the same parameter
field; none follows merely from preserving each local space
against constant rank-one functionals.

At weight four, determinant weights at the ends of an
unsaturated run can both exceed one. Those connectors are
not covered by the two classifications above. The unrestricted
balanced two-pair question, the next higher-degree finite-
variable cases, and the polynomial-in-growing-width isolation
gate remain separate. No historical priority claim is made.


\subsection{I-44. Torus-chain connectors and the next balanced degree windows}
\label{I-44}

\paragraph{Statement.}
Fix arbitrary distinct nodes $a_1,\ldots,a_n$ over a field
$E$, and use the balanced rational two-pole map of I-43.
The following assertions hold in an unknown witnessing
ROABP order.

First, let $f$ be irreducible, homogeneous, multilinear,
common-translation invariant, and of width at most three,
and remove its unused variables. Construct the actual minimal
affine cores from its coefficient support spaces. If every
run of rank-three cores has length at most five, the common
one-pole center $\alpha_i=t+z/(a_i-y)$ preserves all strict
internal interval ranks. This assertion holds in every
characteristic. Consequently balanced two-pole evaluation
hits $f$ whenever those runs have length at most four, or
at most five if $\operatorname{char}E\ne2$.

Second, without an irreducibility or run promise, every
nonzero homogeneous multilinear common-translation-invariant
width-at-most-three polynomial is hit when
\[
 0\le n-2\deg f\le4,
 \quad\hbox{or when}\quad
 0\le n-2\deg f\le5\ \hbox{and }\operatorname{char}E\ne2.
 \tag{I.172}
\]
Together with P-19, this proves the balanced two-pole
conclusion for all nonzero multilinear width-at-most-three
polynomials on at most twelve variables in every
characteristic, and on at most thirteen variables outside
characteristic two.

All the scalar conclusions have the same explicit finite
grid as I-40:
\[
 (n+1)^2\binom{D+3}{3},\qquad
 |E|\ge\max\{2n+1,D+1\},
 \tag{I.173}
\]
for total degree at most $D$, with evaluation over a stated
working extension if necessary. These are conclusions for
the balanced rational map with common slope one. No ordinary
centered-$\mathrm{SV}_2$ or arbitrary-slope promotion is made.

\paragraph{1. Torus spectra that need no anchor classification.}
Use I-40's actual minimal cores. Normalize the constant
matrices in any full-rank run to the identity. If the initial
three integer torus weights are $\lambda_1,\lambda_2,
\lambda_3$, the successive weight lists are obtained by
adding one to every entry. Each linear matrix $N_i$ has an
entry only from a column weight to a row weight two larger.

Unless the three weights are distinct and have the form
\[
 (q+2,q,q-2),
 \tag{I.174}
\]
all the allowed entries have a common one-dimensional image
or a common one-dimensional domain, or there is at most one
allowed entry. Indeed, with three distinct weights two
composable steps require exactly (I.174); without such a
chain there is at most one step. If a weight is repeated,
the only nonzero case has two levels of multiplicities one
and two. All products $N_iN_j$ then vanish, and their linear
span consists of matrices of rank at most one.

It follows that the run product is $I+\sum_i x_iN_i$ and
its second compound is affine in all run variables. After
multiplication by the two anchor compounds it has the form
\[
 P_0(u,v)+\sum_i x_iP_i(u,v).
 \tag{I.175}
\]
The original connector has rank two, so this polynomial
matrix is nonzero. The unique-internal-pole argument of
I-43, part~4, applies without a restriction on the number of
summands: at the node of $x_i$, the two anchor values are
independent linear forms in $t,z$. A nonzero $P_i$ therefore
has nonzero residue; if only $P_0$ remains, pair dominance
applies. This proves preservation for every spectrum other
than (I.174), in every characteristic. The differences of
weights here are integer characters of the torus, not
integers reduced modulo the characteristic.

\paragraph{2. All minimal singular anchors into a three-weight chain.}
Assume (I.174), in the displayed order. A preceding singular
anchor $A(u)$ has three columns and two or three rows. Both
its row and column coefficient tensors are independent over
$E$, by actual minimality, and its generic matrix rank is
two. Torus covariance permits preceding row weights only
from
\[
 q+3,\ q+1,\ q-1,\ q-3.
 \tag{I.176}
\]
Its columns have respectively the following possible
coefficient vectors in those row-weight spaces:
\[
 (u a,b,0,0)^{\mathsf T},\quad
 (0,u c,d,0)^{\mathsf T},\quad
 (0,0,u e,f)^{\mathsf T}.
 \tag{I.177}
\]

No preceding weight can have multiplicity two. At either
extreme of (I.176) all row coefficients use a single vector,
contrary to row independence. At weight $q+1$, multiplicity
two requires $b,c$ independent. A nonzero third column must
then use one of the two lower weights, and the resulting
three-column minor is nonzero, contrary to rank two.
The argument at weight $q-1$ is the same in reverse. Three
copies of a weight cannot be spanned by its at most two
coefficient vectors either.

Thus the preceding weights are distinct. The complete list
of possibilities, with rows ordered by decreasing weight,
is as follows. Every displayed scalar is nonzero, except
that in the middle two-row case $c,d$ may vanish individually
but not both.
\[
\begin{array}{c|c}
 \hbox{preceding weights}&A(u)\\ \hline
 (q+3,q-1)&
 \begin{pmatrix}a u&0&0\\0&d&e u\end{pmatrix}\\[3pt]
 (q+1,q-1)&
 \begin{pmatrix}b&c u&0\\0&d&e u\end{pmatrix}\\[3pt]
 (q+1,q-3)&
 \begin{pmatrix}b&c u&0\\0&0&f\end{pmatrix}\\[3pt]
 (q+3,q+1,q-1)&
 \begin{pmatrix}a u&0&0\\b&0&0\\0&d&e u\end{pmatrix}\\[3pt]
 (q+1,q-1,q-3)&
 \begin{pmatrix}b&c u&0\\0&0&e u\\0&0&f\end{pmatrix}
\end{array}
 \tag{I.178}
\]
To verify completeness, the other two-row subsets miss an
entire column. Among the four three-row subsets, the two
omitted subsets give a nonzero triangular determinant.
For the two retained three-row subsets, rank two forces
the middle diagonal entry to vanish; row and column
independence force all the entries shown as nonzero.
This proves the list directly from (I.177).

Over $E(u)$ the row spaces in (I.178) have one of three
forms:
\[
\begin{aligned}
 G_1(u)&=\begin{pmatrix}b&c u&0\\0&d&e u\end{pmatrix},
       &&bce\ne0,\quad d\in E,\\
 G_2(u)&=\begin{pmatrix}b&c u&0\\0&0&1\end{pmatrix},
       &&bc\ne0,\\
 B(u)&=\begin{pmatrix}1&0&0\\0&d&e u\end{pmatrix},
       &&de\ne0.
\end{aligned}
 \tag{I.179}
\]
The passage to these row spaces uses only full-rank
univariate factors and nonzero univariate denominators.
All remain valid after the one-pole substitution, since
$\alpha_i$ is transcendental over $E$.

\paragraph{3. Either of the first two anchor forms suffices.}
The normalized run product in the chain case is
\[
 M=\begin{pmatrix}1&A&B_0\\0&1&C\\0&0&1\end{pmatrix},
 \quad
 A=\sum_i p_ix_i,\quad C=\sum_i q_ix_i,\quad
 B_0=\sum_{i<j}p_iq_jx_ix_j.
 \tag{I.180}
\]
Here each individual core is
$I+x_i(p_iE_{12}+q_iE_{23})$; no quadratic term occurs
within a single core.

Let the opposite, unreduced minimal anchor have rows
$D_1(v),D_2(v),D_3(v)$. They are independent affine
coefficient tensors, and the matrix $D(v)$ has generic
rank two. If the left row space is $G_1$, the coefficient
of $u^2$ in the exterior row of $G_1MD$ is
$ce\,D_2\wedge D_3$. If this is nonzero, its unique double
pole preserves the connector. If $D_2\wedge D_3=0$,
I-43's primitive affine factorization gives
$D_2=p_2(v)q$, $D_3=p_3(v)q$, with $q$ constant and
$D_1\wedge q\ne0$. The full exterior row becomes
\[
 b\,[d p_2+(dC+e u)p_3]\,(D_1\wedge q).
 \tag{I.181}
\]
Its simple pole at the left-anchor node is nonzero.

For $G_2$ the exterior row is
\[
 b\,D_1\wedge D_3+(bA+c u)D_2\wedge D_3.
 \tag{I.182}
\]
If $D_2\wedge D_3\ne0$, use its unique simple pole in
$u$. Otherwise the expression is the nonzero univariate
vector $bD_1(v)\wedge D_3(v)$, which also survives.
These arguments do not bound the length of the middle run.
Order reversal supplies the same conclusion if the opposite
anchor has either of these two forms in the reversed word.

\paragraph{4. Two remaining anchor forms and a quadratic pole test.}
It remains to consider form $B$ at both ends. At the end
of a run of length $\ell$, the left weights are
$(q+\ell+2,q+\ell,q+\ell-2)$. Homogeneity pairs them
antidiagonally, with nonzero constants, with the three
distinct right weights. Apply (I.179) to the reverse anchor
and remove its full-rank univariate outside factor. After
constant column rescaling and ordering, the actual right
factor in the two-dimensional contraction is
\[
 \begin{pmatrix}\alpha v&0\\\beta&0\\0&1\end{pmatrix},
 \qquad \alpha\beta\ne0.
 \tag{I.183}
\]
This is the actual coefficient pairing; it does not assume
that a three-dimensional module is self-dual.

The two-by-two contraction is therefore
\[
 \begin{pmatrix}
  \alpha v+\beta A&B_0\\
  d\beta&dC+e u
 \end{pmatrix},
 \qquad
 F=(\alpha v+\beta A)(dC+e u)-\beta dB_0,
 \tag{I.184}
\]
where $\alpha,\beta,d,e$ are nonzero. In particular $F$ is
nonzero: its $uv$ coefficient is $\alpha e$.

We record the elementary quadratic restriction imposed by
one-pole annihilation. For any homogeneous quadratic in
distinct coordinate variables, a nonzero square coefficient
has a unique double pole at that coordinate's node.
Consequently an annihilating quadratic must be multilinear.
Write it as $\sum_{i<j}c_{ij}X_iX_j$, with symmetric
coefficients off the diagonal. The coefficient of $tz$ in
its one-pole pullback, and then the simple pole of its $z^2$
coefficient at $y=a_i$, give respectively
\[
 \sum_{j\ne i}c_{ij}=0,
 \qquad
 \sum_{j\ne i}\frac{c_{ij}}{a_j-a_i}=0.
 \tag{I.185}
\]
Two distinct possible neighbors give an invertible two-row
matrix in these equations. Thus every nonzero row of an
annihilating quadratic has at least three nonzero neighbors.
This derivation uses no factor of two and holds in every
characteristic.

Apply this to (I.184). The coefficient of $x_i^2$ is
$\beta d p_iq_i$. If any is nonzero, the connector already
survives. Otherwise define the disjoint sets
$P=\{i:p_i\ne0\}$ and $Q=\{i:q_i\ne0\}$.
In the resulting multilinear quadratic, a row indexed by
$i\in P$ has precisely its nonzero edge to $u$ and its
edges to $Q$ positions strictly before $i$. A row indexed
by $i\in Q$ has its edge to $v$ and its edges to $P$
positions strictly after $i$. This follows from
\[
 AC-B_0=\sum_{i>j}p_iq_jx_ix_j
 \quad\hbox{when all }p_iq_i=0.
 \tag{I.186}
\]

Suppose $\ell\le5$. Then at least one of $|P|,|Q|$ is
at most two. If either set is empty, $B_0=0$ and $F$ is
a product of two nonzero linear forms. Such a product
survives because distinct simple poles make the one-pole
map nonzero on every nonzero linear form.

If both sets are nonempty and $|Q|\le2$, (I.185) applied
to any $P$ row forces $|Q|=2$ and forces both $Q$ positions
to precede that row. Hence every $Q$ precedes every $P$.
This again makes $B_0=0$, so the same nonzero product
survives. If $|P|\le2$, apply the symmetric argument to
the $Q$ rows; again every $Q$ must precede every $P$.
Thus (I.184) cannot annihilate the one-pole map for
$\ell\le5$.

This proves that every constrained rank-two connector with
at most five intervening full-rank cores preserves rank.
All coefficient identities, univariate factors and residues
are in the same field $E(t,z,y)$.

\paragraph{5. Completion of the rank and scalar assertions.}
The irreducible/minimal reduction ensures that strict
internal interval ranks are two or three and that the first
and last internal factors are singular, exactly as in I-41.
The two cases of part~1 and parts~2--4 cover every torus
spectrum and every connector of length at most five.
Univariate pivots remain nonzero. Their evaluated rank
factorizations and the invertible two-by-two adjacent
contractions therefore preserve all strict interval ranks,
as proved in I-43, part~5. Pair dominance supplies the
nonzero active endpoint projections. All these assertions
hold simultaneously with the same one-pole parameters.
Small words with no such internal interval are directly
covered by pair or triple dominance.

I-42, part~3, now proves scalar nonvanishing when all runs
have length at most four, or at most five outside
characteristic two. This proves the first assertion.

For (I.172), factor a hypothetical annihilator and recompute
the chosen factor's actual minimal cores, as in I-41.
Each factor is homogeneous and translation-invariant, uses
a disjoint set of variables, and has nonnegative weight.
Its weight is bounded by the total weight. I-40 bounds its
full-rank runs by that weight, so the first assertion applies
to the required factors. The target of the evaluation map
is a field, so no product of these nonzero images can vanish.
This proves (I.172), without assuming that a run promise
for one nonminimal representation passes to its factors.

\paragraph{6. Twelve and thirteen variables.}
I-43 already handles at most eleven variables. On twelve
variables, P-19 and multilinear reversal leave only
homogeneous degrees four through eight. A quartic
annihilator is translation-invariant: a highest nonzero
positive translation coefficient would inherit width at
most three and have degree at most three, contradicting
P-19 and the known minimum degree. Its weight is four,
contradicting (I.172). Reversal excludes degree eight.
For degree five the same leading-translation argument
uses the already excluded lower degrees and produces
translation invariance of weight two; I-41 excludes it,
as well as degree seven by reversal. Degree six similarly
becomes translation-invariant of weight zero and is
excluded by I-40.

On thirteen variables, outside characteristic two, the
same induction starts with a quartic of weight five,
excluded by (I.172). A quintic becomes translation-invariant
of weight three and is excluded by I-43; a sextic has
weight one and is excluded by I-40. Reversal handles their
complementary degrees. In both cases a lowest nonzero
homogeneous component inherits the witnessing width,
so these homogeneous conclusions cover every multilinear
input. The interpolation proof of I-40 gives (I.173)
without a new degree estimate.\hfill$\square$

\paragraph{Remaining scope.}
The rank-preservation conclusion for five-core runs is
valid even in characteristic two. The scalar conclusion
there remains restricted to four-core runs because the
seven-node abstract Hessian of I-42 shows that its three
kernel equations alone are insufficient. In particular,
the weight-five characteristic-two case is reduced to its
possible saturated length-five runs, with their interval
ranks now preserved; no abstract Hessian is asserted to
integrate to a counterexample.

For unrestricted run length, the two remaining anchor forms
lead to the explicit quadratic (I.184). When both internal
types have at least three positions, the short-run argument
above no longer suffices. This is a constrained connector
problem, not a claim that general shared-parameter matrices
compose. The higher-degree and growing-width sampling gates
remain open. No historical priority claim or unperformed
numerical verification is made.
